% concatenated from KN_2_April_2026.tex
\documentclass[11pt, a4paper,leqno]{amsart}
%\documentclass[12pt]{article}
\usepackage{amsmath,amsthm,amscd,amssymb,amsfonts, amsbsy}
\usepackage{latexsym}
\usepackage{txfonts}
\usepackage{exscale}
\usepackage{enumitem}
\setlist[itemize]{leftmargin=2em}
\usepackage{xcolor}
%\usepackage{hyperref}

\usepackage[colorlinks=true,
            linkcolor=blue,citecolor=blue,urlcolor=blue]{hyperref}











\newcommand{\doi}[1]{\href{https://doi.org/#1}{doi: #1}}




\usepackage{pgf}
%\usepackage{color}

\newcommand{\Diag}{\operatorname{Diag}}
\DeclareMathOperator{\sech}{sech}

\newcommand{\Bk}{\color{black}}
\newcommand{\Rd}{\color{red}}
\newcommand{\Bl}{\color{blue}}
\newcommand{\Mg}{\color{magenta}}




\newcommand{\kaj}[1]{\marginpar{\scriptsize \textbf{Kaj:} #1}}




\parskip=3pt
\textwidth=1.25\textwidth
\textheight=1\textheight

%\marginparwidth=65pt


%\renewcommand{\thefootnote}{\fnsymbol{footnote}}




\calclayout
\allowdisplaybreaks

\newtheorem{theorem}{Theorem}[section]
\newtheorem{proposition}{Proposition}[section]
\newtheorem{lemma}{Lemma}[section]
\newtheorem{corollary}{Corollary}[section]
\theoremstyle{definition}
\newtheorem{definition}{Definition}
\newtheorem{remark}{Remark}[section]
\newtheorem{example}{Example}
\newtheorem{problem}{Problem}
\newtheorem{assumption}{Assumption}

% \newenvironment{proof}[1][Proof]{\textbf{#1.} }{\ \rule{0.5em}{0.5em}}
\newcommand{\wb}{\mathbf{w}}


\numberwithin{equation}{section}
\def\barint{\,\Xint -} % \, corrects the \! used in the definition
\def\bariint{\barint_{} \kern-.4em \barint}
\def\bariiint{\bariint_{} \kern-.4em \barint}

\def\doublespaced{\baselineskip=\normalbaselineskip
    \multiply\baselineskip by 2}
\def\doublespace{\doublespaced}
\def\singlespaced{\baselineskip=\normalbaselineskip}
\def\singlespace{\singlespaced}
\newcommand{\hd}{\mbox{H-dim}\;}
\newcommand{\qdb}{\mbox{ $ \bar{C}_d ( x^*, t^* ) $ } }
\newcommand{\qd}{\mbox{ $ B_{cc} ( x^* , t^* ) $ }}
\newcommand{\mm}{\mbox{ $ { \ds \min_{\{ ( y, s ) \in C \}} } \:
\hat{ g } ( y, s )  $ } }
\newcommand{\imm}{\mbox{ $ { \ds \inf_{ ( y, s ) \in C  } }
\: \hat{g} ( y, s ) $ } }
\newcommand{\fb}[1]{{ \bf {#1}}}
\newcommand{\dar}[2]{\mbox{ $ { \ds \partial^{{#1}}{#2}}$ }}
\newcommand{\ina}[1]{- \hs{-.50cm}  \: \int\limits_{{#1}} }
\newcommand{\co}[1]{ \| { #1 } \|_{\mbox{\tiny comm}}}
\newcommand{\lp}{ L^p_{ 1, 1/2} }
\newcommand{\dst}{ I \hspace{-.15cm} D_n }
\newcommand{\dez}{I \hspace{-.15cm} D }
\newcommand{\Int}[2]{\mbox{ $ {\displaystyle \int _{#1}^{#2}} $ }}
\newcommand{\rn}[1]{\mathbf{R}^{#1}}
\newcommand{\Itl}[2]{\mbox{ ${\displaystyle \int\limits_{#1}^{#2}} $}}
\newcommand{\su}[1]{\mbox{${\displaystyle\sum_{#1}}$} }
\newcommand{\Su}[2]{\mbox{${\displaystyle \sum_{#1}^{#2}}$}}
\newcommand{\re}{ \mathbf{R}}%

%\newcommand{\ce}{ \mathbf{C}}
\newcommand{\beq}{\begin{equation}}
\newcommand{\bea}[1]{\begin{array}{#1} }
\newcommand{\eeq}{ \end{equation}}
\newcommand{\ea}{ \end{array}}
\newcommand{\ep}{\epsilon}
\newcommand{\es}{\emptyset}
\newcommand{\itl}[1]{\mbox{ ${\displaystyle \int\limits_{#1}}$}}
\newcommand{\al}{\alpha}
\newcommand{\ga}{\gamma}
\newcommand{\ps}{\psi}
\newcommand{\de}{\delta}
\newcommand{\ds}{\displaystyle}
%\newcommand{\ts}{\textstyle}
\newcommand{\rar}{\mbox{$\rightarrow$}}
\newcommand{\lar}{\mbox{$\leftarrow$}}
\newcommand{\ran}{\rangle}
\newcommand{\lan}{\langle}
\newcommand{\Ga}{\Gamma}
\newcommand{\la}{\lambda}
\newcommand{\La}{\Lambda}
\newcommand{\ar}{\partial}
\newcommand{\si}{\sigma}
\newcommand{\Si}{\Sigma}
\newcommand{\om}{\omega}
\newcommand{\Om}{\Omega}
\newcommand{\be}{\beta}
\newcommand{\ch}{\chi}
\newcommand{\Ps}{\Psi}
\newcommand{\ph}{\phi}
\newcommand{\he}{\theta}
\newcommand{\He}{\Theta}
\newcommand{\Ph}{\Phi}
\newcommand{\hs}[1]{\mbox{$ \hspace{#1}$}}
\newcommand{\vs}[1]{\mbox{$\vspace{#1}$}}
\newcommand{\sem}{\setminus}
\newcommand{\ze}{\zeta}
\newcommand{\De}{\Delta}
\newcommand{\ti}{\widetilde}
\newcommand{\noi}{\noindent}
\def\singlespace{\singlespaced}


\def \MLip {{\text{\rm Lip}_K}}

\def \rnn {{\mathbb {R}}^{N+1}}
\def \rn  {{\mathbb {R}}^{N}}
%\def \R  {{\mathbb {R}}}
%\def \N  {{\mathbb {N}}}
\def \F  {{\mathscr {F}}}
\def \x {{\xi}}
\def \g {{\gamma}}
\def \e {{\varepsilon}}
%\def \eps {{\varepsilon}}
\def \t {{\tau}}
%\def \n {{\nu}}
\def \m {{\mu}}
\def \y {{\eta}}
\def \th {{\theta}}
\def \z {{\zeta}}
%\def \p {{\partial}}
\def \a {{\alpha}}
\def \O {{\Omega}}
\def \gr{\`}
\def \d {{\delta}}
\def \OOO {\bar{\Omega}}
\def \TT {{T}}
\DeclareMathOperator*{\essosc}{\osc}



\def\mean#1{\mathchoice%
          {\mathop{\kern 0.2em\vrule width 0.6em height 0.69678ex depth -0.58065ex
                  \kern -0.8em \intop}\nolimits_{\kern -0.4em#1}}%
          {\mathop{\kern 0.1em\vrule width 0.5em height 0.69678ex depth -0.60387ex
                  \kern -0.6em \intop}\nolimits_{#1}}%
          {\mathop{\kern 0.1em\vrule width 0.5em height 0.69678ex
              depth -0.60387ex
                  \kern -0.6em \intop}\nolimits_{#1}}%
          {\mathop{\kern 0.1em\vrule width 0.5em height 0.69678ex depth -0.60387ex
                  \kern -0.6em \intop}\nolimits_{#1}}}

\def\vintslides_#1{\mathchoice%
          {\mathop{\kern 0.1em\vrule width 0.5em height 0.697ex depth -0.581ex
                  \kern -0.6em \intop}\nolimits_{\kern -0.4em#1}}%
          {\mathop{\kern 0.1em\vrule width 0.3em height 0.697ex depth -0.604ex
                  \kern -0.4em \intop}\nolimits_{#1}}%
          {\mathop{\kern 0.1em\vrule width 0.3em height 0.697ex depth -0.604ex
                  \kern -0.4em \intop}\nolimits_{#1}}%
          {\mathop{\kern 0.1em\vrule width 0.3em height 0.697ex depth -0.604ex
                  \kern -0.4em \intop}\nolimits_{#1}}}

\newcommand{\aveint}[2]{\mathchoice%
          {\mathop{\kern 0.2em\vrule width 0.6em height 0.69678ex depth -0.58065ex
                  \kern -0.8em \intop}\nolimits_{\kern -0.45em#1}^{#2}}%
          {\mathop{\kern 0.1em\vrule width 0.5em height 0.69678ex depth -0.60387ex
                  \kern -0.6em \intop}\nolimits_{#1}^{#2}}%
          {\mathop{\kern 0.1em\vrule width 0.5em height 0.69678ex depth -0.60387ex
                  \kern -0.6em \intop}\nolimits_{#1}^{#2}}%
          {\mathop{\kern 0.1em\vrule width 0.5em height 0.69678ex depth -0.60387ex
                  \kern -0.6em \intop}\nolimits_{#1}^{#2}}}


\newcommand{\Haus}[1]{\mathcal {H}^{#1}}
\newcommand{\tta}{\widetilde{a}}
\newcommand{\tl}{\widetilde{L}}
\newcommand{\tF}{\widetilde{F}}
\newcommand{\Dim}{\textnormal{dim}_{\mathcal{H}}}
%\newcommand{\CC}{\subset\hskip-.2em\subset }
\newcommand{\rif}[1]{(\ref{#1})}
\newcommand{\pezzo}[1]{(1+|z|^2)^\frac{#1}{2}}
\newcommand{\pezzod}[1]{(1+|Du|^2)^\frac{#1}{2}}
\newcommand{\tit}{t_2}
\newcommand\ap{``}
%\newcommand{\tit}{\widetilde{t}}
\newcommand{\pezzoo}{(1+|z|^2)}
\newcommand{\pezzood}{(1+|Du|^2)}
\newcommand{\taus}{\tau_{s,h}}
\newcommand{\esse}{s \in \{1,\ldots,n-1\}}
\newcommand{\essen}{s \in \{1,\ldots,n\}}
\newcommand{\modu}{[u]_{0,\lambda}}
\newcommand{\pezzodh}[1]{(1+|Du_h|^2)^\frac{#1}{2}}
\newcommand{\ts}{\textstyle}
\newcommand{\mmm}[1]{(s^2+|z|^2)^{\frac{p-#1}{2}}}
\newcommand{\ssim}{\sigma(q)}
\newcommand{\BMO}{\textnormal{BMO}}
\newcommand{\VMO}{\textnormal{VMO}}
\newcommand{\tu}{\widetilde{u}}
\newcommand{\tv}{\widetilde{v}}
\newcommand{\pot}[1]{{\bf I}_{#1}^{\mu}}
\newcommand{\tuq}{\widetilde{q}}
 \newcommand{\intav}{-\hskip -1.1em\int}
  \newcommand{\mR}{B_R}



\def\eqn#1$$#2$${\begin{equation}\label#1#2\end{equation}}
\def\charfn_#1{{\raise1.2pt\hbox{$\chi
_{\kern-1pt\lower3pt\hbox{{$\scriptstyle#1$}}}$}}}
\def\diam{\operatorname{diam}}
\def\qq1{q_*}
\def\q2{q_{**}}
\def\dist{\operatorname{dist}}
\def\e{\boldsymbol e}
\def\ep{\varepsilon}
\def\en{\mathbb N}
\def\er{\mathbb R}
\def\erm{\er^m}
\def\ff{\varphi}
\def\co{\operatorname{co}}
\def\id{\operatorname{id}}
\def\lip{\operatorname{lip}}
\def\loc{\operatorname{loc}}
\def\osc{\operatorname{osc}}
\def\rd{\mathbb R^d}
\def\rk{\mathbb R^k}
\def\O{\rm O}
\def\rn{\mathbb R^N}
\def\qe{\mathbb Q}
\def\ma{\mathbb R^{nN}}
\def\dive{\operatorname{div}}
\def\spt{\operatorname{spt}}
\def\ve{\mathbb V}
\def\trinom#1#2#3{\pmatrix #1\\#2\\#3\endpmatrix}
\def\zed{\mathbb Z}
\newdimen\vintbar
\vintbar12pt
\def\vint{-\kern-\vintbar\int}
\def\we{\boldsymbol w}
\def\xiv{\boldsymbol\xi}
\def\zet{\mathbf Z}
\def\unitball{\boldsymbol\alpha}
\def\A{\mathcal A}
\def\B{\mathcal B}
\def\C{\mathcal C}
\def\D{\mathcal D}
\def\E{\mathcal E}
\def\F{\mathcal F}
\def\G{\mathcal G}
\def\H{\mathcal H}
\def\I{\mathcal I}
\def\J{\mathcal J}
\def\L{\mathcal L}
\def\K{\mathcal K}
%\def\N{\mathcal N}
\def\O{\mathcal O}
\def\P{\mathcal P}
\def\Q{\mathcal Q}
%\def\er{\mathcal R}

\newcommand{\rset}{\mathbb{R}}
%\newcommand{\E}{\mathbb{E}}
\newcommand{\Dset}{{\cal X}}
\newcommand{\Dpnt}{x}
\newcommand{\rmd}{\mathrm{d}}
\newcommand{\Bm}{\mathrm{B}}
\newcommand{\Lm}{\mathrm{L}^\alpha}
\newcommand{\Id}{\mathbf{I}}
%\newcommand{\beq}{\begin{eqnarray}}
%\newcommand{\eeq}{\end{eqnarray}}
%\newcommand{\beqs}{\begin{eqnarray*}}
%\newcommand{\eeqs}{\end{eqnarray*}}
\newcommand{\umut}[1]{{\color{cyan} #1}}
\newcommand{\ls}[1]{{\color{purple} #1}}
\newcommand{\mg}[1]{{\color{red}MG: #1}}
\newcommand{\sas}{{\cal S}\alpha{\cal S}}


\newcommand{\dint}{\int\!\!\!\int}

\def\S{\mathcal S}
\def\L{\mathcal L}
\def\T{\mathcal T}
\def\U{\mathcal U}
\def\V{\mathcal V}
\def\W{\mathcal W}
%\def\X{\mathcal X}
\def\Y{\mathcal Y}
\def\Z{\mathcal Z}
\def\0{\boldsymbol 0}
\def\px{p(x)}
\def\pxs{p(x)(1+\sigma)}
\newcommand{\aint}{-\!\!\!\!\!\! \int}
\newcommand{\Wq}{W^{1,q}}
\newcommand{\Wp}{W^{1,p}}
\newcommand{\pq}{(p,q)}
\newcommand{\tom}{\widetilde{\Omega}}
\newcommand{\ratio}{\nu, L}
\newcommand{\ttau}{\pi_{h}}
\newcommand{\MM}{\mathcal{M}}
\newcommand{\Qin}{Q_{\textnormal{inn}}}
\newcommand{\Qou}{Q_{\textnormal{out}}}
\newcommand{\dro}{(\varrho-t)}


\newcommand{\Lq}{L^{q}}
\newcommand{\divo}{\textnormal{div}}
\newcommand{\dia}{\textnormal{diam}}
%\newcommand{\Lp}{L^{p}}
\newcommand{\tq}{\widetilde{Q}}
\newcommand{\tmu}{\widetilde{\mu}}
\newcommand{\deo}[1]{{|#1|}^{\beta}}
\newcommand{\deot}[1]{{|#1|}^{\widetilde{\beta}}}
\newcommand{\tep}{\widetilde{\ep}}
% \newcommand{\mean}[1]{-\hskip-1.08em\int_{#1}}
 \newcommand{\means}[1]{-\hskip-1.00em\int_{#1}}
\newcommand{\Lpx}{L^{\px}}
\newcommand{\di}{\textnormal{dist}}
\newcommand{\omp}{\Omega^{\prime}}
\newcommand{\ompp}{\Omega^{\prime \prime}}
\newcommand{\distc}{\textnormal{dist}(\omp, \partial \ompp)}
\newcommand{\M}{M^*}
%\newcommand{\CC}{\subset\hskip-.2em\subset }
\newcommand{\Ma}{\er^{N\times n}}
\newcommand{\tb}{\widetilde{\beta}}
\newcommand{\cu}{{\rm B}}
\newcommand{\npm}{{\bf W}_{1,p}^{\mu}}
\newcommand{\npa}{{\bf W}_{\alpha,p}^{\mu}}
\newcommand{\npma}{{\bf W}_{\frac{1}{p},p}^{\mu}}
\newcommand{\npmat}{{\bf \widetilde{M}}_{1,p}^{\mu}}
\newcommand{\ww}{{\bf W}_{1,p}^{\mu}}
\newcommand{\npmt}{{\bf \widetilde{M}}_{1,p}^{\mu}}
\newcommand{\npf}{{\bf M}_{1,p}^{f}}
\newcommand{\incl}{\omp \CC \ompp \CC \Omega}
\newcommand{\ce}{x_{\rm B}}
\newcommand{\ca}{u_{\rm B}}
\newcommand{\cav}{v_{\rm B}}
\newcommand{\trif}[1] {\textnormal{\rif{#1}}}
\newcommand{\tref}[1] {\textnormal{\ref{#1}}}
\newcommand{\ro}{\varrho}
\newcommand{\qo}{{\rm Q}_{\deo{h}}}
\newcommand{\qoh}{{\rm Q}_{\deo{h}/2}}
\newcommand{\qof}{{\rm Q}_{\deo{h}/4}}
\newcommand{\sco}{c(n,\gamma_1,\gamma_2,\nu,L)}
\newcommand{\scos}{c(n,\gamma_1,\gamma_2,\nu,L,\sigma)}
\newcommand{\qom}{Q_{8R_0}}
\newcommand{\td}{\widetilde{d}}
\newcommand{\ratiol}{\tl/\nu}
\newcommand{\VVd}{|V(Du)-V(Dv)|^2}
\newcommand{\VVq}{|V(Du)-V(Dv)|^{\frac{2(p-1)}{p}}}
\newcommand{\QQ}{2B_0}

\DeclareMathOperator*{\argmin}{arg\min}
\DeclareMathOperator*{\argmax}{arg\max}


    \newcommand{\SL}{\mathcal L^{1,p}( D)}
\newcommand{\Lp}{L^p( Dega)}
\newcommand{\CO}{C^\infty_0( \Omega)}

\newcommand{\R}{\mathbb R}

\newcommand{\Hn}{\mathbb H^n}
\newcommand{\N}{\mathbb N}
\newcommand{\aB}{\alpha B}
\newcommand{\eps}{\epsilon}
\newcommand{\BVX}{BV_X(\Omega)}
\newcommand{\p}{\partial}
\newcommand{\bG}{\mathbb{G}}
\newcommand{\bg}{\mathfrak g}
\newcommand{\Bux}{\mbox{Box}}
\newcommand{\til}{\tilde}
\newcommand{\nuX}{\boldsymbol{\nu}^X}
\newcommand{\bN}{\boldsymbol{N}}
\newcommand{\nh}{\nabla^H}
\newcommand{\Gp}{G_{D,p}}
\newcommand{\n}{\boldsymbol \nu}
\newcommand{\subsol}{$A$-subsolution of \eqref{theeq} }
\newcommand{\supsol}{$A$-supersolution of \eqref{theeq} }
\newcommand{\sol}{solution to \eqref{theeq} }
\newcommand{\CC}{Carnot-Carath\'eodory }
\newcommand{\wolff}{\mathbf{W}_{p,\X }^{\mu}}
\newcommand{\X}{\mathbf{X}}
%%%%%%%% Bibliography
\newtoks\by
\newtoks\paper
\newtoks\book
\newtoks\jour
\newtoks\yr
\newtoks\pages
\newtoks\vol
\newtoks\publ
\def\et{ \& }
\def\name[#1, #2]{#1 #2}
\def\ota{{\hbox{\bf ???}}}
\def\cLear{\by=\ota\paper=\ota\book=\ota\jour=\ota\yr=\ota
\pages=\ota\vol=\ota\publ=\ota}
\def\endpaper{\the\by, \textit{\the\paper},
{\the\jour} \textbf{\the\vol} (\the\yr), \the\pages.\cLear}
\def\endbook{\the\by, \textit{\the\book},
\the\publ, \the\yr.\cLear}
\def\endpap{\the\by, \textit{\the\paper}, \the\jour.\cLear}
\def\endproc{\the\by, \textit{\the\paper}, \the\book, \the\publ,
\the\yr, \the\pages.\cLear}



\renewcommand{\d}{\, \mathrm{d}} %differential

\begin{document}
\title[Invariant Laws and Fokker-Planck Analysis of Adam-Type Dynamics]{Fokker-Planck Analysis and Invariant Laws for a Continuous-Time Stochastic Model of Adam-Type Dynamics}





\address{Kaj Nystr\"{o}m\\Department of Mathematics, Uppsala University\\
S-751 06 Uppsala, Sweden}
\email{kaj.nystrom@math.uu.se}









\author{Kaj Nystr{\"o}m}\thanks{The author was partially supported by grant 2022-03106 from the Swedish research council (VR)}
\date{\today}
\maketitle
\begin{abstract}
We develop an effective continuous-time model for the long-term dynamics of adaptive stochastic optimization, focusing on bias-corrected Adam-type methods. Starting from a finite-sum setting, we identify a canonical scaling of the learning rate, decay parameters, and gradient noise that gives rise to a coupled, time-inhomogeneous stochastic differential equation for the parameters $x_t$, the first-moment tracker $z_t$, and the per-coordinate second-moment tracker $y_t$. The bias-correction mechanism is retained in the limit through explicit time-dependent coefficients, while the dynamics becomes asymptotically time-homogeneous. We analyze the associated Fokker-Planck equation and under mild regularity
and dissipativity assumptions on the objective function $f$ we prove existence and uniqueness
of invariant measures for the limiting diffusion. Noise propagation is encoded by the matrix $A(x)=\mathrm{Diag}(\nabla f(x))\,H_f(x)$, where $H_f(x)$ denotes the Hessian of $f$, and hypoellipticity of the system may fail on the closed set $\mathcal D_A
\times\mathbb R^m\times(\mathbb R_+)^m$, where $$\mathcal D_A:=
\{x\in\mathbb R^m:\ e_j^\top A(x)=0
\ \mbox{ for some } j\}\subset \{x\in\R^m:\det A(x)=0\}:=\mathcal D_A^\dagger.$$  The critical points of $f$ are contained in $\mathcal D_A$. We conclude that  $\mathcal D_A^\dagger\neq \mathbb R^m$ and we use this to prove that the Markov semigroup $\mu_0P_t$ associated to time-homogeneous system converges exponentially fast, independently of initial distribution $\mu_0$,  to the unique invariant measure. The proof is based on a Harris-type argument, combined with
a minorization condition on Lyapunov sublevel sets, explicit constructions of control skeletons, and the hypoellipticity on $(\mathbb R^m\setminus \mathcal D_A)\times\mathbb R^m\times (\mathbb R_+)^m$. Our results provide a transparent continuous-time perspective on the long-time behavior of Adam-type dynamics.














\medskip





\noindent
{\em Mathematics Subject Classification (2020).} 60H10, 65C30, 37A50, 35Q84, 90C15, 68T07.



\medskip

\noindent\textit{Keywords and phrases:}
Adam optimizer, continuous-time limit, stochastic differential equations, Fokker-Planck equations, invariant measures, stochastic gradient descent (SGD), hypoellipticity, Malliavin calculus, Stroock-Varadhan support theorem, control theory, nonconvex optimization, loss landscape exploration, machine learning, deep neural networks.

\end{abstract}


%\newpage


%\setcounter{tocdepth}{1}

%{\tableofcontents}







    \setcounter{equation}{0} \setcounter{theorem}{0}






\section{Introduction}

The training of deep neural networks, as well as numerous other parametric models in machine learning, can be formulated as the finite-sum optimization problem
\begin{equation}\label{minima}
x^\star \in \operatorname*{arg\,min}_{x \in \mathbb{R}^m}\biggl\{ f(x) := \frac{1}{n} \sum_{i=1}^n f_{i}(x) \biggr\},
\end{equation}
where $x \in \mathbb{R}^m$ denotes the vector of model parameters, $f:\mathbb{R}^m \to \mathbb{R}$ is the empirical risk (loss function), and $f_i:\mathbb{R}^m \to \mathbb{R}$ corresponds to the loss incurred on the $i$-th data point. In typical applications, $f$ is continuously differentiable but non-convex, and $f \geq 0$.

When $n$ is large, evaluating $\nabla f(x)$ exactly is computationally expensive. A standard approach is therefore to employ \emph{stochastic gradient methods}, in which one replaces the full gradient by an unbiased estimator computed from a small random subset (minibatch) of the data. The simplest such method is \emph{stochastic gradient descent} (SGD), which updates the parameter vector in the negative direction of the stochastic gradient, scaled by a prescribed step size.

An important refinement is \emph{SGD with momentum}, in which the update direction is obtained by an exponential moving average of past stochastic gradients. This modification introduces a velocity variable into the dynamics and can be interpreted as a first-order discretization of a damped dynamical system. Empirically, momentum often improves both the speed and stability of convergence relative to vanilla SGD.

Beyond momentum, a major line of research has focused on \emph{adaptive gradient methods}. The first widely used algorithm in this class is \emph{Adagrad} \cite{Duchi2011,McMahanStreeter2010}, which adapts the step size coordinatewise using accumulated squared gradients. While this approach can significantly improve performance in problems with sparse or small gradients, it suffers from an ever-decreasing effective learning rate, which may be detrimental in nonconvex or high-dimensional settings. To address this issue, several refinements have been proposed, including \emph{RMSProp} \cite{TielemanHinton2012}, \emph{Adadelta} \cite{Zeiler2012}, and \emph{Nadam} \cite{Dozat2016}. These methods replace cumulative sums by exponential moving averages, thereby limiting memory to recent gradients. Despite strong empirical performance, such methods may fail to converge in settings where rare but informative gradients are rapidly down-weighted.

A further extension, widely adopted in modern deep learning, is the \emph{Adam} optimizer (ADaptive Moment Estimation) \cite{KingmaBa2015}. Adam maintains exponential moving averages of both stochastic gradients and their coordinatewise squares, yielding adaptive estimates of the first and second moments. These are bias-corrected to compensate for initialization effects. The resulting update can be interpreted as a preconditioned stochastic gradient step, in which each coordinate is scaled inversely proportionally to the estimated standard deviation of its stochastic gradient. Owing to its robustness and ease of tuning, Adam has become one of the most widely used optimization methods in large-scale applications, particularly in computer vision and natural language processing. We refer the reader to~\cite{E2020,Ruder2016,Sun2019} and~\cite{Bach2024,GarrigosGower2023,JentzenIntroDL2023} for further background.



From a mathematical perspective, many stochastic optimization algorithms can be viewed as discretizations of stochastic differential equations (SDEs) with state-dependent drift and diffusion. This continuous-time viewpoint has proved fruitful for analyzing stability, convergence, invariant measures, and long-time behavior via tools from stochastic analysis and partial differential equations, in particular through associated Fokker-Planck equations.



The purpose of this paper is to develop and analyze a continuous-time stochastic framework for adaptive optimization algorithms, with a particular focus on bias-corrected Adam-type methods. Starting from a finite-sum setting, we identify a canonical scaling of the learning rate, decay parameters, and gradient noise that gives rise to a coupled, time-inhomogeneous SDE for the parameter $x_t$, a first-moment tracker $z_t$, and a second-moment tracker $y_t$. In the long-time regime, this system becomes asymptotically time-homogeneous. We study the associated Fokker-Planck equations and establish the existence and uniqueness of invariant measures, thereby contributing to a continuous-time understanding of the stochastic behavior induced by bias correction and adaptivity.



Our focus on Adam reflects both its central role in large-scale learning and gaps in its theoretical understanding. Early worst-case examples demonstrated divergence even for convex objectives, prompting the development of variants with convergence guarantees \cite{Reddi2018,KingmaBa2015}. More recent work shows that \emph{vanilla} Adam can converge under suitable conditions and parameter choices \cite{Zhang2023}. Nevertheless, important open questions remain, including the role of bias correction at finite horizons, convergence rates beyond convex or Polyak-Lojasiewicz regimes, robustness under heavy-tailed or state-dependent gradient noise, the structure of invariant measures induced by coordinatewise preconditioning, and metastability near saddle points in high dimensions. Even in deterministic settings, sharp convergence rates for adaptive methods have only recently been obtained \cite{DereichJentzenRiekert2025}. These developments motivate further continuous-time analysis, with the long-term goal of placing Adam alongside SGD, momentum methods, and RMSProp within a unified analytical framework.



While the literature on continuous-time analysis of Adam-type methods remains limited, a key contribution is~\cite{malladi2022sdes}, where It\^o SDE approximations are derived and used to study batch-size scaling. Building on scaling rules for SGD \cite{goyal2017accurate,li2019stochastic,li2021validity}, the authors show that adaptive methods obey a \emph{square-root} scaling rule, supported by theory and simulations. Our work relates to and complements these results by focusing on the long-time behavior of the resulting stochastic dynamics. In particular, we analyze the associated Fokker-Planck equations, establish existence and uniqueness of invariant measures, and study convergence to equilibrium. To our knowledge, this provides a new and self-contained analytical framework for the long-time behavior of Adam-type dynamics.

\section{The {Adam} optimization algorithm} Consider the optimization problem in \eqref{minima} and suppose that we initialize the parameter vector at $x_0 \in \mathbb{R}^m$ at time $t=0$. We discretize time using a step size $h>0$, and for any stochastic process $s(\cdot)$ we denote by $s_k := s(kh)$ its value at discrete time $t_k = kh$. The \emph{learning rate} is denoted by $\eta$ and it controls the size of each update in the parameter space.


 In  the \emph{Adam} optimization algorithm \cite{KingmaBa2015}, the discrete-time iteration can be written as \begin{align}\label{eqn:sgd_main+}
x_{k+1}^i &= x_k^i - \eta\,\frac{z_{k+1}^i}{\sqrt{y_{k+1}^i} + \varepsilon},\notag\\
z_{k+1} &= \frac{1-\alpha^{k}}{1-\alpha^{k+1}}\,\alpha z_{k}
+ \frac{1-\alpha}{1-\alpha^{k+1}} \bigl( \nabla f(x_k) + \xi_k \bigr), \notag\\
y_{k+1}^i &= \frac{1-\beta^{k}}{1-\beta^{k+1}}\,\beta y_{k}^i
+ \frac{1-\beta}{1-\beta^{k+1}} \bigl| \partial_{x_i} f(x_k) + \xi_k^i \bigr|^2,
\end{align}
for $k \in \{0,1,\dots,K-1\}$ and $i \in \{1,\dots,m\}$. The parameters satisfy $\alpha,\beta \in (0,1)$ and $\varepsilon>0$ is included to avoid division by zero in the denominator of the update. The stochastic term
\begin{align}
\xi_k(x) := \nabla \tilde{f}_{k}(x) - \nabla f(x)
\end{align}
represents the \emph{stochastic gradient noise}, i.e., the discrepancy between the true gradient $\nabla f(x)$ and a stochastic gradient $\nabla \tilde{f}_{k}(x)$ computed, for instance, from a minibatch of the data.

The variables in \eqref{eqn:sgd_main+} are interpreted as follows: $x_k\in\mathbb{R}^m$ denotes the parameter vector at iteration $k$; $z_k\in\mathbb{R}^m$ is the \emph{first–moment estimate} (momentum), i.e., an exponentially weighted moving average (EWMA) of past stochastic gradients; and $y_k\in\mathbb{R}^m$ is the per–coordinate \emph{second–moment estimate}, an EWMA of squared stochastic gradients. The algorithmic parameters are: $\alpha\in(0,1)$, the decay rate for the first–moment estimate (Adam’s $\beta_1$), where values close to $1$ give longer memory; $\beta\in(0,1)$, the decay rate for the second–moment estimate (Adam’s $\beta_2$), with values close to $1$ producing smoother variance tracking; $\varepsilon>0$, a small constant ensuring numerical stability in the preconditioner; and $\xi_k$, the stochastic gradient noise, typically modeled as zero–mean with covariance depending on the minibatch size and the data distribution.


The prefactors ${(1-\alpha)}/{(1-\alpha^{k+1})}$ and ${(1-\beta)}/{(1-\beta^{k+1})}$ act as \emph{bias-correction factors}. To make their role precise, it is convenient to distinguish between the uncorrected exponential moving averages and their bias-corrected counterparts used in Adam. Let $\tilde z_k$ denote the uncorrected first-moment estimate, defined by
\[
\tilde z_{k+1} = \alpha \tilde z_k + (1-\alpha)\bigl(\nabla f(x_k)+\xi_k\bigr),
\]
and define the bias-corrected quantity $z_k = {\tilde z_k}/{(1-\alpha^k)}$. An analogous definition applies to the second-moment estimate $y_k$. Without correction, the uncorrected averages $\tilde z_k$ and $\tilde y_k$ retain a dependence on their initialization, leading to a bias in early iterations whenever the initial values do not match the true moments. This effect is particularly pronounced when $\alpha$ (or $\beta$) is close to~1, since the influence of the initialization then decays only slowly. To illustrate this, consider for simplicity the noiseless and stationary case where $\nabla f(x_k)=g$ is constant. The recursion $\tilde z_{k+1} = \alpha \tilde z_k + (1-\alpha) g$ with general initialization $\tilde z_0$ yields
\[
\tilde z_k = \alpha^k \tilde z_0 + (1-\alpha^k)\, g,
\qquad
\mathbb{E}[\tilde z_k] = \alpha^k \tilde z_0 + (1-\alpha^k)\, g.
\]
Hence,
\[
\mathbb{E}[\tilde z_k] - g = \alpha^k (\tilde z_0 - g),
\]
so the bias arises from the mismatch between the initialization and the true first moment, and decays at rate $\alpha^k$. In the standard choice $\tilde z_0=0$, this reduces to $
\mathbb{E}[\tilde z_k]=(1-\alpha^k)\,g$, which underestimates $g$ for small $k$. The bias-corrected quantity is $z_k = {\tilde z_k}/{(1-\alpha^k)}$  and in the standard initialization $\tilde z_0=0$ this yields $z_k = g$ in the stationary setting, i.e., an unbiased estimator of the first moment. For general initialization, substituting the expression for $\tilde z_k$ gives
\[
z_k = g + \frac{\alpha^k}{1-\alpha^k}\,\tilde z_0,
\]
so that the dependence on $\tilde z_0$ persists through a transient term. However, this term decays exponentially fast as $k \to \infty$, and hence $z_k \to g$. Thus, while the bias correction removes the bias associated with zero initialization, the influence of general initial conditions vanishes asymptotically.  An analogous argument applies to the second-moment estimate $y_k$. Consequently, the bias-correction factors in Adam should be understood as compensating for the bias induced by the standard zero initialization, while more generally the influence of initialization appears as a transient effect that vanishes exponentially fast. Under stationarity or slowly varying gradients and assuming $\mathbb{E}[\xi_k]=0$, the corrected quantities $z_k$ and $y_k$ are therefore approximately unbiased estimators of the first and second moments.







In practical applications, the parameters $\alpha$ and $\beta$ are chosen close to $1$ so that the exponential averaging underlying $z_k$ and $y_k$ retains long-term memory of past gradients and squared gradients, respectively. The most common default values, used for example in the original Adam implementation~\cite{KingmaBa2015}, are $\alpha = 0.9$ and $\beta = 0.999$, which correspond to characteristic memory scales of order $(1-\alpha)^{-1} \approx 10$ steps for the first-moment estimate and $(1-\beta)^{-1} \approx 1000$ steps for the second-moment estimate. These values have been found to work well across a wide range of problems, but they may be tuned depending on the desired trade-off between adaptivity and responsiveness. Smaller values of $\alpha$ or $\beta$ yield faster adaptation to changes in the gradient statistics, while values closer to $1$ provide stronger smoothing.



\section{Heuristic Scaling and Effective Continuous-Time Modeling of Adam-Type Dynamics}\label{heur}

The Adam iteration \eqref{eqn:sgd_main+} can be equivalently written in increment form as
\begin{align}\label{eqn:sgd_main++}
x_{k+1}^i - x_k^i &= -\eta\,\frac{z_{k+1}^i}{\sqrt{y_{k+1}^i} + \varepsilon},\notag\\
z_{k+1} - z_k &= \frac{1-\alpha}{1-\alpha^{k+1}} \Bigl[ -z_k + \bigl( \nabla f(x_k) + \xi_k \bigr) \Bigr], \notag\\
y_{k+1}^i - y_k^i &= \frac{1-\beta}{1-\beta^{k+1}} \Bigl[ -y_k^i + \bigl| \partial_{x_i} f(x_k) + \xi_k^i \bigr|^2 \Bigr],
\end{align}
for $k \in \{0,1,\dots,K-1\}$ and $i \in \{1,\dots,m\}$. This representation is well suited for deriving continuous-time models. To pass formally to the limit as the time-step $h \to 0$, one must (i) introduce a continuous-time interpolation of the discrete iterates, (ii) scale the algorithmic parameters appropriately, and (iii) specify a consistent model for the stochastic gradient noise. Different choices lead to qualitatively different limiting dynamics, ranging from deterministic ODEs to stochastic differential equations (SDEs).

\smallskip
\noindent\emph{Interpolation.}
Let $t_k:=kh$ and define the piecewise-constant (càdlàg) interpolations
\[
X^h(t):=x_k,\quad Z^h(t):=z_k,\quad Y^h(t):=y_k,\qquad t\in[t_k,t_{k+1}).
\]
Then $(X^h,Z^h,Y^h)$ are stochastic processes indexed by continuous time.

\smallskip
\noindent\emph{Learning rate.}
To obtain finite drift terms in the limit, it is natural to set
\[
\eta = \gamma h, \qquad \gamma>0.
\]
This ensures that the increments of $x_k$ are of order $h$, leading to a nontrivial continuous-time evolution. Other scalings either suppress the dynamics ($\eta=o(h)$) or cause it to diverge ($\eta\gg h$).

\smallskip
\noindent\emph{Decay parameters.}
If $\alpha,\beta$ are kept fixed as $h \to 0$, the effective memory windows $(1-\alpha)^{-1}$ and $(1-\beta)^{-1}$ collapse on the time scale $t=kh$, and the moving averages lose their temporal structure. To retain a nontrivial memory, one must therefore let $\alpha,\beta \to 1$ as $h \to 0$. A natural scaling is
\begin{equation}\label{scalla}
\alpha = 1 - a h, \qquad \beta = 1 - b h, \qquad a,b>0,
\end{equation}
for which
\begin{equation}\label{scalla+}
\alpha^{k} \to e^{-at}, \qquad \beta^{k} \to e^{-bt}.
\end{equation}
Accordingly,
\begin{equation}\label{scalla++}
\frac{1-\alpha}{1-\alpha^{k+1}} \sim hc_a(t),
\quad
\frac{1-\beta}{1-\beta^{k+1}} \sim hc_b(t),\qquad\mbox{where}\qquad c_a(t):=\frac {a}{1-e^{-a t}},\quad c_b(t):=\frac {b}{1-e^{-b t}},
\end{equation}
yielding exponential memory kernels in the limit.

\smallskip
\noindent\emph{Noise and limiting regimes.}
The limiting behavior depends critically on the scaling of the stochastic gradient noise $\xi_k$. Indeed, under \eqref{scalla},
\begin{align}\label{Klara1}
z_{k+1}-z_k
&=
\frac{1-\alpha}{1-\alpha^{k+1}}
\Bigl[-z_k+\nabla f(x_k)+\xi_k\Bigr]
\;\sim\;
h\,c_a(t_k)\Bigl[-z_k+\nabla f(x_k)+\xi_k\Bigr],\notag\\
y_{k+1}^i-y_k^i
&=
\frac{1-\beta}{1-\beta^{k+1}}\Bigl[-y_k^i+\big|\partial_{x_i}f(x_k)+\xi_k^i\big|^2\Bigr ]
\;\sim\;
h\,c_b(t_k)\Bigl[-y_k^i+\big|\partial_{x_i} f(x_k)+\xi_k^i\big|^2\Bigr ],\end{align}
and
\begin{align}\label{Klara2}
\bigl|\partial_{x_i}f(x_k)+\xi_k^i\bigr|^2
=
(\partial_{x_i}f(x_k))^2
+
2\,\partial_{x_i}f(x_k)\,\xi_k^i
+
(\xi_k^i)^2.
\end{align}
Hence the stochastic contribution enters through the term $h\xi_k$. The scaling of this quantity determines the nature of the limiting dynamics.

\medskip
\noindent
\textbf{(i) Deterministic (ODE) limit.}
Under the central-limit scaling
\[
\xi_k = \sqrt{h}\,\sigma\,\zeta_k, \qquad \zeta_k \sim N(0,I_m),
\]
one has $h\xi_k = \mathcal{O}(h^{3/2})$, so the stochastic contribution vanishes in the limit. This yields the deterministic, time-inhomogeneous ODE
\begin{align} \label{eq:cts-xode}
\mathrm{d}x_t^i &= -\gamma\,\frac{z_t^i}{\sqrt{y_t^i}+\varepsilon}\,\mathrm{d}t,\quad
\mathrm{d}z_t^i = c_a(t)\bigl(\partial_{x_i}f(x_t)-z_t^i\bigr)\,\mathrm{d}t, \quad
\mathrm{d}y_t^i = c_b(t)\bigl(-y_t^i + (\partial_{x_i}f(x_t))^2 \bigr)\,\mathrm{d}t,
\end{align}
describing averaged dynamics.

\medskip
\noindent
\textbf{(ii) Stochastic (SDE) limit.}
To obtain a nontrivial diffusion term, the random increments $h\xi_k$ must be of order $\sqrt{h}$, since this is the scaling that accumulates to Brownian motion over $1/h$ steps. This requires
\begin{align}\label{klara3}
h\xi_k \sim \sqrt{h},
\qquad\text{that is,}\qquad
\xi_k &= \frac{\sigma}{\sqrt{h}}\,\zeta_k.
\end{align}
Under this scaling, the increments of $z_k$ contain fluctuations of order $\sqrt{h}$, which converge to Brownian motion in the limit, and one formally deduces
\begin{align}
\mathrm{d}z_t^i &= c_a(t)\bigl(\partial_{x_i}f(x_t)-z_t^i\bigr)\,\mathrm{d}t
                 + c_a(t)\sigma\,\mathrm{d}B_t^i. \label{eq:cts-zapa}
\end{align}
 However, inserting this scaling into \eqref{Klara2} leads to a divergence of order $1/h$. Thus, a direct SDE limit of the full Adam recursion is not available without additional modeling assumptions.






\medskip
\noindent
\textbf{(iii) Effective stochastic model.}
To obtain a well-posed stochastic limit under the scaling \eqref{klara3}, one must carefully analyze the second-moment dynamics. Consider the decomposition
\[
h\bigl|\partial_{x_i}f(x_k)+\xi_k^i\bigr|^2
=
h(\partial_{x_i}f(x_k))^2
+
2h\,\partial_{x_i}f(x_k)\,\xi_k^i
+
h(\xi_k^i)^2.
\]
Under the scaling $\xi_k^i=\sigma h^{-1/2}\zeta_k^i$, with $\zeta_k^i$ centered and of unit variance, the three terms exhibit markedly different behavior. The first term is of order $h$ and contributes to the drift in the limit. The cross term is centered and of order $\sqrt{h}$, and hence produces fluctuations that, in principle, accumulate to a stochastic integral. In contrast, the quadratic term satisfies
\[
h(\xi_k^i)^2 = \sigma^2(\zeta_k^i)^2,
\]
which is of order one at each step. Consequently, in the exact scheme, this term produces contributions of order one per step, and thus diverges over $O(h^{-1})$ steps. This shows that, under the strong noise scaling required to obtain a nontrivial diffusion in the $z$–equation, the second-moment recursion does not admit a finite continuous-time limit. To extract a meaningful limiting dynamics, one must therefore modify the second-moment equation at the level of modeling. Writing
\[
(\zeta_k^i)^2 = 1 + \bigl((\zeta_k^i)^2 - 1\bigr),
\]
we decompose
\[
h(\xi_k^i)^2
=
\sigma^2
+
\sigma^2\bigl((\zeta_k^i)^2 - 1\bigr).
\]
The first term corresponds to a deterministic contribution, while the second is centered. When accumulated over time, the centered fluctuations are of lower order (their variance is of order $h$) and vanish in the limit, whereas the mean contributes a finite drift. Similarly, the cross term
\[
2\,\partial_{x_i}f(x_k)\,\xi_k^i
\]
is centered and would formally give rise to an additional multiplicative noise term in the limiting equation for $y_t$. However, this term introduces state-dependent fluctuations and does not contribute to the drift. Moreover, the resulting diffusion coefficient does not vanish at $y_t^i=0$, so the limiting process would not preserve positivity, and $y_t^i$ could become negative with positive probability. This is undesirable, both because $y_t^i$ represents a second-moment quantity and because the update involves $\sqrt{y_t^i}$. Motivated by these observations, we adopt an effective closure and replace the exact second-moment update by the approximation
\begin{align*}
\bigl|\partial_{x_i}f(x_k)+\xi_k^i\bigr|^2
\;\leadsto\;
(\partial_{x_i}f(x_k))^2 + \sigma^2.
\end{align*}
This can be interpreted as an averaging or moment-closure procedure, in which the second-moment dynamics retain the mean-square effect of the noise while discarding fast and state-dependent fluctuations. Under this closure, the prefactor $\theta_k^{(b,h)}={(1-\beta)}/{(1-\beta^{k+1})}\sim h\,c_b(t_k)$ yields contributions of order $h$, leading to a finite drift term in the limit. The resulting continuous-time model thus captures the effective variance induced by stochastic gradients while remaining well posed and analytically tractable.









\medskip
\noindent
\textbf{Canonical scaling.}
Motivated by the above discussion, we adopt the scaling
\begin{align}\label{scale}
\eta = \gamma h, \qquad
\alpha = 1 - a h, \qquad
\beta = 1 - b h, \qquad
\xi_k = \frac{\sigma}{\sqrt{h}}\,\zeta_k,
\end{align}
with $a,b,\gamma,\sigma>0$ fixed and $\{\zeta_k\}$ i.i.d.\ standard Gaussians. In addition, we impose the effective closure
\begin{align}\label{scale+}
\bigl|\partial_{x_i}f(x_k)+\xi_k^i\bigr|^2
\;\leadsto\;
(\partial_{x_i}f(x_k))^2+\sigma^2,
\end{align}
as described above.


\begin{remark}
The scaling $\xi_k=\sigma h^{-1/2}\zeta_k$ should not be interpreted as a literal model of minibatch noise. Rather, it is an effective scaling ensuring that the stochastic increments in the $z_k$-equation are of order $\sqrt{h}$ and hence converge to Brownian motion. In this sense, the resulting SDE captures macroscopic fluctuations of the algorithm rather than the microscopic noise at the discrete level.
\end{remark}

\begin{remark}
The deterministic ODE limit corresponds to a law-of-large-numbers approximation capturing averaged optimization dynamics. In contrast, the SDE model incorporates stochastic effects and is therefore more suitable for studying invariant measures, metastability, and escape phenomena. The SDE should thus be viewed as an effective model rather than a direct limit of the original discrete scheme.
\end{remark}

\begin{remark}
More general noise models may be considered, including state-dependent covariances, heavy-tailed distributions, or temporal correlations. While such extensions may better reflect practical training dynamics, they significantly complicate the analysis. The isotropic Gaussian setting adopted here provides a tractable yet representative framework.
\end{remark}

\begin{remark}
Time-dependent learning rates, such as $\eta_t \propto t^{-1/2}$, can also be incorporated and lead to non-autonomous continuous-time dynamics. The present scaling instead focuses on the refinement limit $h \to 0$, thereby isolating the intrinsic structure of the algorithm.
\end{remark}

\begin{remark}
The continuous-time system derived in this work should be viewed as an intermediate description between the exact discrete-time Adam iteration and purely heuristic stochastic differential equation models. On the one hand, it is rooted in a systematic scaling analysis of the discrete algorithm. On the other hand, it incorporates an effective closure in the second-moment dynamics to obtain a well-posed and analytically tractable limit. In this sense, the resulting SDE captures the essential structural features of Adam-type methods, momentum, adaptivity, and bias correction, while providing a simplified framework for studying their long-time stochastic behavior.
\end{remark}

\subsection{The continuous-time limit} We now formalize the preceding discussion by stating the continuous-time limit in the form of Theorem \ref{Thm1} below. Theorem \ref{Thm1} is proved in Appendix
\ref{aaapp}. The result is derived under a set of assumptions on the objective function $f$. These conditions are specified and discussed in Subsection~\ref{Subcond}. By construction, the scaling laws and approximation in \eqref{scale}-\eqref{scale+} yield, in the small-step limit, a time-inhomogeneous SDE system for $(x_t,z_t,y_t)$ with smooth coefficients. The limit dynamics retain the finite memory of the exponential moving averages as well as the bias-correction mechanism.

In the following,  $\mathbb{D}([0,T];\mathbb{R}^m)$ denotes the Skorokhod space of
$\mathbb{R}^m$-valued càdlàg (right-continuous with left limits) functions
on $[0,T]$, equipped with the standard topology. Additional notation is introduced in Subsection~\ref{Prelim}.




\begin{theorem}\label{Thm1}
Assume that $f$ satisfies condition \textnormal{(A1)} from Subsection~\ref{Subcond}, and fix $\varepsilon>0$. Let $t_k:=kh$, consider the scaling laws and closure approximation in \eqref{scale}-\eqref{scale+}, and define the piecewise-constant interpolations
\[
X^h(t):=x_k,\qquad Z^h(t):=z_k,\qquad Y^h(t):=y_k,
\qquad t\in[t_k,t_{k+1}),
\]
with initial condition
\[
(X^h(0),Z^h(0),Y^h(0))=(x_0,z_0,y_0)\in\mathbb R^{3m}.
\]
Assume that $(x_0,z_0,y_0)$ is deterministic (or, more generally, independent of the noise sequence $\{\zeta_k\}$). Let $B_t=(B_t^1,\dots,B_t^m)$ be an $m$-dimensional Brownian motion with independent components. Then, for every $0<\delta<T<\infty$, the interpolated processes $(X^h,Z^h,Y^h)$ converge in law in $\mathbb{D}([\delta,T];\mathbb{R}^{3m})$ as $h\to0$ to a continuous process $(x_t,z_t,y_t)_{t\in[\delta,T]}$, which is the unique strong solution on $[\delta,T]$ of the system
\begin{align}
\mathrm{d}x_t^i &= -\gamma\,\frac{z_t^i}{\sqrt{y_t^i}+\varepsilon}\,\mathrm{d}t, \label{eq:cts-x}\\
\mathrm{d}z_t^i &= c_a(t)\bigl(\partial_{x_i}f(x_t)-z_t^i\bigr)\,\mathrm{d}t
                 + c_a(t)\sigma\,\mathrm{d}B_t^i, \label{eq:cts-z}\\
\mathrm{d}y_t^i &= c_b(t)\bigl(-y_t^i + (\partial_{x_i}f(x_t))^2 + \sigma^2\bigr)\,\mathrm{d}t, \label{eq:cts-y}
\end{align}
for $i=1,\dots,m$, with initial condition at time $\delta$ given by the law of the limit of $(X^h(\delta),Z^h(\delta),Y^h(\delta))$.  The time-dependent coefficients $c_a(t)$ and $c_b(t)$ are defined in \eqref{scalla++}. Moreover, the limiting process is uniquely defined for all $t>0$, and the convergence holds on every compact interval $[\delta,T]\subset(0,\infty)$.
\end{theorem}




\begin{remark}
The continuous-time system derived in Theorem~\ref{Thm1} should be interpreted as an effective stochastic model for Adam-type dynamics rather than as the literal scaling limit of the exact discrete recursion. The convergence result relies on the moment-closure approximation introduced in the second-moment equation. The relevance of the limiting system lies in the fact that it preserves the defining structural features of Adam: momentum through the variable $z$, adaptive preconditioning through $y$, and bias correction through the time-dependent coefficients $c_a(t)$ and $c_b(t)$. In the limit, stochasticity enters only through the momentum variable, while the second-moment tracker evolves deterministically with an additional drift $\sigma^2$ induced by gradient noise. Thus, the resulting system provides a tractable continuous-time representation of Adam-type dynamics that captures its essential mechanisms while enabling the study of long-time stochastic behavior.
\end{remark}



\begin{remark}\label{suppa}
Note that  at the boundary $y_t^i=0$ we have
\[
\dot y_t^i = c_b(t)\bigl((\partial_{x_i}f(x_t))^2+\sigma^2\bigr) \ge 0,
\]
so the nonnegative orthant $[0,\infty)^m$ is forward invariant.
\end{remark}


\begin{remark}
Since $c_a(t),c_b(t)\sim 1/t$ as $t\downarrow 0$, the limiting system is singular at $t=0$.
In particular, when $\sigma\neq 0$, the process $z_t$ does not admit a finite limit as $t\downarrow 0$, and the solution cannot in general be extended to $t=0$. The limit process is therefore naturally defined on $(0,T]$, and the convergence in the theorem holds on every interval $[\delta,T]$ with $\delta>0$. The family of laws is consistent as $\delta\downarrow 0$, but does not in general define a process with a finite initial condition at $t=0$.
\end{remark}
\begin{remark}
It is natural to ask whether an appropriate choice of initial data can remove the singularity at $t=0$. If $\sigma=0$ and the initial data satisfy
\[
z_0^i=\partial_{x_i}f(x_0),\qquad
y_0^i=(\partial_{x_i}f(x_0))^2,\qquad i=1,\dots,m,
\]
then the singularity at $t=0$ is removable, and the limiting system may be extended to a well-defined solution on $[0,T]$. If $\sigma\neq 0$, however, the stochastic term in the $z$-equation is too singular near $t=0$, and in general no finite initial condition yields a solution starting from $t=0$ in the usual sense.
\end{remark}












\section{Long-term behaviour and invariant measures}\label{Invar}
Having established Theorem~\ref{Thm1}, the rest of the paper is devoted to the long-time behavior of the
inhomogeneous system \eqref{eq:cts-x}-\eqref{eq:cts-y}, and to its time-homogeneous  counterpart. In particular, since $c_a(t)\to a$ and $c_b(t)\to b$ as $t\to\infty$, standard perturbation arguments for SDEs
(see Appendix \ref{LLD}) imply that the asymptotic dynamics of the system \eqref{eq:cts-x}-\eqref{eq:cts-y} is captured by the
time-homogeneous limit
\begin{align}
\mathrm{d} x_t^i &= -\gamma\,\frac{ z_t^i}{\sqrt{ y_t^i}+\varepsilon}\,\mathrm{d}t, \label{eq:cts-x+}\\
\mathrm{d} z_t^i &= a\bigl(\partial_{x_i}f( x_t)- z_t^i\bigr)\,\mathrm{d}t + a\sigma\,\mathrm{d}B_t^i, \label{eq:cts-z+}\\
\mathrm{d} y_t^i &= b\bigl(- y_t^i + (\partial_{x_i}f( x_t))^2 + \sigma^2\bigr)\,\mathrm{d}t, \label{eq:cts-y+}
\end{align}
for $i=1,\dots,m$. Assuming that $f$ satisfies \textnormal{(A1)} stated in Subsection~\ref{Subcond},
it follows that the system in  \eqref{eq:cts-x+}-\eqref{eq:cts-y+}
have unique strong solutions $(x_t, z_t, y_t)_{t\ge0}$ for every initial condition
$(x_0, z_0, y_0)$ with $y_0\in[0,\infty)^m$. We are interested in the long-time behaviour of the system in \eqref{eq:cts-x+}-\eqref{eq:cts-y+}  and in particular in the existence and uniqueness of invariant measures, and as well as the rate of convergence to equilibrium.






In our study of \eqref{eq:cts-x+}-\eqref{eq:cts-y+} we consider initial conditions $(x_0,z_0,y_0)$ drawn from a probability
distribution $\mu_0$ on $\R^{3m}$ supported on $\bar E$, $E:=\R^m\times\R^m\times(0,\infty)^m$.  Unless otherwise stated, we assume
\begin{equation}\label{intimeasure}
\mu_0\big(\bar E\big)=1
\quad\text{and}\quad
\mathbb{E}_{\mu_0}\!\big [V(x_0,z_0,y_0)\big]<\infty,
\end{equation}
where $V$ is a Lyapunov function constructed in Section~\ref{sec:lya}.  We will refer to such $\mu_0$ as an \emph{admissible initial probability distribution} on $\R^{3m}$. As proved in Section~\ref{sec:lya},
\begin{equation}\label{intimeasureapa}
\mathbb{E}_{\mu_0}\!\big [V(x_0,z_0,y_0)\big]<\infty\Leftrightarrow \iiint_E (\|x\|^2+\|z\|^2+\|y\|_1)\,\mu_0(\mathrm d x\mathrm d z\mathrm d y)<\infty.
\end{equation}

The law corresponding to $(x_t,z_t,y_t)$ in \eqref{eq:cts-x+}-\eqref{eq:cts-y+}, with initial distribution $\mu_0$, is
denoted by $\pi_t := \mu_0 P_t$, where $(P_t)_{t\ge0}$ is the Markov semigroup associated with
\eqref{eq:cts-x+}–\eqref{eq:cts-y+}. I.e.,  given the state space
$E$,
\begin{equation}\label{eq:pi_t_def}
\pi_t(A)
\;=\;
(\mu_0 P_t)(A)
\;=\;
\iiint_E P_t(x,z,y,A)\,\mu_0(\mathrm d x\mathrm d z\mathrm d y),
\qquad
A\in\mathcal B(E).
\end{equation}

As $y_0^i\ge 0$, we have $y_t^i\ge (1-e^{-bt})\sigma^2>0$ for all $t>0$, and the support of $\pi_t$ is contained in
$\R^m\times \R^m\times \bigl((1-e^{-bt_0})\sigma^2,\infty\bigr)^m$ for all $t> t_0$. In particular, if $t> t_\ast:=(\ln 2)/b$, then the support of $\pi_t$ is contained in
$\R^m\times \R^m\times \bigl(\sigma^2/2,\infty\bigr)^m$. We introduce the sets
\begin{equation}\label{intimeasurell}
E_{\eta}:=\R^m\times\R^m\times(\eta,\infty)^m
\end{equation}
for all $\eta\geq 0$ and we identify $E=E_0=\R^m\times\R^m\times(0,\infty)^m$.





We prove the following result concerning the existence and uniqueness of invariant measures, and the  convergence to equilibrium.  The total variation distance $\|\cdot\|_{\mathrm{TV}}$
and the $2$-Wasserstein distance $W_2$ are defined in
Section~\ref{Prelim}.

\begin{theorem}\label{CorC}
Assume that $f$ satisfies \textnormal{(A1)} and \textnormal{(A2)} from Subsection~\ref{Subcond}. Fix $a,b,\gamma,\sigma>0$ and $\varepsilon>0$,
and let $(P_t)_{t\ge0}$ denote the Markov semigroup associated with
\eqref{eq:cts-x+}-\eqref{eq:cts-y+}. Then there exists a unique invariant probability measure $\pi_\infty$ on $E$, and for every admissible initial distribution $\mu_0$ we have
\[
\mu_0 P_t \Rightarrow \pi_\infty
\qquad \text{as } t\to\infty .
\]
In particular, the limiting distribution does not depend on $\mu_0$. Moreover, there exist constants $C<\infty$, $\lambda>0$, and $t_\ast>0$,
independent of the admissible initial distribution $\mu_0$, such that
\begin{equation}\label{eq:V-geometric-ergodicity-revisedapa}
\|\mu_0P_t-\pi_\infty\|_{\mathrm{TV}}+\bigl (W_2(\mu_0P_t,\pi_\infty)\bigr)^2
\le
C e^{-\lambda (t-t_\ast)}
\bigl(1+\mathbb{E}_{\mu_0}\!\big[V(x_{t_\ast},z_{t_\ast},y_{t_\ast})\big]\bigr),
\qquad t\ge t_\ast.
\end{equation}
\end{theorem}



\begin{corollary}\label{correa}
Assume the hypotheses of Theorem~\ref{CorC}. Let $\pi_\infty$ be the unique invariant probability measure associated to homogeneous system in \eqref{eq:cts-x+}-\eqref{eq:cts-y+} established in Theorem~\ref{CorC}. Let $X_t=(x_t,z_t,y_t)$ denote the unique strong solution of the inhomogeneous system \eqref{eq:cts-x}-\eqref{eq:cts-y}, started at time $\delta>0$ with admissible initial distribution $\mu_\delta$ on $E$. Then
\begin{equation}\label{converga}
\lim_{t\to\infty}\mathbb{P}_{\mu_\delta}(X_t\in A)
=
\pi_\infty(A),
\qquad \forall\, A\in\mathcal{B}(E).
\end{equation}
\end{corollary}



\begin{remark}
Corollary~\ref{correa} is included to emphasize that the study of invariant measures for the system \eqref{eq:cts-x+}-\eqref{eq:cts-y+} provides a natural framework for understanding the long-term behavior of the inhomogeneous system \eqref{eq:cts-x}-\eqref{eq:cts-y}. At the same time, within the present analysis, Corollary~\ref{correa} follows directly from Theorem~\ref{CorC}, see Lemma~\ref{thm:same-equilibrium} in the appendix.
\end{remark}
\begin{remark}
Although we say that $\pi_\infty$ is an invariant probability measure on $\R^{3m}$, recall that its support lies in
$\R^m\times\R^m\times (\sigma^2/2,\infty)^m$. On this set, assuming \textnormal{(A1)},
the coefficients of \eqref{eq:cts-x+}–\eqref{eq:cts-y+} are $C^\infty$.
\end{remark}



To put Theorem~\ref{CorC} into perspective we note that while the existence of invariant probability measures for dissipative stochastic
dynamical systems can often be established by relatively soft arguments,
the \emph{uniqueness}  and quantitative rates of convergence
to equilibrium are typically considerably more delicate.
Uniqueness generally requires some form of irreducibility or mixing,
which in turn depends on subtle properties of how the noise propagates
through the system.
Establishing exponential convergence further demands a quantitative version
of this mixing, often formulated in terms of minorization,
Doeblin-type conditions, or suitable coupling or Lyapunov-Harris structures.

An important analytical tool in the study of the ergodic properties of SDEs and Markov processes
is the \emph{infinitesimal generator}
together with its the adjoint and the associated
Fokker-Planck (forward Kolmogorov) equation. It is a general principle that if every probability solution of the
Fokker-Planck equation possesses a continuous strictly positive density
with respect to Lebesgue measure, then the corresponding diffusion process
admits at most one invariant probability measure.
For systems with smooth coefficients, the existence of such a density is ensured
either by uniform ellipticity of the diffusion matrix or,
more generally, by Hörmander's bracket condition and the theory of hypoellipticity.










While hypoellipticity ensures the existence of a smooth transition
density,  establishing positivity of
the density is also central. This typically relies on the Stroock-Varadhan support
theorem, see \cite{StroockVaradhan1972,StroockVaradhan1979}, which connects the SDE with its associated skeleton ODE and the reachable sets. In the case of a global strong Hörmander condition (and hence global hypoellipticity),
frequently assumed in the literature, this issue is largely unproblematic
thanks to the Chow-Rashevskii theorem.


In our case, the {infinitesimal generator} ${\mathcal L}$ associated to \eqref{eq:cts-x+}-\eqref{eq:cts-y+}, see
\eqref{generator}, contains a  lower order drift term and Hörmander's bracket condition (taking iterative commutators between the diffusion directions and the first order vector field defined by the lower order terms in ${\mathcal L}$) fails
on $\mathcal D_A\times\mathbb R^m\times(\mathbb R_+)^m$. Here  $\mathcal D_A$ is the set defined in \eqref{matt+++} below and this set contains the ``flat'' directions of $f$ where the $y$-variables decouple from
the noisy coordinates, see Subsection~\ref{Hypoellipticity}. We believe that this lack of global hypoellipticity,  and the subsequent need to construct suitable
control paths,  make Theorem~\ref{CorC} and its proof non-trivial and interesting.

\subsubsection{Proof of Theorem~\ref{CorC}} In the context of quantitative mixing and minorization, the notions of \emph{petite} and \emph{small} sets play a central role, see Subsection~\ref{small}
for the precise definitions. In our case, a set $C\in\mathcal B(E)$ can be
shown to be small if there exist $t_0>0$, $\eta>0$, and a set $W\in\mathcal B(E)$ with finite,
positive Lebesgue measure, such that
\[
\inf_{(x,z,y)\in C} P_{t_0}\bigl((x,z,y),A\bigr) \geq \eta {\operatorname{Leb}(A \cap W)}/{\operatorname{Leb}(W)},\qquad \forall A\, \in\mathcal B(E).
\]
 In particular, there exists a time $t_0>0$ such that the process
$(x_t,z_t,y_t)$, when started from any initial condition in $C$, assigns
a uniformly positive probability to the set $W$, with a lower bound
independent of the initial point in $C$.




Having constructed the Lyapunov function in Section~\ref{sec:lya}, we show in
Theorem~\ref{Thm3-revised-} below that the proof of Theorem~\ref{CorC} can be reduced to verifying that the compact sets $\{C_R\}$, $C_R=\hat C_R\cap E_{\sigma^2/2}$, where \(\hat C_R\) denotes a level set of the Lyapunov function, are small. In the literature, see for instance
\cite{MattinglyStuartHigham2002,HerzogMattingly2015} and the references
therein, the standard route to proving that a set \(C\) is small
is, as mentioned above, through hypoellipticity together with the Stroock-Varadhan support
theorem. In particular, the construction of suitable
control paths plays a central role. Consequently, the  main themes in the proof of Theorem~\ref{CorC} is overcoming the lack of global hypoellipticity,  and the construction of control paths.




In the Lie bracket analysis of ${\mathcal L}$ associated to \eqref{eq:cts-x+}-\eqref{eq:cts-y+}, the matrix
\begin{equation}\label{matt}
A(x):=\mathrm{Diag}(\nabla f(x))\,H_f(x),
\end{equation}
where $H_f(x)$ denotes the Hessian of $f$, arises naturally. As shown in Subsection~\ref{Hypoellipticity},
hypoellipticity for our system and ${\mathcal L}$ may fail on the closed set
\begin{equation}\label{matt+}
\mathcal D_H
:=
\mathcal D_A
\times\mathbb R^m\times(\mathbb R_+)^m,
\end{equation}
where \begin{equation}\label{matt+++}
\mathcal D_A:=
\{x\in\mathbb R^m:\ e_j^\top A(x)=0
\ \mbox{ for some } j\}.
\end{equation}
Therefore, as
\begin{equation}\label{matt+++a}
\{x\in\mathbb R^m:\ \partial_{x_j}f(x)=0
\ \mbox{ for some } j\}\subset\mathcal D_A,
\end{equation}
hypoellipticity may fail at the critical points of $f$.

To compensate for the lack of global hypoellipticity we also work with the set
\begin{equation}\label{dagga}
\mathcal D_A^\dagger:=\{x\in\R^m:\ \det A(x)=0\}.
\end{equation}
Since a matrix with a zero row is necessarily singular, we always have the inclusion
\begin{equation}\label{Inclu}
\mathcal D_A \subset \mathcal D_A^\dagger.
\end{equation}
In general the inclusion in \eqref{Inclu} is strict as a matrix may have no zero row and still fail to have full rank. $\mathcal D_A$ detects \emph{coordinate degeneracy} (loss of at least one forcing direction),
while $\mathcal D_A^\dagger$ detects \emph{rank degeneracy} (failure of the forcing directions
to span $\R^m$). Equivalently,
\[
x\notin\mathcal D_A
\quad\Longleftrightarrow\quad
\text{no row of }A(x)\text{ vanishes},
\]
whereas
\[
x\notin\mathcal D_A^\dagger
\quad\Longleftrightarrow\quad
A(x)\text{ is invertible}.
\]
The latter condition  guarantees,
via the inverse function theorem, local openness of the map
\begin{equation}\label{aaaa+}
x\mapsto\big((\partial_{x_1}f(x))^2,\dots,(\partial_{x_m}f(x))^2\big),
\end{equation}
a map which is present in the dynamics of $y$.

An important insight is that if $f$ satisfies \textnormal{(A1)} and \textnormal{(A2)}, then
\begin{equation}\label{Remma1+}
\mathcal D_A^\dagger \neq \mathbb R^m.
\end{equation}
This is a consequence of an argument based on topological degree and Sard’s theorem, see Lemma \ref{Redu}  below. To prove Theorem~\ref{CorC} we exploit \eqref{Remma1+}  to construct suitable control paths. Indeed,  since \(\mathcal D_A^\dagger\) is closed, \eqref{Remma1+} implies the
existence of an open ball
\[
B(x_\ast,r)\subset \mathbb R^m\setminus\mathcal D_A^\dagger .
\]
In the control argument we then first steer points in \(C_R\) into an open and bounded set
\[
B(x_\ast,r)\times Z\times Y,
\]
on which the Hörmander bracket condition holds.
In fact, this step relies only on the weaker condition
\(B(x_\ast,r)\subset\mathbb R^m\setminus\mathcal D_A\). Second, we exploit the stronger condition
\(B(x_\ast,r)\subset\mathbb R^m\setminus\mathcal D_A^\dagger\)
in a core controllability argument underlying the proof of
Theorem~\ref{CorC}.
This controllability argument consists of a construction of a controlled
ODE skeleton, which allows us to apply the Stroock-Varadhan support
theorem and deduce positivity of the transition density between certain
sets.
A key difficulty in this construction is to control the trajectory of the
ODE in the \(y\)-variables and here the local openness of the map
\eqref{aaaa+} plays a crucial role. In this way we show, using the  geometric nondegeneracy
encoded in \(\mathcal D_A^\dagger\neq \R^m\), that the diffusion is sufficiently
irreducible to guarantee smoothing, a minorization condition, and hence
exponential ergodicity.






















\subsection{Organization of the rest of the paper} In Section~\ref{Prelim} we collect notation, state the standing assumptions on \(f\), we introduce the Fokker-Planck equation and we state the Stroock-Varadhan support theorem in our context. In Subsection~\ref{Hypoellipticity} we discuss hypoellipticity of the infinitesimal generator and Fokker-Planck equations associated to \eqref{eq:cts-x+}-\eqref{eq:cts-y+}.

Existence and uniqueness of invariant measures for SDEs are classical topics, e.g., see the monographs \cite{Hasminskii2012,MeynTweedie2009,BogachevKrylovRoeckner2015,BakryGentilLedoux2014,Villani2009}
and the foundational works \cite{Doob1953,StroockVaradhan1979,Hormander1967,LindvallRogers1986,Eberle2016,Eberle2019,HairerMattingly2011b,MattinglyStuartHigham2002}.
A common approach is via a \emph{Foster–Lyapunov} drift condition formulated in terms of a Lyapunov function \(V\).
Section~\ref{sec:lya} is devoted to the construction of such a function \(V\) adapted to our dynamics under the assumptions on \(f\).


In  Section~\ref{aaaa} we state and prove  Theorem~\ref{Thm3-revised-}. Theorem~\ref{Thm3-revised-} is an auxiliary result which reduces the proof of Theorem \ref{CorC} to verifying that
the compact sets $\{C_R\}$ are small. The proof of Theorem~\ref{Thm3-revised-} is based on a
continuous-time Harris-Meyn-Tweedie argument.   In general, several standard approaches are available for proving uniqueness and
ergodicity of invariant measures. If the semigroup $(P_t)_{t\ge0}$ is strong
Feller and the process is irreducible, then there exists at most one
invariant probability measure \cite{Doob1953}. Strong Feller regularity
typically follows under uniform ellipticity or from Hörmander’s theorem
\cite{Hormander1967}, while irreducibility can be established using support
theorems \cite{StroockVaradhan1979}. In situations where neither uniform ellipticity nor a direct application of
Hörmander’s condition is available  uniqueness can
still be obtained, together with geometric ergodicity, by combining a Lyapunov
drift condition with a minorization (small set) condition and invoking
Harris’ theorem \cite{MeynTweedie2009,HairerMattingly2011b}. We follow this
approach in the proof of Theorem~\ref{Thm3-revised-} as we in the statement of the theorem assume that the sets $\{C_R\}$ are small.


Theorem~\ref{CorC} is proved in Section~\ref{sec:proof-CorC} where we prove that the sets $\{C_R\}$ are small. Once this is done we can invoke Theorem~\ref{Thm3-revised-}. Verifying that the sets $\{C_R\}$ are small is perhaps the most involved part of the paper due to the delicate construction of controlled
ODE skeletons, which allows us to apply the Stroock-Varadhan support
theorem to deduce positivity of the transition density between certain
sets. In our case, the lack of global hypoellipticity makes the construction of suitable
control paths central and non-trivial.

 Section \ref{Conc}  is devoted to some concluding remarks and topics for future research.

At the very end of the paper we have include three appendices, Appendix \ref{aaapp}-\ref{HMI}.
Appendix \ref{aaapp} contains the proof of Theorem~\ref{Thm1}. In Appendix \ref{LLD} we give the details concerning the long-term dynamics for the inhomogeneous system \eqref{eq:cts-x}-\eqref{eq:cts-y}, and Appendix \ref{HMI} is devoted to higher moments and integrability for solutions to the homogeneous system in \eqref{eq:cts-x+}-\eqref{eq:cts-y+}. For experts, there are no surprises in these appendices, and the arguments are fairly standard. The material is included for completeness but could perhaps be omitted.




\section{Preliminaries}\label{Prelim}
Throughout, $\|\cdot\|$ denotes the standard Euclidean $\ell^2$-norm on $\mathbb{R}^d$, while $\|\cdot\|_p$ denotes the $\ell^p$-norm. Sets will in general be open and in particular $B(\cdot,\rho)\subset\R^m$ will denote an open ball in $\R^m$.

For two probability measures $\mu,\nu$ on a measurable state space $(E,\mathcal B(E))$, the
\emph{total variation distance} is defined by
\[
\|\mu-\nu\|_{\mathrm{TV}}
:= \sup_{A \in \mathcal B(E)} |\mu(A)-\nu(A)|
= \frac12 \sup_{\|f\|_\infty \le 1} \left| \int_E f \,\mathrm d\mu - \int_E f \,\mathrm d\nu \right|.
\]
If $\mu,\nu$ are probability measures on $\R^{3m}$ with finite second moments, their $2$-Wasserstein distance is defined by
\[
W_2(\mu,\nu)
:= \inf_{\pi \in \Pi(\mu,\nu)}
\left( \iint_{\R^{3m}\times\R^{3m}} \|u-v\|^2 \,\pi(\mathrm{d}u,\mathrm{d}v) \right)^{1/2},
\]
where $\Pi(\mu,\nu)$ denotes the set of all \emph{couplings} of $\mu$ and $\nu$, i.e.\ probability measures
$\pi$ on $\R^{3m}\times\R^{3m}$ with marginals $\mu$ and $\nu$,
\[
\pi(A\times \R^{3m}) = \mu(A), \qquad
\pi(\R^{3m}\times B) = \nu(B),
\quad A,B \subseteq \R^{3m}\ \text{measurable}.
\]
Intuitively, a coupling describes a joint distribution of two random variables with given laws $\mu$ and $\nu$, and the Wasserstein distance quantifies the minimal expected transport cost needed to move one distribution into the other.





\subsection{Assumptions on $f$} \label{Subcond}
Our main results are established under structural assumptions on the objective function $f$ to be minimized. These conditions are standard in optimization and stochastic analysis, but we recall them here for clarity.

\medskip
\noindent\textbf{(A1) Global $L_f$-smoothness.}
\begin{equation*}
f \in C^2(\R^m)\, (f \in C^\infty(\R^m))
\quad\text{and}\quad
\|\nabla f(x)-\nabla f(\bar x)\| \le L_f \|x-\bar x\|
\quad \forall\, x,\bar x \in \R^m.
\end{equation*}
Equivalently, the Hessian is uniformly bounded
\[
\|\nabla^2 f(x)\|_{\mathrm{op}} \le L_f, \qquad \forall x \in \R^m.
\]
Global smoothness controls the curvature of $f$ and rules out excessively steep growth at infinity. In particular,  $f$ can grow at most \emph{quadratically} in $\|x\|$. More precisely, for any $x \in \R^m$,
\[
f(x) \;\le\; f(0) + \nabla f(0)\cdot x + \frac{L_f}{2}\|x\|^2,
\]
so $f(x) = O(\|x\|^2)$ as $\|x\|\to\infty$. Functions with super-quadratic growth (e.g.\ $f(x) = \|x\|^4$) are not globally $L$-smooth, since their Hessians become unbounded at infinity. $f \in C^2(\R^m)$ is sufficient for this discussion but when discussing hypoellipticity the assumption  $f \in C^\infty(\R^m)$ is the correct one.

\begin{remark}
While (A1) is a strong condition, it greatly simplifies the analysis of the continuous-time dynamics. Frequently one can replace global Lipschitz continuity of the gradient by \emph{local} Lipschitz continuity together with standard linear growth bounds. This weaker setting still ensures existence and uniqueness of solutions to the limiting SDEs but complicates the ergodic analysis.
\end{remark}


\noindent\textbf{(A2$''$) Global $m_f$–strong convexity.}
\begin{equation}\label{eq:dissipativity-}
\exists\, m_f>0 \ \text{ such that }\
\langle \nabla f(x)-\nabla f(\bar x),\,x-\bar x\rangle
\;\ge\; m_f \|x-\bar x\|^2,
\quad \forall\, x,\bar x\in\R^m.
\end{equation}
If $f\in C^2$, this is equivalent to a uniform positive lower bound on the Hessian,
\[
\nabla^2 f(x) \ \succeq\ m_f I_m, \qquad \forall x\in\R^m,
\]
so that $f$ is \emph{$m_f$–strongly convex}.
Strong convexity excludes flat directions and saddle points. It implies that $f$ grows \emph{at least quadratically} at infinity as
\[
f(x) \;\ge\; f(0) + \frac{m_f}{2}\|x\|^2 - C,
\]
for some constant $C$. Thus $f(x)\to\infty$ as $\|x\|\to\infty$, with quadratic lower growth.

\smallskip
\noindent\textbf{(A2$'$) Global dissipativity (one–sided coercivity).}
There exist constants $m_\infty>0$ and $C_\infty\ge 0$ such that
\begin{equation}\label{eq:dissipativity}
\langle \nabla f(x),\,x\rangle
\;\ge\; m_\infty \|x\|^2 - C_\infty,
\qquad \forall x\in\R^m.
\end{equation}
This condition is weaker than strong convexity as it only requires $f$ to dominate a quadratic \emph{at infinity}, without ruling out nonconvexity in bounded regions.
In particular, (A2$'$) guarantees that $f$ is \emph{coercive}, i.e.\ $f(x)\to\infty$ as $\|x\|\to\infty$, so level sets of $f$ are compact.

\smallskip
\noindent\textbf{(A2) Dissipativity at infinity.}
There exist constants $m_f>0$, $c\ge0$, and $R\ge0$ such that
\begin{equation}\label{D}
\langle \nabla f(x),\,x\rangle
\;\ge\; m_f \|x\|^2 - c,
\qquad \text{for all }\|x\|\ge R.
\end{equation}
This condition only enforces quadratic growth for sufficiently large $\|x\|$.
It allows for nonconvex behavior (e.g.\ saddle points, flat regions, multiple local minima) within a bounded domain, while still ensuring that $f(x)\to\infty$ as $\|x\|\to\infty$.




\begin{remark}
If $f\in C^2$ is both \emph{$L_f$-smooth} (A1) amd \emph{strongly convex} (A2$''$), then $f$ satisfies the classical $(L_f,m_f)$-condition.
\end{remark}
\begin{remark}
The growth conditions form a hierarchy
\[
\text{(A2$''$)} \;\implies\; \text{(A2$'$)} \;\implies\; \text{(A2)}.
\]
In particular, (A2$''$) enforces global quadratic curvature (uniform convexity), (A2$'$) ensures quadratic growth at infinity but permits local nonconvexity, (A2) is the weakest condition, requiring only eventual quadratic growth with no restriction on bounded regions.
\end{remark}

\begin{remark} If $f$ satisfies \textnormal{(A2)}, the (pathological)
situation that $f$ is independent of one or more of the coordinates $(x_1,..,x_m)$ can  not occur.
\end{remark}

While the following result may not be stated in this exact form in the literature, it follows from standard arguments combining topological degree and Sard’s theorem.

\begin{lemma}\label{Redu}
Assume that $f$ satisfies \textnormal{(A1)} and \textnormal{(A2)}, and let $\mathcal D_A^\dagger$ be defined as in \eqref{dagga}. Then
\[
\mathcal D_A^\dagger \neq \mathbb R^m.
\]
\end{lemma}

\begin{proof}
Assume, for contradiction, that $\mathcal D_A^\dagger=\mathbb R^m$. Then
\[
0=\det A(x)
=\det(\Diag(\nabla f(x)))\det(H_f(x))
=\Bigl(\prod_{i=1}^m \partial_{x_i} f(x)\Bigr)\det(H_f(x))
\qquad \forall x\in\mathbb R^m.
\]
Set $F:=\nabla f:\mathbb R^m\to\mathbb R^m$. By \textnormal{(A2)}, there exists $R>0$ such that
\[
\langle F(x),x\rangle >0 \qquad \text{for all }x\in \partial B(0,R).
\]
In particular, $F(x)\neq 0$ on $\partial B(0,R)$. Moreover, for every $y\in B(0,\epsilon)$, $\epsilon$ sufficiently small, we also have
\[
\langle F(x)-y,x\rangle >0 \qquad \text{for all }x\in \partial B(0,R),
\]
and hence $F(x)-y\neq 0$ on $\partial B(0,R)$. Therefore, for all $y\in B(0,\epsilon)$, the Brouwer degree
\[
\deg(F-y,B(0,R),0):=\sum_{x\in (F-y)^{-1}(0)} \operatorname{sign}(\det DF(x))
\]
is well-defined and equals the number of solutions $x \in B(0,R)$ of $F(x)=y$, counted with multiplicity given by $\operatorname{sign}(\det DF(x))$. To compute it, consider the homotopy
\[
G_t(x):=t(F(x)-y)+(1-t)(x-y), \qquad t\in[0,1].
\]
By the above estimate, $G_t(x)\neq 0$ for all $x\in\partial B(0,R)$ and all $t\in[0,1]$, and hence the homotopy $G_t$ provides a continuous deformation between $F-y$ and $\mathrm{Id}-y$ which does not vanish on $\partial B(0,R)$, and hence the degree is preserved. Hence, by homotopy invariance of the degree,
\[
\deg(F-y,B(0,R),0)=\deg(\mathrm{Id}-y,B(0,R),0)=1,
\]
where $\mathrm{Id}(x)=x$ is the identity map. In particular, for every $y\in B(0,\epsilon)$, there exists $x\in B(0,R)$ such that $F(x)=y$. Thus a whole neighborhood of the origin is contained in $F(\mathbb R^m)$. Next, recall that a value $y\in\mathbb R^m$ is called a \emph{regular value} of $F$ if for every $x$ with $F(x)=y$, the Jacobian matrix $DF(x)$ is invertible, i.e.\ $\det DF(x)\neq 0$. Since $F\in C^1(\mathbb R^m)$, Sard's theorem implies that the set of regular values of $F$ is dense in $\mathbb R^m$. Since the set
\[
\{y\in\mathbb R^m:\ y_i\neq 0\ \text{for all }i=1,\dots,m\}
\]
is open and dense, we may choose $y$ such that  $y\in B(0,\epsilon)$, hence $y\in F(\mathbb R^m)$, such that $y_i\neq 0$ for all $i=1,\dots,m$, and such that $y$ is a regular value of $F$. Fix such a $y$, and choose $x\in\mathbb R^m$ such that $F(x)=y$, i.e.\ $\nabla f(x)=y$. Then $y_i\neq 0$ for all $i$ implies
\[
\prod_{i=1}^m \partial_{x_i} f(x)\neq 0.
\]
Since $\det A(x)=0$, it follows that
\[
\det(H_f(x))=0.
\]
On the other hand, since $y$ is a regular value of $F=\nabla f$, we must have
\[
\det DF(x)=\det(H_f(x))\neq 0,
\]
which is a contradiction. This contradiction shows that $\mathcal D_A^\dagger\neq \mathbb R^m$.
\end{proof}





\subsection{The Fokker-Planck equation}\label{FP}



A key analytical tool in the study of the ergodic properties of the system
\eqref{eq:cts-x+}–\eqref{eq:cts-y+} is its \emph{infinitesimal generator}
${\mathcal{L}}$, together with the adjoint ${\mathcal{L}}^{\!*}$ and the corresponding
Fokker-Planck (or forward Kolmogorov) equation.




For $\varphi \in C^2(E)$, $E=\mathbb{R}^m\times\mathbb{R}^m\times(\mathbb{R}_+)^m$, and $(x,z,y)\in E$,
the generator ${\mathcal{L}}$ of \eqref{eq:cts-x+}–\eqref{eq:cts-y+} is
\begin{equation}\label{generator}
\begin{aligned}
({\mathcal{L}}\varphi)(x,z,y)
&= \sum_{i=1}^m \Bigg[
\underbrace{-\gamma\,\frac{z_i}{\sqrt{y_i}+\varepsilon}}_{\text{$x$–drift}}\;\partial_{x_i}\varphi
\;+\;
\underbrace{a\big(\partial_{x_i}f(x)-z_i\big)}_{\text{$z$–drift}}\;\partial_{z_i}\varphi
\\[-0.25em]
&\hspace{6em}+\;
\underbrace{b\Big(-y_i+(\partial_{x_i}f(x))^2+\sigma^2\Big)}_{\text{$y$–drift}}\;\partial_{y_i}\varphi
\Bigg]
+\;\frac12\,a^2\sigma^2\sum_{i=1}^m \partial_{z_i z_i}^2 \varphi.
\end{aligned}
\end{equation}



Let $(x_t,z_t,y_t)$ denote the solution of
\eqref{eq:cts-x+}–\eqref{eq:cts-y+} with given initial distribution.
For each time $t \ge 0$, we denote by $\pi_t(\mathrm{d}x\,\mathrm{d}z\,\mathrm{d}y)$
the \emph{law} (probability measure) of $(x_t,z_t,y_t)$ on
$E$. That is, for any bounded measurable test function $\varphi$,
\[
\iiint \varphi(x,z,y)\,\pi_t(\mathrm{d}x\,\mathrm{d}z\,\mathrm{d}y)
\;=\;
\mathbb{E}[\varphi(x_t,z_t,y_t)].
\]
If the process admits an \emph{invariant probability measure}
$\pi_\infty(\mathrm{d}x\,\mathrm{d}z\,\mathrm{d}y)$, then
\[
\pi_\infty P_t = \pi_\infty, \qquad \forall\,t\ge0,
\]
where $(P_t)_{t\ge0}$ is the Markov semigroup generated by the system and $\mathcal{L}$.
Equivalently, if the initial law is $\pi_\infty$, then the distribution of
$(x_t,z_t,y_t)$ remains $\pi_\infty$ for all $t\ge0$. In weak (measure–theoretic) form, the invariance condition can be written as
\begin{equation}\label{weak-inv-measure}
\iiint ({\mathcal{L}}\varphi)(x,z,y)\,
\pi_\infty(\mathrm{d}x\,\mathrm{d}z\,\mathrm{d}y)
\;=\;0,
\qquad
\forall\,\varphi\in C_0^\infty(E).
\end{equation}
This characterization remains meaningful even when $\pi_\infty$ does not admit a density
with respect to Lebesgue measure. When $\pi_t$ and $\pi_\infty$ admit densities with respect to
Lebesgue measure, we write
\[
\pi_t(\mathrm{d}x\,\mathrm{d}z\,\mathrm{d}y)
= p(t,x,z,y)\,\mathrm{d}x\,\mathrm{d}z\,\mathrm{d}y,
\qquad
\pi_\infty(\mathrm{d}x\,\mathrm{d}z\,\mathrm{d}y)
= p_\infty(x,z,y)\,\mathrm{d}x\,\mathrm{d}z\,\mathrm{d}y,
\]
and $p$ and $p_\infty$ are referred to as the \emph{probability densities}
of the process and its stationary distribution, respectively.








The forward Kolmogorov (Fokker-Planck) equation governing the evolution
of the density $p(t,x,z,y)$ associated with $\pi_t$ is
\begin{equation}\label{Fokk}
\partial_t p = {\mathcal{L}}^{\!*} p,
\end{equation}
where ${\mathcal{L}}^{\!*}$ denotes the formal $\mathrm{L}^2$–adjoint of $\mathcal{L}$,
acting as
\begin{equation}\label{Fokkop}
\begin{aligned}
{\mathcal{L}}^{\!*} p
&= -\sum_{i=1}^m \partial_{x_i}\!\Bigl(-\gamma\,\frac{z_i}{\sqrt{y_i}+\varepsilon}\, p\Bigr)
   \;-\;\sum_{i=1}^m \partial_{z_i}\!\Bigl(a(\partial_{x_i}f(x)-z_i)\, p\Bigr)\\
&\quad\;-\;\sum_{i=1}^m \partial_{y_i}\!\Bigl(b\big(-y_i+(\partial_{x_i}f(x))^2+\sigma^2\big)\, p\Bigr)
   \;+\;\frac12\,a^2\sigma^2\sum_{i=1}^m \partial_{z_i z_i}^2  p.
\end{aligned}
\end{equation}
The formal adjoint relation reads
\begin{equation}\label{weak}
\iiint({\mathcal{L}}\varphi)(x,z,y)\, \psi(x,z,y)\,\mathrm{d}x\,\mathrm{d}z\,\mathrm{d}y
\;=\;\iiint \varphi(x,z,y)\, ({\mathcal{L}}^{\!*}\psi)(x,z,y)\,\mathrm{d}x\,\mathrm{d}z\,\mathrm{d}y,
\end{equation}
for smooth compactly supported test functions
$\varphi,\psi\in C_0^\infty(E)$, or, more generally, for sufficiently
regular and rapidly decaying functions so that boundary terms vanish when integrating by parts. Thus ${\mathcal{L}}^{\!*}$ governs the time evolution of the probability density
$p(t,x,z,y)$, and a stationary density $p_\infty$ satisfies ${\mathcal{L}}^{\!*} p_\infty = 0$ in the sense of distributions, that is,
\[
\iiint ({\mathcal{L}}\varphi)(x,z,y)\, p_\infty(x,z,y)\,\mathrm{d}x\,\mathrm{d}z\,\mathrm{d}y = 0,
\quad \forall\,\,\varphi\in C_0^\infty(E).
\]
This expresses the balance of drift and diffusion in equilibrium, the continuous-time
analogue of the invariant measure condition $\pi_\infty P_t = \pi_\infty$ for Markov chains.




\subsection{Hypoellipticity}\label{Hypoellipticity}

As discussed in Subsection \ref{FP}, a key analytical tool in the study of the ergodic properties of the system
\eqref{eq:cts-x+}–\eqref{eq:cts-y+} is its \emph{infinitesimal generator}
${\mathcal{L}}$, together with the adjoint ${\mathcal{L}}^{\!*}$ and the corresponding
Fokker-Planck (or forward Kolmogorov) equation. It is natural to discuss the regularity theory for ${\mathcal{L}}$ and ${\mathcal{L}}^\ast$, in particular since
these operators have degenerate features.  The smoothness of the coefficients of ${\mathcal{L}}$ and ${\mathcal{L}}^{\!*}$ is dictated by that of $f$. Based on the resemblance between ${\mathcal{L}}$, ${\mathcal{L}}^\ast$ and kinetic Fokker-Planck, it is relevant to note that ${\mathcal{L}}$ and ${\mathcal{L}}^\ast$ can be expressed in a compact form as sums of squares of vector fields with a lower order drift terms. Indeed,
\[
{\mathcal{L}} = \frac{1}{2}\sum_{i=1}^m X_i^2+X_0,
\qquad
{\mathcal{L}}^{\!*}=   \frac{1}{2}\sum_{i=1}^m X_i^2 -X_0\; + \; \text{lower-order terms} ,
\]
where the diffusion vector fields are
\[
X_i = a\sigma\,\partial_{z_i}, \qquad i=1,\dots,m,
\]
and the drift field is
\[
X_0= \Bigl(-\gamma\,\frac{z}{\sqrt{y}+\varepsilon},\;\; a(\nabla f(x)-z),\;\; b(-y+g(x))\Bigr),
\quad
g_i(x):=(\partial_{x_i}f(x))^2+\sigma^2.
\]
By Hörmander’s theorem \cite{Hormander1967}, hypoellipticity of ${\mathcal{L}}$ is equivalent to that of its adjoint ${\mathcal{L}}^\ast$. The central structural condition is the following condition where (H) stands for Hypoellipticity.
\[
\mbox{\textbf{(H)}}\quad\{X_1,\dots,X_m,X_0\} \;\;\text{and  commutators thereof span $\mathbb{R}^m\times\mathbb{R}^m\times(\mathbb{R}_+)^m$ at each point.}
\]


Hypoellipticity therefore depends on whether the Lie algebra generated by
$\{X_1,\dots,X_m\}$ and their iterated commutators with $X_0$ spans the full
tangent space. In our setting this condition may fail, in particular along
flat directions of $f$, where the $y$–dynamics becomes effectively decoupled
from the noisy $z$–variables. Consequently, it is not possible to guarantee
hypoellipticity at every point $(x,z,y)\in\mathbb{R}^m\times\mathbb{R}^m\times
(\mathbb{R}_+)^m$.


 In our case a direct computation shows that
\[
[X_i,X_0] \;=\; \frac{\gamma a\sigma}{\sqrt{y_i}+\varepsilon}\,\partial_{x_i}
\;+\; a^2\sigma\,\partial_{z_i},
\]
for each $i$, hence the first commutators already generate all $x$-directions from the noisy $z$-directions.
To produce $y$-directions, note that the $y$-component of $X_0$ depends on $x$ through $g(x)=(g_1(x),..,g_m(x))$, $g_i(x)=(\partial_{x_i}f(x))^2+\sigma^2$.
Writing $H(x) =H_f(x)= \nabla^2 f(x)$ for the Hessian, we have
\[
\frac{\partial g_j}{\partial x_i}(x) \;=\; 2\,\bigl(\partial_{x_j} f(x)\bigr)\,H_{ij}(x),
\]
so the Jacobian $Dg(x)$ has entries $2\,(\partial_{x_j} f(x))\,H_{ij}(x)$.
To compute $[[X_i,X_0],X_0]$ we set
\[
\Gamma_i(y)\;:=\;\frac{\gamma a\sigma}{\sqrt{y_i}+\varepsilon},\qquad
\Gamma_i'(y)\;=\;-\frac{\gamma a\sigma}{2}\,\frac{1}{\sqrt{y_i}\,(\sqrt{y_i}+\varepsilon)^2},
\qquad
B\;:=\;a^2\sigma,
\]
for  $i\in\{1,\dots,m\}$. Recall that $X_i=a\sigma\,\partial_{z_i}$ and write
\[
X_0=\big(X_0^x,X_0^z,X_0^y\big)
=\Big(-\gamma\,\frac{z}{\sqrt{y}+\varepsilon},\; a(\nabla f(x)-z),\; b\big(-y+g(x)\big)\Big),
\quad g_j(x)=(\partial_{x_j}f(x))^2+\sigma^2.
\]
We have already established that
\[
[X_i,X_0]\;=\;\Gamma_i(y)\,\partial_{x_i}\;+\;B\,\partial_{z_i}.
\]
A direct coordinate computation of the Lie bracket
\[
[[X_i,X_0],X_0]=[\,\Gamma_i\partial_{x_i}+B\partial_{z_i}\,,\,X_0\,]
\]
yields the following components:

\medskip
\noindent\emph{$x$-component.} For $j\neq i$ it is zero; for $j=i$,
\[
\big([[X_i,X_0],X_0]\big)^{x_i}
\;=\; -\,\frac{B\,\gamma}{\sqrt{y_i}+\varepsilon}\;-\;b\big(-y_i+g_i(x)\big)\,\Gamma_i'(y).
\]

\noindent\emph{$z$-component.} For each $j=1,\dots,m$,
\[
\big([[X_i,X_0],X_0]\big)^{z_j}
\;=\; a\,\Gamma_i(y)\,H_{j i}(x)\;-\;a\,B\,\delta_{ij}.
\]

\noindent\emph{$y$-component.} For each $j=1,\dots,m$,
\[
\big([[X_i,X_0],X_0]\big)^{y_j}
\;=\; 2\,b\,\Gamma_i(y)\,\big(\partial_{x_j} f(x)\big)\,H_{j i}(x).
\]

\medskip
\noindent
Using the expression for $[[X_i,X_0],X_0]^{\,y_j}$ we see that $[[X_i,X_0],X_0]$ has a $y_j$-component proportional to
$\big(\partial_{x_j} f(x)\big)\,H_{j i}(x)$. Hence, as $b\neq 0$ and $\Gamma_i$ is different from zero, we see that $\partial_{y_j}$ can be generated at $(x,z,y)$ if and only if
\[
\exists\, i\in\{1,\dots,m\} \quad \text{such that} \quad
(\partial_{x_j} f(x))\,H_{ji}(x)=(\partial_{x_j} f(x))\,H_{ij}(x) \neq 0.
\]
In particular, in this context the matrix
\begin{equation}\label{matt}
A(x):=\Diag(\nabla f(x))\,H_f(x)
\end{equation}
arises naturally, and consequently hypoellipticity may break down on the closed set
\begin{equation}\label{matt+}
\mathcal{D}_H:=(\R^m\times\R^m\times(\R_+)^m)\setminus\mathcal{G}_H,
\end{equation}
where
\begin{equation}\label{matt++}
\mathcal{G}_H := \{(x,z,y): e_j^\top A(x)\neq 0 \ \ \forall j=1,\dots,m\}.
\end{equation}
Equivalently, $(x,z,y)\in\mathcal{G}_H$ if and only if $A(x)$ has no zero row. The \emph{degenerate set} where Hörmander’s condition  fails is thus the closed set $\mathcal{D}_H$ introduced in \eqref{matt+}. Equivalently,  $\mathcal{D}_H$ can be expressed as
$$\mathcal{D}_H=\mathcal{D}_H^1\cup ....\cup\mathcal{D}_H^m,$$
where
\begin{align}\label{hypc}
\mathcal{D}_H^j
&:=  \{(x,z,y): e_j^\top A(x)=0\}= \left\{(x,z,y) \in \mathbb{R}^m\times\mathbb{R}^m\times(\mathbb{R}_+)^m:
 \ (\partial_{x_j} f(x)) = 0 \ \text{or} \ H_{j\cdot}(x) = 0
\right\}.
\end{align}
In particular, points where $\nabla f(x)=0$ belong to $\mathcal{D}_H$. We summarize our finding in a lemma.

\begin{lemma}\label{lemmahyp} Let $\mathcal{D}_H$, $\mathcal{G}_H$, be the set introduced in \eqref{matt+} and \eqref{matt++}, respectively. Then, on  $\mathcal{G}_H$ the Lie algebra generated by $\{X_0,X_1,\dots,X_m\}$ spans all coordinate
directions and hence condition (H) is satisfied. In particular,  ${\mathcal{L}}$ and ${\mathcal{L}}^{\!*}$ are hypoelliptic on $\mathcal{G}_H$.
\end{lemma}




\begin{remark}
One may attempt to recover the missing $y$–directions on $\mathcal D_H$
by considering higher-order Lie brackets. However, a direct computation
shows that the resulting expressions become increasingly complicated
and still vanish at points where $\nabla f(x)=0$ and $z=0$.
Consequently, Hörmander's condition cannot in general be recovered
through higher commutators at such points. In particular,
hypoellipticity cannot be guaranteed on $\mathcal D_H$.
\end{remark}









\begin{remark}
Failure of Hörmander’s condition at a point does not necessarily destroy smoothness of the invariant law.
Indeed, Hörmander’s theorem provides a \emph{sufficient} local criterion for $C^\infty$ regularization of transition
densities, but it is not necessary. If Hörmander’s condition holds on an open set and the process is irreducible,
then the invariant density $p_\infty$ is absolutely continuous and $C^\infty$ on that set. Regularity may still
persist across points where the bracket condition fails. Genuine singularities typically arise only from
structural degeneracies of the dynamics, for instance when flat directions of $f$ disconnect the $y$–dynamics
from the noisy $z$–variables.
\end{remark}




\subsection{The support theorem of Stroock-Varadhan}   For $T>0$, the controlled ODE (skeleton system) associated to \eqref{eq:cts-x+}–\eqref{eq:cts-y+} and starting at $(x_0,z_0,y_0)$ is
\begin{equation}\label{eq:skeleton-components}
\left\{
\begin{aligned}
\dot x^i(t) &= -\gamma\,\frac{z^i(t)}{\sqrt{y^i(t)}+\varepsilon},\\
\dot z^i(t) &= a\big(\partial_{x_i}f(x(t)) - z^i(t)\big) + a\sigma\,h_i(t),\\
\dot y^i(t) &= b\big(-y^i(t) + (\partial_{x_i}f(x(t)))^2 + \sigma^2\big),
\end{aligned}\right.
\qquad i=1,\dots,m
\end{equation}
where $(x(0),z(0),y(0))=(x_0,z_0,y_0)$ and  $h=(h_1,\dots,h_m)\in \mathrm L^2([0,T];\R^m)$ is a control.
Given $T$ and $U_0:=(x_0,z_0,y_0)$ we let
$$\mathcal R_T(U_0)=\bigl\{(x(T),z(T),y(T)):\,\mbox{$h=(h_1,\dots,h_m)\in \mathrm L^2([0,T];\R^m)$ is a control}\bigr\}.$$
$\mathcal R_T(U_0)$ is the set of all possible endpoints $(x(T),z(T),y(T))$ for the system in \eqref{eq:skeleton-components}, when starting at $U_0$ and using controls $h\in \mathrm L^2([0,T];\R^m)$ for which the solution exists on $[0,T]$.

In our context the support theorem of Stroock-Varadhan \cite{StroockVaradhan1979} can be stated as follows.

\begin{theorem}\label{thm:support} Assume that $f$ satisfies \textnormal{(A1)}  from Subsection~\ref{Subcond}.
Fix $a,b,\gamma,\sigma>0$ and $\varepsilon>0$, and let $(P_t)_{t\ge0}$ denote the Markov semigroup
associated with the time-homogeneous system \eqref{eq:cts-x+}–\eqref{eq:cts-y+}. Fix $T>0$ and $U_0\in E$. Then the topological support of the transition kernel $P_T(U_0,\cdot)$ equals the closure of the skeleton reachable set, i.e.,
\[
\mathrm{supp}\,P_T(U_0,\cdot) \;=\; \overline{\mathcal R_T(U_0)}.
\]
In particular, if $\mathcal R_T(U_0)$ contains a nonempty open set $W\subset E$, then $P_T(U_0,W)>0$.
\end{theorem}


\begin{remark}
Note that the Itô and Stratonovich formulations of \eqref{eq:cts-x+}-\eqref{eq:cts-y+} coincide, since the diffusion coefficient is state-independent and the Itô–Stratonovich correction term vanishes.
\end{remark}

\begin{remark}\label{thm:supportrem}
Consider vector fields $\{X_1,\dots,X_m,X_0\}$ on $\mathbb{R}^N$.
In the literature, see for instance the books \cite{AgrachevSachkov2004,Jurdjevic1997}, the following connectivity hypothesis (here labeled Hypothesis~[CH]) frequently appears:


\noindent For every $(x,t),(y,s)\in \mathbb{R}^{N+1}$ with $t>s$, there exists an absolutely continuous path $\gamma : [0,t-s]\to \mathbb{R}^N$
such that
\begin{equation}\label{eq:hypothesis-C}
\begin{cases}
\displaystyle
\dot{\gamma}(\tau)
=
\sum_{k=1}^m \omega_k(\tau)\,X_k(\gamma(\tau))
+
X_0(\gamma(\tau)),
\\[0.5em]
\gamma(0)=x, \qquad \gamma(t-s)=y,
\end{cases}
\end{equation}
with $\omega_1,\dots,\omega_m \in L^\infty([0,t-s])$.

\noindent The relevance of Hypothesis~[CH] in the probabilistic setting is clarified
by the Stroock-Varadhan support theorem which describes the support of the law of the diffusion associated
with $\{X_1,\dots,X_m,X_0\}$ as the closure of the set of endpoints of
solutions to the controlled ODE \eqref{eq:hypothesis-C}. Importantly, the support theorem itself does not require any Hörmander-type
bracket condition, nor does it assume Hypothesis~[CH].
It applies under general regularity assumptions on the coefficients and
provides a control-theoretic characterization of the support. Hypothesis~[CH] is instead a global controllability condition ensuring that
the reachable set of the controlled system is large (in particular, that any
point can be reached in positive time).
When it holds, the support theorem implies that the diffusion has full
topological support.
\end{remark}
\begin{remark}\label{thm:supportrem+}
 In contrast to controllability/support theorems, Hörmander-type bracket conditions are used to establish
hypoellipticity, yielding smoothness of transition
densities, but they do not in general imply global accessibility. It is well known that the condition
\begin{equation}\label{eq:rank-spatial}
\operatorname{rank}\,\mathrm{Lie}\{X_1,\dots,X_m\}(x)
=
N,
\qquad \forall x\in \mathbb{R}^{N},
\end{equation}
implies Hypothesis~[CH] of Remark \ref{thm:supportrem}. This is usually referred to as the \emph{strong Hörmander condition}. If one instead requires that
\begin{equation}\label{eq:rank-full}
\operatorname{rank}\,\mathrm{Lie}\{X_1,\dots,X_m,X_0\}(x)
=
N,
\qquad \forall x\in \mathbb{R}^{N},
\end{equation}
then one speaks of the \emph{weak Hörmander condition}.
In general, the weak Hörmander condition does not imply Hypothesis~[CH].
\end{remark}


\subsection{Small sets}\label{small}

In the context of quantitative mixing and minorization, the notions of \emph{petite} and \emph{small} sets play a central role.
A set $C \in \mathcal B(E)$ is called \emph{small} if there exist constants
$t_0 > 0$, $\varepsilon > 0$, and a probability measure $\varphi$
on $(E,\mathcal B(E))$ such that for all $(x,z,y) \in C$ and all
$A \in \mathcal B(E)$,
\begin{align}\label{petta-}
P_{t_0}((x,z,y),A) \ge \varepsilon\,\varphi(A).
\end{align}
In this case, $\varphi$ is called a \emph{minorizing measure} for $C$. To prove that \(C\) is a small set, it is sufficient to show that there exist
a time \(t_0>0\), a constant \(\eta>0\), and a set \(W\in\mathcal B(E)\) with finite, positive Lebesgue measure such that
\begin{align}\label{petta-+}
P_{t_0}((x,z,y),A)
\ge
\eta\,\frac{\operatorname{Leb}(A\cap W)}{\operatorname{Leb}(W)},
\qquad
\forall (x,z,y)\in C,\ \forall A\in\mathcal B(E).
\end{align}
Indeed, if \eqref{petta-+} holds, then \eqref{petta-} follows by defining
\[
\varphi(A):={\operatorname{Leb}(A\cap W)}/{\operatorname{Leb}(W)},
\qquad A\in\mathcal B(E).
\]
Condition \eqref{petta-+} expresses that whenever the process starts from
\((x_0,z_0,y_0)=(x,z,y)\in C\), its law at time \(t_0\) dominates, uniformly in the initial condition, a fixed positive multiple of the normalized Lebesgue measure on \(W\).
In particular, the process reaches \(W\) with probability at least \(\eta\), uniformly over all initial points in \(C\).










\section{Construction of Lyapunov functions} \label{sec:lya}

In this section we construct the Lyapunov functions underlying much of the analysis of the paper and frequently referred to in the statement of our main results. The subsequent analysis will use the condition on $f$ stated in \textnormal{(A1)}, i.e.,
for some constant $0<L_f<\infty$, we have
\begin{equation}\label{D-}
\|\nabla f(x)-\nabla f(\bar x)\| \le L_f \|x-\bar x\|
\quad \forall\, x,\bar x \in \R^m,
\end{equation}
and the coercivity/dissipativity condition on $f$ stated in \textnormal{(A2)}, i.e., for some constants $m_f>0$, $c\ge0$, $R\ge0$, we have
\begin{equation}\label{D}
\langle \nabla f(x),x\rangle \;\ge\; m_f\|x\|^2 - c,
\qquad \text{for all }\|x\|\ge R.
\end{equation}

While we for the proof of Theorem~\ref{CorC} mainly need estimates on sets $\{E_\eta\}$ introduced in \eqref{intimeasurell}, $\eta>0$, we here prefer to make constructions
valid on all of $\mathbb R^m\times \mathbb R^m\times\mathbb R^m$. Since we develop pointwise estimates, we then have the freedom
to restrict to subsets. In line with this we will in the following construct Lyapunov functions in the context of the operator
\begin{align}\label{Regloss1}
\mathcal{L}_\epsilon \;&:=\; \sum_{i=1}^m \Bigg[{-\gamma\,\frac{z_i}{\sqrt{|y_i|}+\varepsilon}}\;\partial_{x_i}
+a\big(\partial_{x_i}f(x)-z_i\big)\;\partial_{z_i}+b\big(-y_i+(\partial_{x_i}f(x))^2+\sigma^2\big)\;\partial_{y_i}
\Bigg]\notag\\
&\quad +\;\frac12\,a^2\sigma^2\sum_{i=1}^m \partial_{z_i z_i}^2+\epsilon\sum_{i=1}^m \partial_{x_i x_i}^2+\epsilon\sum_{i=1}^m \partial_{y_i y_i}^2,
\end{align}
on $\R^m\times\R^m\times\R^m$ and where $\epsilon>0$.  By construction
\[
\mathcal{L}_\epsilon \;=\; \mathcal{L} \;+\; \epsilon\,\Delta_x\;+\; \epsilon\,\Delta_y,\qquad \epsilon>0,
\]
on $\R^m\times\R^m\times(\R_+)^m$. $\mathcal{L}_\epsilon$ is an extension of $\mathcal{L}$ to all of $\R^m\times\R^m\times\R^m$  and it is also a regularization of $\mathcal{L}$ as $\mathcal{L}_\epsilon$ is uniformly elliptic  on $\R^m\times\R^m\times\R^m$. We can  view $\mathcal L$ as a vanishing-viscosity limit of $\mathcal{L}_\epsilon$ on $\R^m\times\R^m\times(\R_+)^m$ and on this set we identify $\mathcal{L}=\mathcal{L}_0$.  At the SDE level the construction adds Brownian noise to the $y$–equation (and to the $x$-equation), destroying the positivity constraint on $y_t$ and potentially rendering the $x$–drift $(\sqrt{y_t}+\varepsilon)^{-1}$ ill–defined, therefore we replace this term with $(\sqrt{|y_t|}+\varepsilon)^{-1}$ in the definition of $\mathcal{L}_\epsilon$.

We define a candidate Lyapunov function as
\begin{equation}\label{V}
V(x,z,y)=V_\upsilon(x,z,y)
= \theta\bigl(f(x)-f_\ast\bigr) + \frac\alpha2\|z\|^2 - \beta\,x\cdot z + \delta \|y\|_{1,\upsilon}
\qquad \theta,\alpha,\beta,\delta, \upsilon>0,
\end{equation}
where $f_\ast := \min_{x} f(x)$ and  $$\|y\|_{1,\upsilon}:=\sum_{i=1}^m\,\sqrt{|y_i|^2+\upsilon^2}.$$
We refer to Remark \ref{uurem-a-} for a brief discussion concerning the choice of structure for $V$. Note that
$$|y|\leq \sqrt{y^2+\upsilon^2}\to |y|\quad\mbox{for all $y\in\mathbb R$ and as $\upsilon\to 0$}, $$
and $\|y\|_{1,\upsilon}$ should be seen as a regularization of the $\mathrm{L}^1$-norm. In the following $\upsilon\in (0,1)$ will be fixed throughout the section, while $\theta,\alpha,\beta,\delta$ are degrees of freedom. In particular, the estimates derived will be uniform in $\upsilon\in (0,1)$ and hence will apply, by limiting arguments, when $\upsilon\to 0$.



\begin{lemma}\label{bog} Assume that $f$ satisfies \textnormal{(A1)}, i.e., $\eqref{D-}$. Let $V=V_\upsilon(x,z,y)$ be defined as in \eqref{V}. Then  there exist constants $c_1,c_2>0$ (depending on $\theta,\alpha,\beta,\delta$ and the data of the problem ($f$,  $L_f$) such that
\begin{equation}\label{Vlla}
V(x,z,y)\leq c_1\bigl(\|x\|^2+\|z\|^2+\|y\|_{1,\upsilon}\bigr) + c_2, \qquad \forall (x,z,y)\in\R^{m}\times \R^m\times \R^m,
\end{equation}
\end{lemma}
\begin{proof} Since $f$ satisfies \textnormal{(A1)}, the descent lemma implies that for all $x,\bar x\in\R^m$,
\[
f(x)\le f(\bar x)+\nabla f(\bar x)\cdot(x-\bar x)
+\frac{L_f}{2}\|x-\bar x\|^2 .
\]
Let $x_\ast$ be a minimizer of $f$. Then $\nabla f(x_\ast)=0$ and $f(x_\ast)=f_\ast$. Choosing $\bar x=x_\ast$ yields
\[
f(x)-f_\ast
\le \frac{L_f}{2}\|x-x_\ast\|^2\le L_f\|x\|^2 + L_f\|x_\ast\|^2\le C_f^1 \|x\|^2 + C_f^2,
\]
for some constants $C_f^1,C_f^2>0$.  Combining this with  Young's inequality we obtain, for every $\epsilon>0$
\[
\begin{aligned}
V(x,z,y)
&\le \theta\bigl(C_f^1\|x\|^2 + C_f^2\bigr)
+\frac{\alpha}{2}\|z\|^2
+\frac{\beta\epsilon}{2}\|x\|^2
+\frac{\beta}{2\epsilon}\|z\|^2
+\delta\|y\|_{1,\upsilon} .
\end{aligned}
\]
Grouping the quadratic terms gives
\[
V(x,z,y)
\le
\Bigl(\theta C_f^1 + \frac{\beta\epsilon}{2}\Bigr)\|x\|^2
+
\Bigl(\frac{\alpha}{2}+\frac{\beta}{2\epsilon}\Bigr)\|z\|^2
+
\delta\|y\|_{1,\upsilon}
+
\theta C_f^2 .
\]
Choosing $\epsilon>0$ arbitrarily (for instance $\epsilon=1$) and defining
\[
c_1
:=
\max\Bigl\{
\theta C_f^1 + \frac{\beta\epsilon}{2},
\frac{\alpha}{2}+\frac{\beta}{2\epsilon},
\delta
\Bigr\},
\qquad
c_2:=\theta C_f^2,
\]
we obtain
\[
V(x,z,y)
\le c_1\bigl(\|x\|^2+\|z\|^2+\|y\|_{1,\upsilon}\bigr)
+ c_2 ,
\]
for all $(x,z,y)\in\R^m\times\R^m\times\R^m$. The constants $c_1,c_2>0$ depend only on
$\theta,\alpha,\beta,\delta$ and the data of the problem
($f$, $L_f$), and are independent of $\upsilon\in(0,1)$.
\end{proof}


\begin{lemma}\label{lowb}  Assume that $f$ satisfies \textnormal{(A2)}, i.e., $\eqref{D}$. Let $V=V_\upsilon(x,z,y)$ be defined as in \eqref{V} and assume that
\begin{equation}\label{Vcond}({\theta m_f}/{4}-{\beta^2}/{\alpha})>0.
\end{equation} Then  there exist constants $\hat c_1,\hat c_2>0$ (depending on $\theta,\alpha,\beta,\delta$ and the data of the problem ($f$,  $m_f$, $c$, and $R$)) such that
\begin{equation}\label{eq:drift-ineql}
V(x,z,y)\;\ge\; \hat c_1\bigl(\|x\|^2+\|z\|^2+\|y\|_{1,\upsilon}\bigr) - \hat c_2,
\qquad \forall (x,z,y)\in\R^{m}\times \R^m\times\R^m.
\end{equation}
\end{lemma}
\begin{proof}  Note that the constant $f_\ast$ ensures that the first term in $V$ is nonnegative. However, $V$ is not
necessarily nonnegative globally because of the cross term $-\beta\,x\cdot z$. Under the asymptotic dissipativity condition \eqref{D},  the function $f$ dominates a quadratic at infinity. Indeed, fix any $m_f'\in(0,m_f)$, e.g., $m_f'=m_f/2$. For $\|x\|\ge R$, consider the function
\[
g(r) := \min_{\|x\|=r} f(x), \qquad r\ge R.
\]
Differentiating along rays $x=ru$ with $\|u\|=1$, we obtain
\[
\frac{\mathrm{d}}{\mathrm{d}r} f(ru) \;=\; \langle \nabla f(ru),\,u\rangle
\;\ge\; \frac{1}{r}\langle \nabla f(ru),\,ru\rangle
\;\ge\; m_f r - \frac{c}{r}, \qquad r\ge R.
\]
Integrating this inequality from $R$ to $r$ yields
\[
f(ru) - f(Ru) \;\ge\; \frac{m_f}{2}(r^2-R^2) - c\log\frac{r}{R},
\qquad r\ge R, \;\|u\|=1.
\]
Thus, for all $\|x\|\ge R$,
\[
f(x) \;\ge\; \frac{m_f'}{2}\|x\|^2 - C,
\]
for some $C>0$ depending on $f$, $m_f'$, $m_f$, $c$, and $R$.
In particular, $f$ is coercive and grows at least quadratically outside a compact set. Consequently,  we also have (repeating the argument with $f$ replaced by $f-f_\ast$), for all $\|x\|\ge R$,
\[
\theta\bigl(f(x)-f_\ast\bigr) \;\ge\; \frac{\theta m_f'}{2}\|x\|^2 - \theta C.
\]
Therefore, the term $\theta\bigl(f(x)-f_\ast\bigr)$ controls $\|x\|^2$ outside a compact set and the term can be used to absorb the cross term $-\beta\,x\cdot z$. Indeed, by Young's inequality, for any $\chi>0$,
\[
-\beta\,x\cdot z \;\ge\; -\frac{\beta^2}{2\chi}\|x\|^2 - \frac{\chi}{2}\|z\|^2.
\]
Combining these bounds, and noting that the $\frac\alpha 2\|z\|^2$ and $\delta\|y\|_{1,\upsilon}$ terms already contribute positively, we see that  if  $(x,z,y)\in \R^{m}\times \R^m\times \R^m$ with $\|x\|\ge R$, then
\begin{align*}
V(x,z,y)&\geq \frac{\theta m_f'}{2}\|x\|^2 - \theta C+\frac \alpha 2\|z\|^2+\delta \|y\|_{1,\upsilon}-\frac{\beta^2}{2\chi}\|x\|^2 - \frac{\chi}{2}\|z\|^2\notag\\
&=\bigl (\frac{\theta m_f'}{2}-\frac{\beta^2}{2\chi}\bigr )\|x\|^2 +\bigl (\frac \alpha 2-\frac{\chi}{2}\bigr )\|z\|^2+\delta \|y\|_{1,\upsilon} - \theta C.
\end{align*}
Letting $\chi=\alpha/2$, $m_f'=m_f/2$, and imposing the restriction in \eqref{Vcond} we can conclude that  there exist constants $\hat c_1,\hat c_2>0$ (depending on $\theta,\beta,\delta,\chi$ and the data of the problem ($f$, $m_f'$, $m_f$, $c$, and $R$)) such that
\begin{equation*}
V(x,z,y)\;\ge\; \hat c_1\bigl(\|x\|^2+\|z\|^2+\|y\|_{1,\upsilon}\bigr) - \hat c_2,
\qquad \forall (x,z,y)\in\R^{m}\times \R^m\times \R^m.
\end{equation*}
Hence $V$ is a coercive Lyapunov function, even though it need not be nonnegative everywhere.
\end{proof}







We next establish a Foster-Lyapunov drift inequality for $\mathcal{L}_\epsilon$ and  $V$ by proving the following lemma.

\begin{lemma}\label{drift} Assume that $f$ satisfies \textnormal{(A1)} and \textnormal{(A2)}. Let $\upsilon\in (0,1)$ be fixed, let $\epsilon\in [0,\upsilon]$, and consider
the operator $\mathcal{L}_\epsilon$ introduced \eqref{Regloss1}. Let $V=V_\upsilon$ be given by \eqref{V} with $\alpha=1$.  Let $R\gg 1$ be large enough and consider
the set $$\hat C_R:=\{(x,z,y)\in \R^{m}\times \R^m\times \R^m : V(x,z,y)\le R\}.$$
There exist parameters $\beta,\theta,\delta>0$ and constants $\lambda>0$, $K<\infty$,
(depending only on $a,b,\gamma,\varepsilon,\sigma,L_f,$ the data in \textnormal{(A2)}, and the chosen $\beta,\theta,\delta$, $R$,) such that
\begin{equation}\label{eq:drift-ineq}
{\mathcal L}_\epsilon V(x,z,y)
\;\le\; -\,\lambda\,V(x,z,y) + K\,\mathbf 1_{\hat C_R},
\qquad \forall (x,z,y)\in \R^{m}\times \R^m\times \R^m.
\end{equation}
\end{lemma}
\begin{proof} For $V$ we have
\[
\nabla_x V = \theta \nabla f(x)-\beta z,\quad,
\Delta_x V = \theta \Delta f(x),
\]
\[
\nabla_z V = z-\beta x,\quad
\Delta_z V = m.
\]
and
\[
\nabla_y V = \delta \bigl(y_1/\sqrt{y_1^2+\upsilon^2},...,y_m/\sqrt{y_m^2+\upsilon^2}\bigr),\quad
\Delta_y V = \delta \sum_{i=1}^m\frac {\upsilon^2}{(y_i^2+\upsilon^2)^{3/2}}.
\]
Hence, using the convention that the division by $(\sqrt{|y|}+\varepsilon)$ is done componentwise,
\[
\begin{aligned}
{\mathcal L}_\epsilon V
&= (\theta\nabla f-\beta z)\cdot\Bigl(-\gamma\,\frac{z}{\sqrt |y|+\varepsilon}\Bigr)
+ (z-\beta x)\cdot a(\nabla f-z) \\
&\quad +  {\delta b}\Bigl (\sum_{i=1}^m(y_i/\sqrt{y_i^2+\upsilon^2})(-y_i + (\partial_{x_i}f)^2 + \sigma^2)\bigr )\\
&\quad+ \frac 12 a^2\sigma^2\,m+\epsilon \theta \Delta f(x)+ {\epsilon\delta} \sum_{i=1}^m\frac {\upsilon^2}{(y_i^2+\upsilon^2)^{3/2}}.
\end{aligned}
\]
Note that
\[
\begin{aligned}
\sum_{i=1}^m\Bigl(\frac{y_i}{\sqrt{y_i^2 + \upsilon^2}}\Bigr)
\left(-y_i + (\partial_{x_i} f)^2 + \sigma^2\right)
&= -\|y\|_{1,\upsilon}+\sum_{i=1}^m\frac{\upsilon^2}{\sqrt{y_i^2 + \upsilon^2}}+ \sum_{i=1}^m\frac{y_i\big((\partial_{x_i} f)^2 + \sigma^2\big)}{\sqrt{y_i^2 + \upsilon^2}},
\end{aligned}
\]
and
\begin{equation}\label{eq:LV-expanded}
\begin{aligned}
{\mathcal L}_\epsilon V
&= T_1+T_2+T_3- \delta b\|y\|_{1,\upsilon}
+  \frac12 a^2\sigma^2\, m+R,
\end{aligned}
\end{equation}
where
\begin{align*}
T_1&:= \Bigl[a\,z\cdot\nabla f - \theta\gamma \sum_{i=1}^m \frac{z_i\,\partial_{x_i}f(x)}{\sqrt{|y_i|}+\varepsilon}\Bigr],\\
T_2&:= \Bigl[-a\|z\|^2 + \beta\gamma \sum_{i=1}^m \frac{z_i^2}{\sqrt{|y_i|}+\varepsilon}\Bigr], \\
T_3&:= \Bigl[-a\beta\, x\cdot\nabla f + a\beta\, x\cdot z\Bigr],
\end{align*}
and where
\begin{align*}
R&:=\epsilon \theta \Delta f(x)+{\epsilon\delta} \sum_{i=1}^m\frac {\upsilon^2}{(y_i^2+\upsilon^2)^{3/2}}+
\sum_{i=1}^m{\delta b}\frac{\upsilon^2}{\sqrt{y_i^2 + \upsilon^2}}+ \sum_{i=1}^m {\delta b}\frac{y_i\big((\partial_{x_i} f)^2 + \sigma^2\big)}{\sqrt{y_i^2 + \upsilon^2}}.
\end{align*}
Note that
$$R\leq C+ {\delta b}\|\nabla f(x)\|^2,$$
for a constant $C$ which depends only on $\theta, \delta, b$, and $\sigma$.  We next estimate $T_1+T_2$, and then $T_3$. Write $w_i:= (\sqrt{|y_i|}+\varepsilon)^{-1}\in(0,\varepsilon^{-1}]$. Then
$$T_1+T_2=\sum_i(T_1+T_2)_i$$
where
\[
(T_1+T_2)_i= -\underbrace{(a-\beta\gamma w_i)}_{=:c_{1,i}}\,z_i^2
\;+\;\underbrace{(a-\theta\gamma w_i)}_{=:c_{2,i}}\,z_i\,\partial_{x_i}f(x).
\]
Assume
\begin{equation}\label{eq:beta-small}
0<\beta<\frac{a\varepsilon}{2\gamma}\quad\Longrightarrow\quad c_{1,i}\ge a-\frac{\beta\gamma}{\varepsilon}=:c_{1,\min}>\frac{a}{2}\,.
\end{equation}
Using the inequality $-c_1 z^2+c_2 zh\le -\tfrac{c_1}{2}z^2+\tfrac{c_2^2}{2c_1}h^2$, we get
\[
(T_1+T_2)_i \;\le\; -\frac{c_{1,i}}{2}z_i^2 \;+\; \frac{c_{2,i}^2}{2c_{1,i}}\,(\partial_{x_i}f(x))^2
\;\le\; -\frac{c_{1,\min}}{2}z_i^2 \;+\; \frac{\bigl(a-\theta\gamma w_i\bigr)^2}{2c_{1,\min}}\,(\partial_{x_i}f(x))^2.
\]
Summing over $i$ and using $w_i\le\varepsilon^{-1}$, we obtain
\begin{equation}\label{eq:T12}
T_1+T_2 \;\le\; -\frac{a-\beta\gamma/\varepsilon}{2}\,\|z\|^2
\;+\; \frac{\bigl(a-\theta\gamma/\varepsilon\bigr)^2}{2\,(a-\beta\gamma/\varepsilon)}\,\|\nabla f(x)\|^2.
\end{equation}
We next estimate $T_3$. By (A2) and Young’s inequality, for any $\eta_z>0$,
\[
-a\beta\,x\cdot\nabla f(x)\;\le\;-a\beta\,(m_f\|x\|^2-c),
\qquad
a\beta\,x\cdot z\;\le\;\frac{\eta_z}{2}\|z\|^2+\frac{a^2\beta^2}{2\eta_z}\|x\|^2.
\]
Fix
\begin{equation}\label{eq:eta-choice}
\eta_z:=\frac{a-\beta\gamma/\varepsilon}{2}\,,
\end{equation}
which is positive by \eqref{eq:beta-small}. Then
\begin{equation}\label{eq:T3}
T_3 \;\le\; -\Bigl(a\beta m_f-\frac{a^2\beta^2}{a-\beta\gamma/\varepsilon}\Bigr)\frac{\|x\|^2}{2}
\;+\; \frac{a-\beta\gamma/\varepsilon}{4}\,\|z\|^2 \;+\; a\beta c.
\end{equation}
Combining \eqref{eq:T12} and \eqref{eq:T3},
\begin{align}\label{eq:T123}
T_1+T_2+T_3
\;\le&\; -\underbrace{\frac{a-\beta\gamma/\varepsilon}{4}}_{=:c_z}\|z\|^2
\;-\;\underbrace{\frac{1}{2}\Bigl(a\beta m_f-\frac{a^2\beta^2}{a-\beta\gamma/\varepsilon}\Bigr)}_{=:c_x}\|x\|^2\notag\\
\;&+\;\underbrace{\frac{(a-\theta\gamma/\varepsilon)^2}{2\,(a-\beta\gamma/\varepsilon)}}_{=:c_g}\|\nabla f(x)\|^2
\;+\;a\beta c.
\end{align}
Note that $c_z>0$ by \eqref{eq:beta-small}, and $c_x>0$ for all sufficiently small $\beta>0$ (since the second term is $O(\beta^2)$). From \eqref{eq:LV-expanded}, \eqref{eq:T123} and the estimate of $E$,
\[
\mathcal{L}_\epsilon V \;\le\; -c_z\|z\|^2 - c_x\|x\|^2 + (c_g+\delta b)\|\nabla f(x)\|^2
\;-\;\delta b\,\|y\|_{1,\upsilon} \;+\; C_0,
\]
where $C_0:=C + \tfrac12 a^2\sigma^2\, m + a\beta c$.
Using (A1), $\|\nabla f(x)\|\le L_f\|x\|+\|\nabla f(0)\|$, hence
\[
(c_g+\delta b)\|\nabla f(x)\|^2 \;\le\; 2L_f^2(c_g+\delta b)\|x\|^2 + C_1,
\]
for some $C_1<\infty$. Thus
\begin{equation}\label{eq:preLyap}
\mathcal{L}_\epsilon V \;\le\; -c_z\|z\|^2 \;-\;\bigl(c_x-2L_f^2(c_g+\delta b)\bigr)\|x\|^2
\;-\;\delta b\,\|y\|_{1,\upsilon} \;+\; C_2,
\end{equation}
for another constant $C_2$. We now choose parameters to make the $x$-coefficient negative.
Pick $\beta>0$ small so that \eqref{eq:beta-small} holds and $c_x\ge \tfrac{a\beta m_f}{4}$. With $\beta$ fixed, we pick $\theta\in(0,a\varepsilon/\gamma)$ so that
\[
c_g=\frac{(a-\theta\gamma/\varepsilon)^2}{2\,(a-\beta\gamma/\varepsilon)}
\;\le\;\frac{a\beta m_f}{16L_f^2}.
\]
Finally, pick $\delta>0$ so that $\delta b\le \tfrac{a\beta m_f}{16L_f^2}$.
Then $c_x-2L_f^2(c_g+\delta b)\ge \tfrac{a\beta m_f}{8}=:c_x'>0$.
With these choices, \eqref{eq:preLyap} becomes
\begin{equation}\label{eq:good-drift}
\mathcal{L}_\epsilon V \;\le\; -\,c_z\|z\|^2 \;-\;c_x'\|x\|^2 \;-\;\delta b\,\|y\|_{1,\upsilon} \;+\; C_2.
\end{equation}
Let $\bar c:=\min\{c_z,c_x',\delta b\}>0$.
Then from \eqref{eq:good-drift},
\begin{align}\label{Adda-}
\mathcal{L}_\epsilon V \;\le\; -\bar c \bigl(\|x\|^2+\|z\|^2+\|y\|_{1,\upsilon}\bigr) + C_2.
\end{align}
Using Lemma \ref{bog} we have
\begin{equation}\label{eq:drift-ineqlapa}
-V(x,z,y)\;\geq\;  -c_1\bigl(\|x\|^2+\|z\|^2+\|y\|_{1,\upsilon}\bigr) -c_2,
\qquad \forall (x,z,y)\in\R^{m}\times \R^m\times\R^m.
\end{equation}
Hence,
\begin{align}\label{Adda-}
\mathcal{L}_\epsilon V \;\le\; \frac {\bar c}{c_1} (-V(x,z,y)+c_2) + C_2=-\hat C_1V(x,z,y)+\hat C_2.
\end{align}
Using this we first see that on the compact set $\hat C_R=\{(x,z,y): V(x,z,y)\le R\}$, the function $\mathcal L V$ is bounded above by continuity, i.e., $\mathcal L_\epsilon V(x,z,y)\ \le K$. In particular, if we choose $R\gg 1$ large enough then we see from Lemma \ref{bog} and Lemma \ref{lowb} that
\begin{align}\label{Adda}
-\bar c \bigl(\|x\|^2+\|z\|^2+\|y\|_{1,\upsilon}\bigr) + C_2&\sim -\bar c \bigl(\|x\|^2+\|z\|^2+\|y\|_{1,\upsilon}\bigr)\sim -\hat  cV(x,z,y),
\end{align}
for all $(x,z,y)\in (\R^{m}\times \R^m\times \R^m)\setminus \hat C_R$. Consequently,
we obtain
\[
\mathcal L_\epsilon V(x,z,y)\ \le\ -\lambda\,V(x,z,y)\ +\ K\,\mathbf 1_{\hat C_R}(x,z,y),
\]
which is \eqref{eq:drift-ineq}.
\end{proof}




\begin{remark}\label{uurem-a-}
In the construction of $V$ in \eqref{V}, the dependence on $y$ enters through the quantity $\|y\|_{1,\upsilon}$. While the remaining components of $V$ are relatively natural (or ``guessable''), one may wonder whether alternative choices for the dependence on $y$ could also be reasonable. This is indeed the case in Lemma \ref{bog} and Lemma \ref{lowb}. However, in the proof of Lemma \ref{drift} (already in the case $\epsilon=0$), terms of the form
\begin{equation}\label{terma}
\sum_{i=1}^m G_i(y_i)\big((\partial_{x_i} f)^2 + \sigma^2\big)
\end{equation}
naturally arise, where each $G_i$ reflects the specific way in which $V$ depends on $y$. Unless one makes more substantial modifications to the structure of $V$, it is desirable that the expression in \eqref{terma} can be controlled by a quantity of the form $\|\nabla f(x)\|^2 + 1$. This requirement essentially forces the dependence on $y$ to be compatible with such a bound, which explains our use of the $\mathrm{L}^1$-type norm in $y$.
\end{remark}

\begin{remark}\label{uurem} Note that by Lemma \ref{lowb} we have  $V:\R^{m}\times \R^m\times \R^m\to [-c_2,\infty)$. However, replacing
$V$ by $V:=V+A$ where $A$ is a non-negative constant, does not the change the validity of \eqref{Adda-} nor of \eqref{Adda}. Consequently,  we can construct a Lyapunov function that takes values in $[1,\infty)$.
\end{remark}

\begin{remark}\label{uuremla} By construction $V=V_\upsilon$, $\upsilon\in (0,1)$, but  the
quantitative results proved above are uniform in $\upsilon\in (0,1)$. In particular, all estimates holds for $\mathcal{L}_0$ uniform in $\upsilon\in (0,1)$.
\end{remark}



By \eqref{eq:drift-ineql}, i.e., coercivity in $(x,z,y)$,  the sublevel sets $\{V_\upsilon\le M\}$ are compact in
$\R^{m}\times\R^m\times\R^m$. Choosing $M$ large enough,  \eqref{eq:drift-ineq} can be stated, for all $\epsilon\in [0,\upsilon]$, as
\begin{equation}\label{eq:HL-drift}
{\mathcal L}_\epsilon V_\upsilon\;\le\;-\lambda' V_\upsilon + K'\,\mathbf 1_{C_\upsilon},\qquad C_\upsilon:=\{(x,z,y)\in \R^{m}\times \R^m\times\R^m : V_\upsilon(x,z,y)\le M\}\ \text{compact},
\end{equation}
for some $\lambda'\in(0,\lambda]$, $K'\ge K$. This is the standard reduction as outside $\{V_\upsilon\le M\}$ the $-\lambda V_\upsilon$ term dominates any bounded perturbation. As \eqref{eq:HL-drift} applies with $\epsilon=0$ we see that
\begin{equation}\label{eq:HL-drift+-}
{\mathcal L} V_\upsilon(x,z,y)={\mathcal L}_0 V_\upsilon(x,z,y)\;\le\;-\lambda' V_\upsilon(x,z,y) + K'\,\mathbf 1_{C_\upsilon}(x,z,y)
\end{equation}
holds for all $(x,z,y)\in \R^{m}\times \R^m\times\R^m$. Finally, note that $C_\upsilon\subset C:=C_0$ for all $\upsilon\in [0,1]$. Hence,
\begin{equation}\label{eq:HL-drift+}
{\mathcal L} V_\upsilon(x,z,y)={\mathcal L}_0 V_\upsilon(x,z,y)\;\le\;-\lambda' V_\upsilon(x,z,y) + K'\,\mathbf 1_{C}(x,z,y)
\end{equation}
holds for all $(x,z,y)\in \R^{m}\times \R^m\times\R^m$.











\subsection{Existence of invariant measures using vanishing viscosity}\label{exista} As previously mentioned,  the existence of invariant probability measures for dissipative stochastic
dynamical systems can often be established by relatively soft arguments. We here show that though our system is degenerate, the existence of invariant measures can be establish using the framework with vanishing viscosity outlined above.

Assume that $f$ satisfies \textnormal{(A1)} and \textnormal{(A2)} from Subsection~\ref{Subcond}. Fix $a,b,\gamma,\sigma>0$ and $\varepsilon>0$, and let $(P_t)_{t\ge0}$ denote the Markov semigroup
associated with the time-homogeneous system \eqref{eq:cts-x+}–\eqref{eq:cts-y+}. Let $\mu_0$ be an admissible initial probability distribution on $\R^{3m}$ in the sense of~\eqref{intimeasure}. Then, writing $\pi_t := \mu_0 P_t$, we intend to prove that the family $(\pi_t)_{t\ge0}$ converges weakly, as $t\to\infty$,
to an invariant probability measure $\pi_\infty$ on $\R^{3m}$ associated with $(P_t)$,
\[
\pi_\infty P_t = \pi_\infty, \qquad t \ge 0,
\]
that is, the law of $(x_t,z_t,y_t)$ remains $\pi_\infty$ whenever the process is initialized with
$(x_0,z_0,y_0)\sim\pi_\infty$.


 To start the argument, consider the operator $\mathcal{L}_\epsilon$ introduced in \eqref{Regloss1} on $\R^m\times\R^m\times\R^m$, $\epsilon>0$. By construction
\[
\mathcal{L}_\epsilon \;=\; \mathcal{L} \;+\; \epsilon\,\Delta_x\;+\; \epsilon\,\Delta_y,\qquad \epsilon>0,
\]
on $\R^m\times\R^m\times(\R_+)^m$. $\mathcal{L}_\epsilon$ is an extension of $\mathcal{L}$ to all of $\R^m\times\R^m\times\R^m$  and it is also a regularization of $\mathcal{L}$ as $\mathcal{L}_\epsilon$ is uniformly elliptic  on $\R^m\times\R^m\times\R^m$.









Consider $$U_R:=\{(x,z,y)\in \R^m\times\R^m\times\R^m:\, \|x\|^2+\|z\|^2+\|y\|^2\leq R^2\}$$ for $R\gg 1$ large. If $(x,z,y)\in \mathbb R^{3m}\setminus U_R$, then we have  $\|x\|^2+\|z\|^2+\|y\|_{1,\upsilon}\geq R/2$. Hence, using Lemma \ref{lowb} and  \eqref{eq:HL-drift} we see, for $R$ large enough, that
$$\mathcal{L}_\epsilon V\leq -C/2\quad\mbox{whenever}\quad (x,z,y)\in \mathbb R^{3m}\setminus U_R,$$
for some constant $C\gg1 $ independent of $\epsilon$. Furthermore,
$$|(\mathcal{L}_\epsilon -\mathcal{L}_0)V|\leq c_V\epsilon \to 0\quad\mbox{as}\quad \epsilon\to 0.$$
Using Theorem 2.4.1 and Corollary 2.4.2 in \cite{BogachevKrylovRoeckner2015} we can conclude there are probability
measures $\{\mu_\epsilon\}$, with positive continuous densities $\{\rho_\epsilon\}$, solving the equations
$$\mathcal{L}_\epsilon^\ast\mu_\epsilon=0\quad\mbox{on $\mathbb R^{3m}$.}$$
 We renormalize these solutions so that $\mu_\epsilon(U_R)=1$. By Theorem 2.3.2 in \cite{BogachevKrylovRoeckner2015}, the renormalized sequence $\{\mu_\epsilon\}$  is still bounded, since the functions
$\{\mathcal{L}_\epsilon V\}$ are
uniformly bounded on the ball $U_R$. Hence $\{\mu_\epsilon\}$ contains a subsequence that converges
weakly on every ball in $\R^{3m}$ to a  measure $\nu$. The measure $\nu$ obtained in the limit is bounded and
nonnegative, but is not identically zero as $\nu(U_R) = 1$.

Let $\varphi\in C_0^\infty\big(\R^m\times\R^m\times\R^m\big)$. Then $\mathcal L_\epsilon\varphi\to \mathcal L_0\varphi$ uniformly as $\epsilon\to 0$,
and $\sup_\epsilon \|\mathcal L_\epsilon\varphi\|_\infty<\infty$. Since $\iiint \mathcal L_\epsilon\varphi\,\mathrm d\mu_\epsilon=0$ for each $\epsilon$,
passing to the limit gives
\begin{equation}\label{watt}
\iiint \mathcal L_0\varphi\,\mathrm d\nu \;=\; \lim_{\epsilon\to0}\iiint \mathcal L_\epsilon\varphi\,\mathrm d\mu_\epsilon \;=\; 0.
\end{equation}
Thus $\mathcal L_0^\ast\nu=0$ in the sense of distributions on
$\R^m\times\R^m\times\R^m$. Renormalizing $\nu$ we get a probability measure on $\R^{3m}$.





To complete the proof we show that the support of $\nu$ is contained in
$\R^m\times\R^m\times(\sigma^2/2,\infty)^m$, and hence that $\nu$ is an invariant measure for
$\mathcal L$ as well. We have already established, see \eqref{watt}, that
\[
\iiint \mathcal L_0\varphi\,\mathrm d\nu = 0
\qquad\text{for all }\varphi\in C_0^\infty(\R^{3m}),
\]
where $\mathcal L_0$ is the full-space operator in~\eqref{Regloss1} with $\epsilon=0$.
For $\varphi$ depending only on $y$, we have
\[
\mathcal L_0\varphi(y)
=\sum_{i=1}^m b\bigl(-y_i+(\partial_{x_i}f(x))^2+\sigma^2\bigr)\,\partial_{y_i}\varphi(y),
\]
since the $x$- and $z$-derivatives of $\varphi$ vanish and there is no diffusion in $y$. Fix $i\in\{1,\dots,m\}$ and choose $\eta\in C^\infty(\R)$ such that
\[
\eta'(r)\ge 0,\qquad
\eta'(r)=1\ \text{for }r\le\sigma^2/4,\qquad
\eta'(r)=0\ \text{for }r\ge\sigma^2/2.
\]
For each $R>0$, let $\chi_R\in C_0^\infty(\R^{3m})$ be a standard smooth cutoff satisfying
$\chi_R\equiv1$ on $\{(x,z,y):\|(x,z,y)\|\le R\}$ and $\chi_R\equiv0$ on
$\{(x,z,y):\|(x,z,y)\|\ge 2R\}$.  Set
\[
\varphi_R(x,z,y):=\eta(y_i)\chi_R(x,z,y).
\]
Then $\varphi_R\in C_0^\infty(\R^{3m})$, and we may use it as a test function in the stationary equation.
Expanding the identity $\iiint \mathcal L_0\varphi_R\,\mathrm d\nu=0$ we have
\begin{equation}\label{watt+}
\iiint b\bigl(-y_i+(\partial_{x_i}f(x))^2+\sigma^2\bigr)\,\eta'(y_i)\chi_R(x,z,y)\,
\mathrm d\nu(x,z,y)
+ \iiint \mathcal R_R(x,z,y)\,\mathrm d\nu(x,z,y)=0,
\end{equation}
where $\mathcal R_R$ collects the additional terms involving derivatives of $\chi_R$.
Because $\chi_R\equiv1$ on $\|(x,z,y)\|\le R$ and $\eta'$ is bounded,
the remainder $\mathcal R_R$ is supported in $\{R\le \|(x,z,y)\|\le 2R\}$ and satisfies
$|\mathcal R_R|\le C/R$. Hence
\[
\lim_{R\to\infty}\iiint \mathcal R_R\,\mathrm d\nu=0,
\]
and we may pass to the limit $R\to\infty$ in \eqref{watt+} to obtain
\[
0=\iiint b\bigl(-y_i+(\partial_{x_i}f(x))^2+\sigma^2\bigr)\,\eta'(y_i)\,\mathrm d\nu(x,z,y).
\]
Because $(\partial_{x_i}f(x))^2+\sigma^2\ge\sigma^2$, on the set $\{y_i\le\sigma^2/2\}$ we have
\[
-y_i+(\partial_{x_i}f(x))^2+\sigma^2 \;\ge\; \sigma^2/2 \;>\; 0.
\]
Therefore the integrand
$b\bigl(-y_i+(\partial_{x_i}f(x))^2+\sigma^2\bigr)\eta'(y_i)$
is nonnegative and strictly positive wherever $y_i\le\sigma^2/4$.
Since its integral is zero, we must have
\[
\nu\bigl(\{y_i\le\sigma^2/4\}\bigr)=0.
\]
As $i$ was arbitrary and we can approximate the indicator of $\{y_i<\sigma^2/2\}$
from below by such $\eta'$, we conclude that
\[
\nu\bigl(\{y:\,y_i<\sigma^2/2\}\bigr)=0
\quad\text{for every }i,
\]
hence
\[
\mathrm{supp}(\nu)
\subset
\R^m\times\R^m\times\bigl([\sigma^2/2,\infty)\bigr)^m.
\]

On the region $\{y_i>0\}$ (in fact on $\{y_i\geq\sigma^2/2\}$) the two operators coincide
\[
\mathcal L_0=\mathcal L
\qquad\text{on}\quad
\R^m\times\R^m\times\bigl([\sigma^2/2,\infty)\bigr)^m,
\]
because $\sqrt{|y_i|}=\sqrt{y_i}$ there. Let
$\varphi\in C_0^\infty\bigl(\R^m\times\R^m\times(\R_+)^m\bigr)$
be any test function. By the support property just proved, we may extend $\varphi$ by zero
to a function in $C_0^\infty(\R^{3m})$ without changing the integral against $\nu$.
Hence
\[
\iiint \mathcal L\varphi\,\mathrm d\nu
= \iiint \mathcal L_0\varphi\,\mathrm d\nu
= 0.
\]
Thus $\mathcal L^\ast\nu=0$ in the sense of distributions on
$\R^m\times\R^m\times(\R_+)^m$. Finally, since the SDE coefficients are continuous with at most linear growth,
the associated Markov semigroup $(P_t)_{t\ge0}$ is Feller.
For Feller semigroups the generator $\mathcal L$ (on a core containing $C_0^\infty$)
characterizes invariance: $\mathcal L^\ast\nu=0$ implies
\[
\nu P_t = \nu,\qquad \forall t\ge0.
\]
We have proved that $\nu$ is an invariant probability measure for the original
diffusion with generator~$\mathcal L$.\qed



\begin{remark}
In the argument above we establish the \emph{existence} of at least one invariant probability measure
$\pi_\infty$ as a weak limit of $(\pi_t)_{t\ge0}$. The proof combines vanishing viscosity
regularization, Lyapunov functions ($V_\upsilon$), and results from~\cite{BogachevKrylovRoeckner2015}.
Note that we only assume \textnormal{(A1)}-\textnormal{(A2)} and no Hörmander/Malliavin
nondegeneracy (see \cite{BouleauHirsch1991,Nualart2006}). The limit $\pi_\infty$ may in principle depend on $\mu_0$, and no rate of convergence
is inferred, in particular, uniqueness is not claimed.
\end{remark}


\section{Uniqueness of invariant measures}\label{aaaa}

Recall that the total variation distance $\|\cdot\|_{\mathrm{TV}}$
and the $2$-Wasserstein distance $W_2$ are defined in
Section~\ref{Prelim}.  With $V=V_0\geq 0$ being the Lyapunov function with $\upsilon=0$ and constructed  in Section~\ref{sec:lya}, we introduce
\begin{align}\label{pettabol}\hat C_R:=\Bigl\{(x,z,y)\in \mathbb R^m\times\mathbb R^m\times\mathbb R^m:
V(x,z,y)\le R\Bigr\}
\end{align}
for every $R>0$.  Recall the sets $\{E_\eta\}$ introduced in \eqref{intimeasurell} and let $t_\ast:=(\ln 2)/b$.

The proof of  Theorem \ref{CorC} is based the following abstract result proved in this section.
\begin{theorem}\label{Thm3-revised-}
Assume that $f$ satisfies \textnormal{(A1)} and \textnormal{(A2)} from Subsection~\ref{Subcond}. Fix $a,b,\gamma,\sigma>0$ and $\varepsilon>0$,
and let $(P_t)_{t\ge0}$ denote the Markov semigroup associated with
\eqref{eq:cts-x+}-\eqref{eq:cts-y+}. Consider the level sets $\{\hat C_R\}$ introduced in \eqref{pettabol} for $R>0$. Assume that for every $R>0$ the compact set
\begin{align}\label{petta}
C_R :=\hat C_R\cap E_{\sigma^2/2}\quad\mbox{is small,}
\end{align}
in the sense that there exist $t_R= t_\ast+\delta_R$, $\delta_R> 0$,  $\varepsilon_R > 0$, and a probability measure $\nu_R$
on $(E,\mathcal B(E))$ such that for all $(x,z,y) \in C_R$ and all
$A \in \mathcal B(E)$,
\begin{align}\label{petta-rep}
P_{t_R}((x,z,y),A) \ge \varepsilon_R\,\nu_R(A).
\end{align}
Then there exists a unique invariant probability measure $\pi_\infty$ on $E$, and for every admissible initial distribution $\mu_0$ we have
\[
\mu_0 P_t \Rightarrow \pi_\infty
\qquad \text{as } t\to\infty .
\]
In particular, the limiting distribution does not depend on $\mu_0$. Moreover, there exist constants $C<\infty$ and $\lambda>0$, independent of the admissible initial distribution $\mu_0$, such that
\begin{equation}\label{eq:V-geometric-ergodicity-revised}
\|\mu_0P_t-\pi_\infty\|_{\mathrm{TV}}
\le
C e^{-\lambda (t-t_\ast)}
\bigl(1+\mathbb{E}_{\mu_0}\!\big[V(x_{t_\ast},z_{t_\ast},y_{t_\ast})\big]\bigr),
\qquad t\ge t_\ast,
\end{equation}
and
\begin{equation}\label{eq:W2-erg-revised}
\bigl (W_2\bigl(\mu_0P_t,\pi_\infty\bigr)\bigr )^2
\le
C e^{-\lambda (t-t_\ast)}
\bigl(1+\mathbb{E}_{\mu_0}\!\big[V(x_{t_\ast},z_{t_\ast},y_{t_\ast})\big]\bigr),
\qquad t\ge t_\ast.
\end{equation}
\end{theorem}
\begin{remark} Note that $t_\ast=(\ln 2)/b$ is tied to the set $E_{\sigma^2/2}$ and $(x_t,z_t,y_t)\in E_{\sigma^2/2}$ for all $t>t_\ast$. However, there is nothing really special with the choice of $t_\ast$ in the sense that for any $\hat t>0$ there is a $\hat\eta>0$ such that $(x_t,z_t,y_t)\in E_{\hat \eta}$ for all $t>\hat t$. In particular, the theorem remains true with $C_R$ replaced with $C_R:=\hat C_R\cap E_{\hat\eta}$ and with $t_\ast$ replaced by $\hat t$. Still the constants in the quantitative estimates will depend on $\hat t$ and $\hat\eta$.
\end{remark}


We present the proof of  Theorem \ref{Thm3-revised-} in two subsections.  The idea is to apply a continuous-time version of the  Harris-Meyn-Tweedie (HMT) theorem tailored to our diffusion. In particular,  we will use the following version of the classical Harris theorem \cite{Harris1956,meyn1993foster, MeynTweedie2009,HairerMattingly2011b} stated as Theorem 4.2 in \cite{HairerMattinglyScheutzow2011}.

\begin{theorem}\label{thm:HMT-cts}
Let \( (\tilde P_t)_{t \ge 0} \) be a Markov semigroup over a Polish space \( \mathcal{X} \) such that
there exists a Lyapunov function \( V \) with the additional property that the level sets
\(C_R:=\{ X \in \mathcal{X} : V(X) \le R \}\) satisfies the following. Given \( R > 0 \), there exist a time
\( t_0> 0 \) and a constant \( \varepsilon > 0 \) such that
\begin{equation}\label{keycond}
\| \tilde  P_{t_0}(X, \cdot) - \tilde  P_{t_0}(\hat X, \cdot) \|_{\mathrm{TV}} \le 1 - \varepsilon,\qquad\mbox{for every \( X, \hat X \in C_R \).}
\end{equation}
 Then \( \tilde  P_t \) has a unique invariant measure \( \mu_* \) and
\[
\| \tilde  P_t(X, \cdot) - \mu_* \|_{\mathrm{TV}}
\le C e^{-\lambda t} (1 + V(X))\qquad\forall X\in\mathcal{X},\, \forall\, t\geq 0,
\]
for some positive constants \( C \) and \( \lambda \).
\end{theorem}


\begin{remark}
Recall that the total variation distance between two probability
measures is equal to \( 1 \) if and only if the two measures are mutually singular. \eqref{keycond} therefore states that the transition probabilities starting from any two points in the set $C_R$
have a ``common part'' of mass at least \( \varepsilon \).
\end{remark}






\subsection{Proof of existence/uniqueness and \eqref{eq:V-geometric-ergodicity-revised}} We intend to apply Theorem \ref{thm:HMT-cts} to the process $X_t=(x_t,z_t,y_t)$ defined by \eqref{eq:cts-x+}-\eqref{eq:cts-y+}.  $X_t$ is Markov process on
$E=\R^m\times\R^m\times(\R_+)^m$  with transition semigroup $(P_t)_{t\ge0}$. Note that $E_{\sigma^2/2}$ is a Polish space.

We have established a  Foster–Lyapunov drift condition in the sense that for every $R>0$, by \eqref{eq:HL-drift} there exist $\lambda'>0$, $K'<\infty$ and a measurable $V:E_{\sigma^2/2}\to[1,\infty)$ (we may assume $V\ge1$ by Remark~\ref{uurem}) such that
\begin{equation}\label{eq:HMT-drift-used}
\mathcal{L}_0V \;\le\; -\lambda' V \;+\; K'\,\mathbf 1_{C_R},
\end{equation}
 $\mathcal L_0$ is the \emph{extended} generator of $X$,  coinciding with $\mathcal L$ on $E$. Hence \eqref{eq:HMT-drift-used} is in the standard HMT form.

 Let $\tilde P_{t}(\xi,\cdot):=P_{t+t_\ast}(\xi,\cdot)$ for $t\geq 0$. Using the minorization assumption stated in  \eqref{petta} we have that there exists $t_0=\delta_R>0$ such that
\begin{equation}\label{alla+}
\mu(\cdot) \ge \varepsilon \nu(\cdot), \qquad \mu'(\cdot) \ge \varepsilon \nu(\cdot)
\end{equation}
for a measure $\nu$, whenever  \( \mu(\cdot) = \tilde P_{t_0}(\xi,\cdot) \) and \( \mu'(\cdot) = \tilde P_{t_0}(\hat \xi,\cdot) \) for some $\xi,\hat\xi\in C_R=\hat C_R\cap E_{\sigma^2/2}$. We can without loss of generality assume that $\nu$ is a probability measure.  Define the residual probabilities
\[
\tilde\mu := \frac{\mu - \varepsilon \nu}{1 - \varepsilon},
\qquad
\tilde\mu' := \frac{\mu' - \varepsilon \nu}{1 - \varepsilon}.
\]
Then \( \tilde\mu, \tilde\mu' \) are probability measures and
\[
\mu = \varepsilon \nu + (1 - \varepsilon)\tilde\mu,
\qquad
\mu' = \varepsilon \nu + (1 - \varepsilon)\tilde\mu'.
\]
We now construct a coupling \( (Y,Y') \) in the following way. First, with probability \( \varepsilon \), draw \( Y = Y' \sim \nu \). Second,  with probability \( 1 - \varepsilon \), draw \( Y \sim \tilde\mu \) and \( Y' \sim \tilde\mu' \) independently.
This is a valid coupling of \( \mu, \mu' \) and
\begin{align*}
\mathbb P(Y\neq Y')
&= \varepsilon\,\underbrace{\mathbb P(Y\neq Y'\mid \text{common draw from }\nu)}_{=\,0}\\
&\quad \;+\;
(1-\varepsilon)\,\mathbb P\bigl(Y\neq Y'\mid Y\sim\tilde\mu,\ Y'\sim\tilde\mu'\ \text{i.i.d.}\bigr)
\;\le\; 1-\varepsilon,
\end{align*}
since the conditional probability is at most $1$. Equivalently,
\[
\mathbb P(Y=Y') \ge \varepsilon
\quad\Rightarrow\quad
\mathbb P(Y\neq Y') \le 1-\varepsilon.\] Furthermore,  for any measurable set $A$,
\[
\mu(A)-\mu'(A)
= \mathbb P(Y\in A)-\mathbb P(Y'\in A)
= \mathbb P(Y\in A, Y'\notin A)-\mathbb P(Y\notin A, Y'\in A)
\le \mathbb P(Y\neq Y').
\]
Swapping the roles of $\mu$ and $\mu'$ gives $\mu'(A)-\mu(A)\le \mathbb P(Y\neq Y')$,
so $|\mu(A)-\mu'(A)|\le \mathbb P(Y\neq Y')$ for all $A$. Taking the supremum over
$A$ yields the coupling inequality
\[
\|\mu-\mu'\|_{\mathrm{TV}}
= \sup_{A} |\mu(A)-\mu'(A)|
\le \mathbb P(Y\neq Y').
\]
In particular, in our case we have
\[
\|\mu - \mu'\|_{\mathrm{TV}} \le \mathbb{P}(Y \neq Y') \le 1 - \varepsilon.
\]
Hence, given $R$ there  exists a time
\( t_0 > 0 \) and a constant \( \varepsilon > 0 \) such that
\begin{equation}\label{keycond+}
\| \tilde P_{t_0}(\xi, \cdot) - \tilde P_{t_0}(\hat\xi, \cdot) \|_{\mathrm{TV}} \le 1 - \varepsilon,\qquad\mbox{for every \( \xi, \hat \xi \in C_R \).}
\end{equation}
Hence, using Theorem \ref{thm:HMT-cts} we can conclude that $X_t$ is positive Harris recurrent, that its admits a unique invariant probability measure $\mu_\ast$, and that there exist $C,\lambda>0$ such that
\begin{equation}\label{eq:TV-ergprev}
\|\tilde P_t((x,z,y),\cdot)-\mu_\ast\|_{\mathrm{TV}}
\ \le\ C\,e^{-\lambda t}\,\bigl(1+V(x,z,y)\bigr),
\qquad \forall (x,z,y)\in E_{\sigma^2/2},\ \forall t\ge 0.
\end{equation}
Consequently, renaming $\mu_\ast$ to $\pi_\infty$ we see that
\begin{equation}\label{eq:TV-erg}
\| P_{t}((x,z,y),\cdot)-\pi_\infty\|_{\mathrm{TV}}
\ \le\ C\,e^{-\lambda (t-t_\ast)}\,\bigl(1+V(x_{t_\ast},z_{t_\ast},y_{t_\ast})\bigr),
\qquad \forall (x,z,y)\in E,\ \forall t\ge t_\ast.
\end{equation}
From \eqref{eq:TV-erg} we obtain a corresponding bound for general initial
distributions. Indeed, let \(\mu_0\) be a probability measure on \(E\). Then
\[
\mu_0P_t - \pi_\infty
=
\iiint_E \bigl(P_t((x,z,y),\cdot)-\pi_\infty\bigr)\,\mu_0(\d(x,z,y)).
\]
Since the total variation norm is a norm on the space of finite signed measures, the triangle inequality yields
\[
\|\mu_0P_t-\pi_\infty\|_{\mathrm{TV}}
\le
\iiint_E \|P_t((x,z,y),\cdot)-\pi_\infty\|_{\mathrm{TV}}\,
\mu_0(\d(x,z,y)).
\]
Applying \eqref{eq:TV-erg} yields
\[
\|\mu_0P_t-\pi_\infty\|_{\mathrm{TV}}
\le
C e^{-\lambda (t-t_\ast)}
\iiint_E \bigl(1+V(x_{t_\ast},z_{t_\ast},y_{t_\ast})\bigr)\,\mu_0(\d(x,z,y)).
\]
In particular,
\[
\|\mu_0P_t-\pi_\infty\|_{\mathrm{TV}}
\le
C e^{-\lambda (t-t_\ast)}(1+\mathbb E_{\mu_0}[V(x_{t_\ast},z_{t_\ast},y_{t_\ast})]).
\]
Hence \(\mu_0P_t \to \pi_\infty\) exponentially fast in total variation. \qed



\begin{remark}
The Foster–Lyapunov function from Lemma~\ref{lowb} and the drift inequality of Lemma~\ref{drift} are crucial to the argument. Note that to apply Theorem \ref{thm:HMT-cts} no global PDE regularity (e.g.\ hypoellipticity or smooth transition densities) is required. The only analytic inputs are (a) the extended–generator inequality for $V$, see \eqref{eq:HMT-drift-used}, and (b) a (local) minorization, see \eqref{petta}.
\end{remark}






\subsection{Proof of \eqref{eq:W2-erg-revised}} We have proved, see \eqref{eq:V-geometric-ergodicity-revised}, that
\begin{equation}\label{HM}
\|\pi_t-\pi_\infty\|_{\mathrm{TV}}
\ \le\ C\,e^{-\lambda (t-t_\ast)}\,\bigl(1+\mathbb E_{\mu_0}[V(x_{t_\ast},z_{t_\ast},y_{t_\ast})]\bigr),\quad t\geq t_\ast.
\end{equation}
We now intend to pass from convergence in total variation to convergence in $W_2$.  In order to do so we use the following truncation bound.


\begin{lemma}\label{kkaa} Let $\mu,\nu$ be probability measures on $\R^{d}$ with finite second moments.
Then, for any $R>0$,
\[
\bigl (W_2(\mu,\nu)\bigr )^2
\;\le\; 8R^2\,\|\mu-\nu\|_{\mathrm{TV}}
\;+\;4\!\!\int_{\{\|u\|>R\}}\!\!\|u\|^2\,\mu(\mathrm du)
\;+\;4\!\!\int_{\{\|v\|>R\}}\!\!\|v\|^2\,\nu(\mathrm dv).
\]
\end{lemma}

\begin{proof}
Let $(U,V)$ be a maximal coupling of $(\mu,\nu)$ so that
$\mathbb P(U\neq V)=\delta:=\|\mu-\nu\|_{\mathrm{TV}}$.
On the event $\{U\neq V,\ \|U\|\le R,\ \|V\|\le R\}$ we have $\|U-V\|\le 2R$, hence
\[
\mathbb E[\|U-V\|^2:\,\|U\|\le R,\|V\|\le R]\le 4R^2\,\delta.
\]
On the complementary event, we use $\|U-V\|^2\le 2\|U\|^2+2\|V\|^2$ and integrate
\[
\mathbb E[\|U-V\|^2:\,\{\|U\|>R\}\cup\{\|V\|>R\}]
\le 2\!\int_{\{\|u\|>R\}}\|u\|^2\,\mu(\mathrm du)
+2\!\int_{\{\|v\|>R\}}\|v\|^2\,\nu(\mathrm dv).
\]
Combining gives the stated inequality.
\end{proof}

 In Appendix \ref{HMI} below we prove, see \eqref{eq:2nd-moment-unif}, that\begin{equation*}
\sup_{t\ge0}\,\mathbb E_{\mu_0}\bigl[\| x_t\|^2+\| z_t\|^2+\| y_t\|\bigr]\ <\ \infty.
\end{equation*}
However, in the proof of \eqref{eq:W2-erg-revised} we need higher integrability of $ X_t=( x_t, z_t,  y_t)$ and  Appendix \ref{HMI}  is devoted to the proof of such estimates. The results of Appendix \ref{HMI} imply
\[
\sup_{t\ge0}\,\mathbb E\bigl[\| X_t\|^{2+\delta}\bigr] < \infty
\quad\text{for some }\delta>0.
\]
This higher integrability implies, by the de la Vallée–Poussin criterion, that the family $\{\pi_t:t\ge0\}$ is uniformly integrable in $\mathrm L^2$. Define
\[
M_{2+\delta}\;:=\;\sup_{t\ge0}\iiint \|u\|^{2+\delta}\,\pi_t(\mathrm du)\;<\;\infty,
\]
where we now start to use the compact notation $u=(x,z,y)$ to, in the end, convey with the notation used in Lemma \ref{kkaa}. Furthermore, using the same Foster-Lyapunov drift in Lemma~\ref{lem:power-Lyap} we deduce integrability under the invariant distribution and that
\[
M_{2+\delta}^\infty\;:=\;\iiint \|v\|^{2+\delta}\,\pi_\infty(\mathrm dv)\;<\;\infty.
\]
Now fix  $R>0$. Using the elementary bound $\mathbf 1_{\{\|u\|>R\}}\|u\|^2\le R^{-\delta}\|u\|^{2+\delta}$, we obtain
\[
\sup_{t\ge0}\ \iiint_{\{\|u\|>R\}}\|u\|^2\,\mathrm d\pi_t(u)
\;\le\; R^{-\delta}\,\sup_{t\ge0}\iiint \|u\|^{2+\delta}\,\mathrm d\pi_t(u)
\;=\; M_{2+\delta}\,R^{-\delta},
\]
and similarly
\[
\iiint_{\{\|v\|>R\}}\|v\|^2\,\mathrm d\pi_\infty(v)
\;\le\; M_{2+\delta}^\infty\,R^{-\delta}.
\]
Hence one may take the explicit tail modulus
\[
\omega(R)\;=\;M_\star\,R^{-\delta},
\qquad
M_\star:=\max\{M_{2+\delta},\,M_{2+\delta}^\infty\},
\]
which satisfies $\omega(R)\downarrow0$ as $R\to\infty$ and yields
\begin{equation}\label{eq:UI}
\sup_{t\ge0}\ \iiint_{\{\|u\|>R\}}\|u\|^2\,\mathrm d\pi_t(u) \;\le\; \omega(R),
\qquad
\iiint_{\{\|v\|>R\}}\|v\|^2\,\mathrm d\pi_\infty(v) \;\le\; \omega(R).
\end{equation}



We now apply the truncation lemma, i.e., Lemma \ref{kkaa}, with $\mu=\pi_t$, $t\geq t_\ast$,  $\nu=\pi_\infty$. Hence, for any $R>0$,
\begin{align}\label{eq:W2-TV}
\bigl (W_2(\pi_t,\pi_\infty)\bigr )^2
\;&\le\; 8R^2\,\|\pi_t-\pi_\infty\|_{\mathrm{TV}}
\;+\;4\!\!\iiint_{\{\|u\|>R\}}\!\!\|u\|^2\,\mathrm d \pi_t(u)
\;+\;4\!\!\iiint_{\{\|v\|>R\}}\!\!\|v\|^2\, \mathrm d \pi_\infty(v)\notag\\
&\le \; 8R^2\,\|\pi_t-\pi_\infty\|_{\mathrm{TV}}+8\omega(R).
\end{align}
Next, for every $\varepsilon>0$ we can choose $R=R(\varepsilon)$ so large that
\[
\bigl (W_2(\pi_t,\pi_\infty)\bigr )^2
\;\le\; 8(R(\varepsilon))^2\,\|\pi_t-\pi_\infty\|_{\mathrm{TV}}
\;+\;4\varepsilon.
\]
Now fix this $R(\varepsilon)$.  Using \eqref{HM} we see that
\begin{equation}\label{eq:W2-TV-final}
\bigl (W_2(\pi_t,\pi_\infty)\bigr )^2
\ \le\ 8\,(R(\varepsilon))^2\,{C}\,e^{-\lambda (t-t_\ast)}\,\bigl(1+\mathbb E_{\mu_0}[V(x_{t_\ast},z_{t_\ast},y_{t_\ast})]\bigr)
\;+\;4{\varepsilon}.
\end{equation}
Letting $t\to\infty$ first (with $R(\varepsilon)$ fixed) gives
\[
W_2(\pi_t,\pi_\infty)\;\longrightarrow\; 2\sqrt{\varepsilon},
\]
and since $\varepsilon>0$ is arbitrary we conclude
\[
\lim_{t\to\infty} W_2(\pi_t,\pi_\infty)=0.
\]
Being more carefully we see that
$$\varepsilon/2=\omega(R(\varepsilon))=M_\star\,(R(\varepsilon))^{-\delta}\implies R(\varepsilon)=C_\star\varepsilon^{-\delta}.$$
Hence we can express \eqref{eq:W2-TV-final} as
\begin{equation}\label{eq:W2-TV-final+}
W_2(\pi_t,\pi_\infty)
\ \le\ 2\sqrt{2}\, C_\star\varepsilon^{-\delta}\,\sqrt{C}\,e^{-\lambda (t-t_\ast)/2}\,\bigl(1+\mathbb E_{\mu_0}[V(x_{t_\ast},z_{t_\ast},y_{t_\ast})]\bigr)^{1/2}
\;+\;2\sqrt{\varepsilon}.
\end{equation}
Moreover, the explicit estimate \eqref{eq:W2-TV-final+} shows an \emph{exponential} rate of convergence (with exponent $\lambda/2$), up to an arbitrarily small additive term $2\sqrt{\varepsilon}$, and it holds for all $\varepsilon>0$. Given $t$, $\varepsilon>0$ is a degree of freedom which we can choose to decay with $t$ to remove the last additive term in the final statement. To outline this,  estimate \eqref{eq:W2-TV-final+} gives for every fixed $\varepsilon>0$
\[
W_2(\pi_t,\pi_\infty)
\ \le\  \hat C_\star\varepsilon^{-\delta}e^{-\lambda (t-t_\ast)/2}\,\bigl(1+\mathbb E_{\mu_0}[V(x_{t_\ast},z_{t_\ast},y_{t_\ast})]\bigr)^{1/2}\;+\;2\sqrt{\varepsilon},
\qquad t\ge t_\ast,
\]
where $ \hat C_\star$ is independent of $\varepsilon$. We now let $\varepsilon=\varepsilon(t-t_\ast)$ to depend on $(t-t_\ast)$, and chosen so that
\[
2\sqrt{\varepsilon(t-t_\ast)} \;\le\;  \hat C_\star(\varepsilon(t-t_\ast))^{-\delta}\,e^{-\lambda (t-t_\ast)/2} \Longleftrightarrow  (\varepsilon(t-t_\ast))^{1/2+\delta}\leq
\frac 1 2 C_\star\,e^{-\lambda (t-t_\ast)/2}\qquad t\ge t_\ast.
\]
For example, if we choose $\varepsilon(t-t_\ast):=e^{-\lambda (t-t_\ast)/(1/2+\delta)}$,  then we can conclude the stated inequality  for $t-t_\ast$ sufficiently large. With this choice, the bound simplifies to
\[
W_2(\pi_t,\pi_\infty)
\;\le\; C'e^{-\lambda (t-t_\ast)/2}\,\bigl(1+\mathbb E_{\mu_0}[V(x_{t_\ast},z_{t_\ast},y_{t_\ast})]\bigr)^{1/2},
\qquad t\ge t_\ast,
\]
for some constant $C'>0$ independent of the initial condition.
This removes the additive term and yields a clean exponential convergence rate in $W_2$. \qed









\section{Mixing and uniqueness: Proof of  Theorem \ref{CorC}}\label{sec:proof-CorC}



Having constructed the Lyapunov function, Theorem~\ref{Thm3-revised-} reduces the proof of Theorem \ref{CorC} to verifying that
the compact sets $\{C_R\}=\{\hat C_R\cap E_{\sigma^2/2}\}$, see \eqref{petta}, are small. In our case the latter is not immediate as the system
\eqref{eq:cts-x+}-\eqref{eq:cts-y+}
is \emph{degenerate} in the sense that the diffusion acts only in the $z$-variables, and as global hypoellipticity can not be ensured.

To start the proof, we fix $R>0$ and we intend to prove that $C_R$ is small. Since $\mathcal D_A^\dagger\neq\R^m$, see \eqref{Remma1+}, and as  $\mathcal D_A^\dagger$ is closed, there exists $x_\ast\in\R^m\setminus\mathcal D_A^\dagger$ and  $r>0$ such that
\begin{equation}\label{good-ball}
\overline{B(x_\ast,2r)}\subset\R^m\setminus\mathcal D_A^\dagger.
\end{equation}
Since $C_R$ is compact,
there exists $M_R>0$ such that
\begin{equation}\label{Bg1}
|x_0|+|z_0|+|y_0|\le M_R
\qquad
\text{for all }(x_0,z_0,y_0)\in C_R.
\end{equation}
We first prove the following lemma.
\begin{lemma}\label{LemJ1}  There exist positive constants $(\rho,c_R,d_R)$, a time $T_R>0$, and a constant $\eta_R>0$, such that if we introduce
\begin{equation*}
 W':=O\times Z\times Y:=B(x_\ast,r)\times B(0,\rho)\times (c_R/2,2d_R)^m,
\end{equation*}
then
\begin{equation}\label{eqbo2}
\inf_{u\in C_R}
P_{T_R}(u,{W'})\ge\eta_R.
\end{equation}
\end{lemma}
\begin{proof}
Starting from $(x_0,z_0,y_0)\in C_R$ we first construct a control, in two steps, based on the skeleton system \eqref{eq:skeleton-components}.

%\medskip
\noindent
{\bf Step 1.} First, we steer $z_0$ to $0$. Fix $\tau>0$ and define
\[
z(t)=(1-{t}/{\tau})z_0,
\qquad
t\in[0,\tau].
\]
Then $z(0)=z_0$ and  $z(\tau)=0$. Given $z(t)$ we define $(x(t),y(t))$ on $[0,\tau]$ as the solution of
\[
\dot x^i(t)=
-\gamma\frac{z^i(t)}{\sqrt{y^i(t)}+\varepsilon},
\qquad
\dot y^i(t)=
b\bigl(-y^i(t)+(\partial_{x_i}f(x(t)))^2+\sigma^2\bigr),\qquad (x(0),y(0))=(x_0,y_0).
\]
The control producing this trajectory is
\[
h_i(t)=\frac1{a\sigma}
\bigl(\dot z^i(t)-a(\partial_{x_i}f(x(t))-z^i(t))\bigr).
\]
By the definition of the set $C_R$ we know that $y^i(t)\geq \sigma^2/2$ for all $t\geq 0$.  Still, for our purpose we here simply use the conservative estimate
\[
|x(t)-x_0|
\le
\frac{\gamma}{\varepsilon}
\int_0^{\tau} |z(s)|\,\d s
\le \frac{\gamma}{\varepsilon}M_R\tau .
\]
Thus, starting at $(x_0,z_0,y_0)\in C_R$ after time $\tau$ we obtain, simply using condition (A1) and elementary ODE estimates, a state $(x_1,0,y_1)$ with
\begin{align}\label{exa}
|x_1|\le R_1,
\qquad
y_1\in[a_{R},b_{R}]^m,\quad a_{R}\geq \sigma^2/2,
\end{align}
where $R_1$ and $a_{R},b_{R}$ only depends on $(x_0,z_0,y_0)$ through $M_R$ (see \eqref{Bg1}) and $R$.

%\medskip
\noindent
{\bf Step 2.} Second, we steer $x_1$ to $x_\ast$. Fix $S>0$ and choose a smooth interpolation
\[
x(t)=x_1+\eta(t-\tau)(x_\ast-x_1),
\qquad
t\in[\tau,\tau+S],
\]
with $\eta(0)=0$, $\eta(S)=1$ and $\dot\eta(0)=\dot\eta(S)=0$. Then,
$$x(\tau)=x_1,\quad x(\tau+S)=x_\ast,\quad \dot x(\tau)=0,\quad \dot x(\tau+S)=0.$$
Define $y(t)$ on $[\tau,\tau+S]$
as the solution  to
\[
\dot y^i(t)=
b\bigl(-y^i(t)+(\partial_{x_i}f(x(t)))^2+\sigma^2\bigr),
\qquad
y(\tau)=y_1 .
\]
Given $R$, and for $x_\ast$, $\tau$ and $S$ fixed, we can conclude that there exists constants $c_R<d_R$, we here only emphasize the dependence  on $R$, such that $c_R\geq \sigma^2/2$ and such that
\begin{align}\label{terra}
c_R\le y^i(t)\le d_R\qquad
t\in[0,T],\quad T:=\tau+S.
\end{align}
$T_R=T$ is the notation appearing in the statement of the lemma, but we here for simplicity stick to $T$.  Define $z(t)$ on $[\tau,\tau+S]$ according  to
\[
z^i(t)=
-\frac{\sqrt{y^i(t)}+\varepsilon}{\gamma}\dot x^i(t), \qquad
z^i(\tau)=0,
\]
and choose the control
\[
h_i(t)=\frac1{a\sigma}
\Bigl(\dot z^i(t)-a(\partial_{x_i}f(x(t))-z^i(t))\Bigr).
\]
We obtain, starting at $(x_1,0,y_1)$ satisfying \eqref{exa}, we obtain after time $S$ a state $(x_\ast,0,y)$
where $y$ satisfies \eqref{terra}.


%\medskip
\noindent
By construction and concatenation we have constructed a triple path $(x(t),z(t),y(t))$ which solves the skeleton system
on $[0,T]$, and satisfies $(x(0),z(0),y(0))=(x_0,z_0,y_0)$ and
\begin{equation}\label{eqbo}
x(T)=x_\ast,
\qquad
z(T)=0, \qquad y(T)\in (c_R,d_R)^m.
\end{equation}
Note that the path $t\to (x(t),z(t),y(t))$ as well as the path of the control $t\to h(t)$ are continuous.

%\medskip
To proceed we next construct an open target set. Given $x_\ast\in \R^m\setminus \mathcal D_A^\dagger$ and $r>0$ as above, we let $ O:=B(x_\ast,r)$. From the explicit skeleton construction above and concluded in \eqref{eqbo}, we have that for every $u=(x_0,z_0,y_0)\in C_R$ there exists a (continuous) control $h_u\in \mathrm L^2([0,T];\R^m)$ such that the corresponding
controlled trajectory satisfies
\[
\Phi_T(u,h_u)=\bigl(x^{u,h_u}(T),z^{u,h_u}(T),y^{u,h_u}(T)\bigr)
\in O\times Z\times Y,
\]
where $Z:=B(0,\rho)$ and $Y:=(c_R/2,2d_R)^m$. In particular, both $Z$ and $Y$ are open and connected sets chosen in the construction. Let
\[
W':=O\times Z\times Y.
\]
Then $W'\subset E$ is a nonempty, bounded, open, and connected set.


Fix \(u\in C_R\). Since \(\Phi_T(u,h_u)\in W'\) and \(W'\) is open, continuity of the
skeleton endpoint map
\[
v\mapsto \Phi_T(v,h_u)
\]
implies that there exists an open neighborhood \(N_u\) of \(u\) in \(E\) such that
\[
\Phi_T(v,h_u)\in W'
\qquad\text{for all }v\in N_u.
\]
In particular,
\[
\mathcal R_T(v)\cap W'\neq\varnothing
\qquad\text{for all }v\in N_u.
\]
Therefore, by the Stroock-Varadhan support theorem (see Theorem~\ref{thm:support}),
\begin{equation}\label{eqbo1}
P_T(v,W')>0
\qquad\text{for all }v\in N_u.
\end{equation}
Note that in \eqref{eqbo1} both the finite time \(T>0\) and the open set \(W'\)
are independent of \(u\).

Since \(W'\) is open and the Markov semigroup is Feller, the map
\[
v\mapsto P_T(v,W')
\]
is lower semicontinuous. Hence, from \eqref{eqbo1} it follows that there exist
an open neighborhood (still denoted) \(N_u\) of \(u\) and a constant \(\eta_u>0\)
such that
\[
P_T(v,W')\ge \eta_u
\qquad \text{for all } v\in N_u.
\]

The family \(\{N_u\}_{u\in C_R}\) forms an open cover of the compact set \(C_R\).
We may therefore extract a finite subcover \(N_{u_1},\dots,N_{u_k}\).
Defining
\[
\eta:=\min_{1\le j\le k}\eta_{u_j}>0,
\]
we conclude that
\[
\inf_{u\in C_R} P_T(u,W')\ge \eta .
\]
This proves the lemma.
\end{proof}


By construction, the Hörmander bracket condition holds on the open and connected set
\[
W''':=B(x_\ast,2r)\times \mathbb R^m\times (\mathbb R_+)^m .
\]
Moreover, with $W'$ as constructed in Lemma \ref{LemJ1},  \(\overline{W'}\) is compactly contained in \(W'''\).
Hence the transition kernel \(P_t(v,\cdot)\) admits a \(C^\infty\)-density $p(t,v,w)$
on \((0,\infty)\times W'''\times W'''\). To complete the argument and the proof of Theorem \ref{CorC}, it remains to show that there exist \(T>0\) and
a nonempty open, connected set \(W''\subset W'\), with \(\overline{W''}\)
compactly contained in \(W'\), such that every point in \(W''\) is reachable
by the skeleton system starting from any point in \(W'\).
This is achieved in the following lemma.

\begin{lemma}\label{lem:access-xzy-openV-3stagefinal}
There exist \(T>0\), \(y_\ast\in\mathbb R^m\), and positive constants
\((\rho_x,\rho_z,\rho_y)\) such that if we introduce
\[
W'' :=
B(x_\ast,\rho_x)\times B(0,\rho_z)\times B(y_\ast,\rho_y)
\subset W'
\]
then every point in \(W''\) is an attainable
endpoint at time \(T\) of the skeleton system \eqref{eq:skeleton-components}
starting from any initial condition \(U_0\in W'\). Moreover, the corresponding controlled trajectory satisfies
\[
(x(t),z(t),y(t))\in W'''
\qquad \text{for all } t\in[0,T].
\]
\end{lemma}

Assuming Lemma~\ref{lem:access-xzy-openV-3stagefinal}, the Stroock-Varadhan
support theorem implies that
\[
P_T(v,U)>0
\]
for every open set \(U\subset W''\) and every \(v\in W'\).
Since the density \(p(t,v,w)\) is smooth on \(W'''\times W'''\), it also follows
that the map
\[
(v,w)\mapsto p(T,v,w)
\]
is jointly continuous and strictly positive on \(W'\times W''\).

We now choose now an open set \(W\subset W''\) such that
\(\overline{W}\) is compactly contained in \(W''\).
Since \(p(T,\cdot,\cdot)\) is continuous and strictly positive on the compact
set \(\overline{W'}\times \overline{W}\), we obtain
\[
\kappa :=
\inf_{v\in\overline{W'},\, w\in\overline{W}} p(T,v,w)
>0 .
\]
Using this bound we deduce, for any Borel set \(A\subset E\),
\[
P_T(v,A)
\ge
\kappa\,\operatorname{Leb}(A\cap W)
=
(\kappa\,\operatorname{Leb}(W))
\frac{\operatorname{Leb}(A\cap W)}{\operatorname{Leb}(W)}
=: \kappa'
\frac{\operatorname{Leb}(A\cap W)}{\operatorname{Leb}(W)}
\qquad
\text{for all } v\in W'.
\]
Let \(t_\ast:=T_R+T\), where \(T_R\) is as in Lemma~\ref{LemJ1} and \(T\) is
as in Lemma~\ref{lem:access-xzy-openV-3stagefinal}. Combining the above with
the Chapman-Kolmogorov equation, we deduce for any Borel set \(A\subset E\)
that
\begin{align}\label{eq:local-minor-++a}
P_{t_\ast}(u,A)
&\ge
\iiint_{W'} P_{T_R}(u,\mathrm{d}v)\,P_T(v,A)\ge
\eta_R\,\kappa'
\frac{\operatorname{Leb}(A\cap W)}{\operatorname{Leb}(W)}
\qquad
\text{for all } u\in C_R .
\end{align}
In particular, defining the probability measure
\[
\nu(A):=
\frac{\operatorname{Leb}(A\cap W)}{\operatorname{Leb}(W)},
\]
we obtain
\begin{equation}\label{eq:minor-far-style}
P_{t_\ast}(u,\cdot)
\ \ge\
\eta_R\kappa'\,\nu(\cdot),
\qquad \forall u\in C_R .
\end{equation}
Thus \(C_R\) is a small set, and the proof of Theorem~\ref{CorC} is complete
modulo the proof of Lemma~\ref{lem:access-xzy-openV-3stagefinal}.\qed





\subsection{Proof of  Lemma \ref{lem:access-xzy-openV-3stagefinal}} We let $\rho_x>0$ be such that
\begin{equation}\label{eq:Ux-final}
U_x':=B(x_\ast,4\rho_x)
\subset B(x_\ast,r)\subset\mathbb R^m\setminus \mathcal D_A^\dagger\implies
 \det A(x)\neq0 \ \text{on } U_x'.
\end{equation}
Define
\[
g(x)=\big((\partial_{x_1}f(x))^2,\dots,(\partial_{x_m}f(x))^2\big).
\]
A direct computation yields $Dg(x)=2A(x)$, hence
$\det Dg(x)\neq0$ on $U_x'$. After shrinking $\rho_x$ if necessary,
\begin{equation}\label{eq:g-diffeo-final}
g:U_x'\to g(U_x')
\quad\text{is a $C^1$ diffeomorphism onto an open set.}
\end{equation}
From here on $\rho_x$ is fixed. Let
\[
y_\ast := g(x_\ast)+\sigma^2\mathbf 1,
\]
and note that as $g(B(x_\ast,2\rho_x))$ is open,
there exists $\rho_y>0$ such that
\[
U_y
:=
B(y_\ast,\rho_y)
\subset
g\big(B(x_\ast,2\rho_x)\big)
+
\sigma^2\mathbf 1.
\]
We are going to construct $W''$ as
\[
W'' := U_x \times U_z \times U_y \subset \R^{3m},
\]
i.e., as a product of open sets, where
\[
U_x:=B(x_\ast,\rho_x),\qquad
U_z=B(0,\rho_z),
\qquad
U_y:=B(y_\ast,\rho_y),
\]
and where  $\rho_z$ will be specified later. Note that we can without loss of generality assume $B(0,2\rho_z)\subset Z$ and $B(y_\ast,2\rho_y)\subset Y$, where the sets $Z$ and $Y$ are as introduced in Lemma \ref{LemJ1}.




To prove Lemma \ref{lem:access-xzy-openV-3stagefinal} we are going to construct controls on three time intervals $[0,T_1]$, $[T_1,T_1+S]$, and $[T_1+S,T_1+S+\delta]$, where  $T_1>0$ is small, $S>0$ is large, and $\delta>0$ is small, all to be fixed and specified. We set
\begin{equation}\label{frame1}
T:=T_1+S+\delta.
\end{equation}
We construct the controls to prove that there exists, for each target triple
\begin{equation}\label{frame2}
(x^\sharp,z^\sharp,y^\sharp)\in
U_x\times U_z\times U_y,
\end{equation}
 a controlled skeleton system $(x(t),z(t),y(t))$ starting at  $U_0=(x_0,z_0,y_0)\in W'$,  and
ending up at $(x^\sharp,z^\sharp,y^\sharp)$ at $t=T$, i.e., $(x(T),z(T),y(T))=(x^\sharp,z^\sharp,y^\sharp)$. Moreover, we are going to do the construction so that the controlled trajectory satisfies
\[
(x(t),z(t),y(t))\in W'''
\qquad \text{for all } t\in[0,T].
\]
We construct the control in steps dividing the argument into a number of lemmas, Lemma \ref{lem:initial-hermite-path}-Lemma \ref{CdG+} below. These lemmas are used to complete the proof of Lemma \ref{lem:access-xzy-openV-3stagefinal}, and their proofs are postponed to subsequent subsections.


We first prove the following lemma, see Subsection \ref{Pr0} for its proof.
\begin{lemma}\label{lem:initial-hermite-path}
Fix \(T_1>0\), \(U_0=(x_0,z_0,y_0)\in W'\), and \(x^c\in B(x_\ast,2\rho_x)\).
Define
\[
v_0^i:=-\gamma\frac{z_0^i}{\sqrt{y_0^i}+\varepsilon},
\qquad i=1,\dots,m.
\]
Then there exists a \(C^\infty\)-curve $x:[0,T_1]\to \mathbb R^m$ such that
\[
x(0)=x_0,\qquad x(T_1)=x^c,\qquad \dot x(0)=v_0,\qquad \dot x(T_1)=0.
\]
Moreover, \(T_1>0\) can be chosen sufficiently small so that, for every
\(U_0\in W'\) and every \(x^c\in B(x_\ast,2\rho_x)\),
\[
x([0,T_1])\subset B(x_\ast,2r).
\]
\end{lemma}




Using Lemma~\ref{lem:initial-hermite-path}, we in the following choose \(T_1>0\) sufficiently small so that, for every
\(U_0\in W'\) and every prescribed terminal point \(x^c\in B(x_\ast,2\rho_x)\), the constructed curve $x([0,T_1])$, $x(T_1)=x^c$, satisfies
\begin{equation}\label{contain}
x([0,T_1])\subset B(x_\ast,2r).
\end{equation}
Using the curve \(x(\cdot)\) we let \(y(\cdot)\) be the solution of
\[
\dot y^i(t)=b\bigl(-y^i(t)+(\partial_{x_i}f(x(t)))^2+\sigma^2\bigr),
\qquad y(0)=y_0.
\]
Since \(U_0\in W'=O\times Z\times Y\), \(Y=(c_R/2,2d_R)^m\),  and due to \eqref{contain}, the regularity of $f$  yields the uniform bounds
\[
0<c_1\le y^i(t)\le d_1<\infty,
\qquad t\in[0,T_1],\ i=1,\dots,m,
\]
for suitable constants \(c_1,d_1\) uniformly with respect to \(U_0\in W'\). Given $(x(t),y(t))$, define
\[
z^i(t):=-\frac{\sqrt{y^i(t)}+\varepsilon}{\gamma}\dot x^i(t).
\]
Then, by construction,
\[
z(0)=z_0,\qquad z(T_1)=0.
\]
Finally, again define
\[
h_i(t)=\frac1{a\sigma}\Bigl(\dot z^i(t)-a(\partial_{x_i}f(x(t))-z^i(t))\Bigr).
\]
Since \(x\) is smooth, \(y\) is \(C^1\), and \(z\) is therefore \(C^1\) on
\([0,T_1]\), the control \(h\) belongs to \(\mathrm L^2([0,T_1];\mathbb R^m)\).
Moreover, the triple \((x(t),z(t),y(t))\) solves the skeleton system on
\([0,T_1]\), and at time \(T_1\) we obtain a state
\begin{equation}\label{state1}
x(T_1)=x^c,\qquad z(T_1)=0,\qquad y(T_1)=:y^c.
\end{equation}
By construction $(x(t),z(t),y(t))\in W'''$ for all $t\in [0,T_1]$.



\begin{remark}
To explain the notation \(x^c\in B(x_\ast,2\rho_x)\), the superscript
\(^c\) denotes an \emph{intermediate control point} (or \emph{control target})
in the construction of the skeleton path. It does not indicate a power or
complement, but serves as a label distinguishing this point from the initial
state \(x_0\) and the final target \(x^\sharp\). More precisely, for each
initial condition \(U_0=(x_0,z_0,y_0)\in W'\) and each desired terminal value
\(y^\sharp\in U_y\), the point
\[
x^c=x^c_{U_0,y^\sharp}\in B(x_\ast,2\rho_x)
\]
will be chosen so that, by holding \(x(t)\equiv x^c\) over a suitable time interval,
the \(y\)-component can be steered exactly to \(y^\sharp\). Thus, \(x^c\) acts
as a control parameter enabling precise matching of the terminal value in the
\(y\)-coordinates.
\end{remark}

To proceed towards the target in \eqref{frame2} on the time interval
\eqref{frame1}, we therefore first focus on the component \(y^\sharp\).
Lemma~\ref{keylem} below shows that we can select, given \(y^\sharp\in U_y\),
a corresponding point \(x^c\in B(x_\ast,2\rho_x)\) such that, starting from
the state \eqref{state1} at time \(T_1\), the system can be steered on
\([T_1,T_1+S]\) to the configuration \((x^c,0,y^\sharp)\). This reduction
of the problem to the choice of \(x^c\) is the key step in controlling the
\(y\)-dynamics. We defer the proof of Lemma~\ref{keylem} to Subsection~\ref{Pr1}.






\begin{lemma}\label{keylem}
Fix \(T_1>0\) as in Lemma~\ref{lem:initial-hermite-path}. Then there exists
\(S>0\) such that the following holds for every \(U_0=(x_0,z_0,y_0)\in W'\). For every \(y^\sharp\in U_y\), there exists
\[
x^c=x^c_{U_0,y^\sharp}\in B(x_\ast,2\rho_x)
\]
such that, if on \([0,T_1]\) we use the controlled trajectory constructed in
Lemma~\ref{lem:initial-hermite-path} with endpoint \(x^c\), and if on
\([T_1,T_1+S]\) we set
\[
x(t)\equiv x^c,\qquad z(t)\equiv 0,\qquad
h(t)=-\sigma^{-1}\nabla f(x^c),
\]
then
\[
y(T_1+S)
=
e^{-bS}\,y^c
+
(1-e^{-bS})\bigl(g(x^c)+\sigma^2\mathbf 1\bigr)
=
y^\sharp,
\]
where \(y^c:=y(T_1)\).
\end{lemma}





Now fix a target triple $(x^\sharp,z^\sharp,y^\sharp)$ as in \eqref{frame2}. Let \(\delta>0\) be a small positive
number, to be chosen, and set
\[
T_2:=T_1+S\implies
T=T_2+\delta.
\]
Here \(S>0\) is as in Lemma~\ref{keylem}. Given \(U_0=(x_0,z_0,y_0)\in W'\) and \(y^\sharp\in U_y\), Lemma~\ref{keylem}
provides a point
\[
x^c=x^c_{U_0,y^\sharp}\in B(x_\ast,2\rho_x)
\]
such that, at time \(T_2\), the corresponding controlled trajectory satisfies
\begin{equation}\label{state2}
x(T_2)=x^c,\qquad z(T_2)=0,\qquad y(T_2)=y^\sharp.
\end{equation}


We next focus on achieving $(x^\sharp,z^\sharp)$ without losing the control achieved in $y$. Given $(z^\sharp,y^\sharp)$ we define the desired terminal velocity for $x$ (at $T$)
\[
v^\sharp
:=
-\gamma\frac{z^\sharp}{\sqrt{y^\sharp}+\varepsilon}.
\]
As this stage we fix $\rho_z$ determining the ball $U_z=B(0,\rho_z)$. Indeed, we let $\rho_z=\min\{\varepsilon/\gamma,1\}$ and then note that
\begin{equation}\label{tribb}
z\in B(0,\rho_z)\implies \|v^\sharp\|\leq \gamma \|z^\sharp\|/\varepsilon<1.
\end{equation}
With $U_z$ fixed, the set $W''=U_x\times U_z\times U_y$ is now completely specified.

In the next lemma we construct a curve \(x^0:[T_2,T]\to\mathbb R^m\), $T=T_2+\delta$, satisfying the four endpoint conditions
\begin{equation}\label{eq:bc-x0}
x^0(T_2)=x^c,\qquad \dot x^0(T_2)=0,\qquad x^0(T)=x^\sharp,\qquad \dot x^0(T)=v^\sharp.
\end{equation}
Hence, by moving along the curve $x^0[T_2,T]$ we have $x^c\to x^\sharp$ and the terminal velocity of $x^0$ at $T$ equals $v^\sharp$. We postpone the proof of the lemma for now, see Subsection \ref{Pr2}.

\begin{lemma}\label{CdG}
Let \(T_2\in\mathbb R\), \(\delta>0\), and \(T=T_2+\delta\). Given $x^c,x^\sharp,v^\sharp\in\mathbb R^m$,
there exists a curve \(x^0:[T_2,T]\to\mathbb R^m\) satisfying the four endpoint
conditions in \eqref{eq:bc-x0}. Moreover, if
\[
x^c\in B(x_\ast,2\rho_x),\qquad x^\sharp\in B(x_\ast,\rho_x),
\]
and
\begin{equation}\label{tribb+}
\delta\|v^\sharp\|\le \frac{27}{8}\rho_x,
\end{equation}
then
\[
x^0([T_2,T])\subset U_x'=B(x_\ast,4\rho_x).
\]
\end{lemma}



To be able to use Lemma~\ref{CdG} to achieve the target $(x^\sharp,z^\sharp,y^\sharp)$ we need more degrees of freedom and we therefore next introduce a finite–dimensional perturbation of the reference curve
constructed in Lemma~\ref{CdG}. We choose a fixed $\psi\in C_c^\infty((0,1))$ such that $\int_0^1 \psi(r)\,\d r\neq 0$,
and define the rescaled bump on $[T_2,T]$ by
\[
\psi_\delta(t):=\frac{1}{\delta}\,\psi\!\Big(\frac{t-T_2}{\delta}\Big),
\qquad t\in[T_2,T],\quad T=T_2+\delta.
\]
For each parameter
\[
\alpha=(\alpha_1,\dots,\alpha_m)\in\R^m,
\]
define
\[
x^\alpha(t)=x^0(t)+\psi_\delta(t)\,\alpha,
\qquad t\in[T_2,T].
\]
Thus \(\alpha\) acts as an \(m\)-dimensional parameter controlling an
interior perturbation of the reference curve \(x^0\) constructed in Lemma~\ref{CdG}.
Because \(\psi\in C_c^\infty((0,1))\), we have \(\psi\equiv0\) near \(0\)
and \(1\). Consequently,
\[
\psi_\delta(T_2)=\psi_\delta(T)=0,
\qquad
\psi_\delta'(T_2)=\psi_\delta'(T)=0.
\]
Hence, the initial and terminal conditions for \(x^\alpha\) coincide with those
of \(x^0\), independently of \(\alpha\), i.e.,
\begin{equation}\label{alphastate}
x^\alpha(T_2)=x^c,\qquad
\dot x^\alpha(T_2)=0,
\qquad
x^\alpha(T)=x^\sharp,
\qquad
\dot x^\alpha(T)=v^\sharp,\qquad \forall\, \alpha=(\alpha_1,\dots,\alpha_m)\in\R^m.
\end{equation}
 Moreover,
\[
\sup_{t\in[T_2,T]}\|x^\alpha(t)-x^0(t)\|
\le \|\psi_\delta\|_\infty\,\|\alpha\|
= \frac{\|\psi\|_\infty}{\delta}\,\|\alpha\|.
\]
As $(x^\sharp,z^\sharp,y^\sharp)\in W''=U_x\times U_z\times U_y$, using \eqref{tribb} we see that \eqref{tribb+} is satisfied if $\delta\leq 27\rho_x/8$ where $\rho_x$ has been fixed at the very beginning of our argument. Hence, subject to this restriction we have $x^0([T_2,T])\subset U_x'=B(x_\ast,4\rho_x)$, see Lemma \ref{CdG}, and since $U_x'$ is open there exists $\eta>0$ such that if
\begin{equation}\label{Bg2}
\|\alpha\|\le r_\alpha(\delta)
:=
\frac{\eta\,\delta}{\|\psi\|_\infty},
\end{equation}
then
\[
x^\alpha([T_2,T])\subset U_x' .
\]
For \(\alpha\) as in \eqref{Bg2} the corresponding solution \(y^\alpha\) of the
\(y\)-equation is well defined on \([T_2,T]\).
We therefore obtain the endpoint map
\[
\Phi_\delta:
B(0,r_\alpha(\delta))\to\mathbb R^m,
\qquad
\Phi_\delta(\alpha):=y^\alpha(T),
\]
where
\[
\Phi_\delta(\alpha)
=
e^{-b\delta}\,y^\sharp
+
b\int_{T_2}^{T}
e^{-b(T-s)}
\big(
g\big(x^\alpha(s)\big)
+
\sigma^2\mathbf 1
\big)\,\d s .
\]
By \eqref{alphastate} we know that $x^\alpha(T_2)=x^c$, $x^\alpha(T)=x^\sharp$. Starting  at $y^\sharp$ at time $T_2$, and following the perturbed reference curve $x^\alpha$,  $\Phi_\delta(\alpha)$ produces the corresponding terminal value $y^\alpha(T)$, i.e., $y^\sharp\to y^\alpha(T)$. The following lemma shows that the parameter \(\alpha\) can be chosen so that, while starting at $y^\sharp$, the desired terminal condition \(y^\alpha(T)=y^\sharp\) is satisfied. For the proof of the lemma we refer to Subsection \ref{Pr3}.
\begin{lemma}\label{CdG+}
There exists $\delta_0\in (0,27\rho_x/8)$, which can be chosen uniformly with respect to
$U_0\in W'$ and $y^\sharp\in U_y$, such that for every $\delta\in(0,\delta_0)$
there exists $\alpha^\ast=\alpha^\ast(\delta)$ with
\[
\|\alpha^\ast\|\le r_\alpha(\delta)=O(\delta)
\]
such that
\[
\Phi_\delta(\alpha^\ast)
=
y^{\alpha^\ast}(T)=
y^\sharp.
\]
\end{lemma}




We now have all the parts of the construction to complete the proof of Lemma \ref{lem:access-xzy-openV-3stagefinal}. To do so, we let $\delta=\delta_0/2$ and we fix \(\alpha^\ast=\alpha^\ast(\delta)\) as in Lemma \ref{CdG+}. We then consider $(x(t),y(t))=(x^{\alpha^\ast}(t),y^{\alpha^\ast}(t))$, and we define $(z(t),h(t))=(z^{\alpha^\ast}(t),h^{\alpha^\ast}(t))$ for \(t\in[T_2,T]\) according to
\[
z^i(t)
=
-\frac{\sqrt{y^i(t)}+\varepsilon}{\gamma}\dot x^i(t),
\qquad
h_i(t)
=
\frac{1}{a\sigma}
\Big(\dot z^i(t)-a(\partial_{x_i}f(x(t))-z^i(t))\Big).
\]
By construction $x^{\alpha^\ast}(T_2)=x^c$, $x^{\alpha^\ast}(T)=x^\sharp$, and $y^{\alpha^\ast}(T_2)=y^\sharp$, $y^{\alpha^\ast}(T)=y^\sharp$. Concerning, $z$ we have $z^{\alpha^\ast}(T_2)=0$ and, with a slight abuse of notation,
$$z(T)=z^{\alpha^\ast}(T)=
-\frac{\sqrt{y^{\alpha^\ast}(T)}+\varepsilon}{\gamma}\dot x^{\alpha^\ast}(T)=-\frac{\sqrt{y^\sharp}+\varepsilon}{\gamma} v^\sharp=\bigl(-\frac{\sqrt{y^\sharp}+\varepsilon}{\gamma}\bigr)\bigl (-\gamma\frac{z^\sharp}{\sqrt{y^\sharp}+\varepsilon}\bigr )=z^\sharp.$$
Thus we have constructed a skeleton system  on \([T_2,T]\), and by construction
\[
(x(T),z(T),y(T))=(x^\sharp,z^\sharp,y^\sharp).
\]

Joining the constructions, given $U_0=(x_0,z_0,y_0)\in W'$ and $(x^\sharp,z^\sharp,y^\sharp)\in W''$, we have constructed controls steering
$$(x_0,z_0,y_0)\to (x^c_{U_0,y^\sharp},0,y^c_{U_0,y^\sharp})\to (x^c_{U_0,y^\sharp},0,y^\sharp)\to (x^\sharp,z^\sharp,y^\sharp)$$
on the intervals $[0,T_1]$, $[T_1,T_2]$, and $[T_2,T]$ respectively. In the last display, the notation $x^c_{U_0,y^\sharp}$ emphasizes that,
given the initial condition $U_0$ and the target value $y^\sharp$,
one can select a point $x^c_{U_0,y^\sharp}$ such that, if the system is steered from $x_0$ to $x^c_{U_0,y^\sharp}$ over $[0,T_1]$, then the subsequent evolution on $[T_1,T_2]$ ensures that $y(T_2)=y^\sharp$.




We can conclude that the set $W''=U_x\times U_z\times U_y$
is open and every point of \(W''\) is attainable from every
\(U_0\in W'\).
This completes the proof of Lemma~\ref{lem:access-xzy-openV-3stagefinal}
modulo the auxiliary lemmas. \qed

\subsection{Proof of Lemma \ref{lem:initial-hermite-path}}\label{Pr0} Set $s:={t}/{T_1}\in[0,1]$.  We use the standard cubic Hermite basis functions
\[
h_{00}(s)=2s^3-3s^2+1,\qquad
h_{10}(s)=s^3-2s^2+s,
\]
\[
h_{01}(s)=-2s^3+3s^2,\qquad
h_{11}(s)=s^3-s^2.
\]
Define
\begin{equation}\label{eq:initial-hermite}
x(t):=
h_{00}(s)\,x_0
+
T_1 h_{10}(s)\,v_0
+
h_{01}(s)\,x^c,
\qquad s={t}/{T_1}.
\end{equation}
Since \(h_{00},h_{10},h_{01}\) are polynomials, the curve \(x\) is \(C^\infty\)
on \([0,T_1]\). We verify the endpoint conditions. Using
\[
h_{00}(0)=1,\quad h_{10}(0)=0,\quad h_{01}(0)=0,
\]
and
\[
h_{00}(1)=0,\quad h_{10}(1)=0,\quad h_{01}(1)=1,
\]
we obtain
\[
x(0)=x_0,\qquad x(T_1)=x^c.
\]
Differentiating \eqref{eq:initial-hermite}  we get
\[
\dot x(t)
=
\frac1{T_1}\Big(h_{00}'(s)x_0 + T_1 h_{10}'(s)v_0 + h_{01}'(s)x^c\Big).
\]
Since
\[
h_{00}'(0)=0,\quad h_{10}'(0)=1,\quad h_{01}'(0)=0,
\]
and
\[
h_{00}'(1)=0,\quad h_{10}'(1)=0,\quad h_{01}'(1)=0,
\]
it follows that
\[
\dot x(0)=v_0,\qquad \dot x(T_1)=0.
\]
It remains to prove that \(x([0,T_1])\subset B(x_\ast,2r)\) for \(T_1>0\) chosen
sufficiently small, uniformly with respect to \(U_0\in W'\) and
\(x^c\in B(x_\ast,2\rho_x)\). From \eqref{eq:initial-hermite} and the identity
\[
h_{00}(s)+h_{01}(s)=1,
\]
we deduce
\[
x(t)-x_\ast
=
h_{00}(s)(x_0-x_\ast)
+
h_{01}(s)(x^c-x_\ast)
+
T_1 h_{10}(s)v_0.
\]
Taking norms and using \(0\le h_{00}(s),h_{01}(s)\le 1\), we obtain
\[
\|x(t)-x_\ast\|
\le
h_{00}(s)\|x_0-x_\ast\|
+
h_{01}(s)\|x^c-x_\ast\|
+
T_1 |h_{10}(s)|\,\|v_0\|.
\]
Since \(x_0\in B(x_\ast,r)\) and
\(x^c\in B(x_\ast,2\rho_x)\subset B(x_\ast,r)\), we have
\[
\|x_0-x_\ast\|<r,\qquad \|x^c-x_\ast\|<r.
\]
Therefore,
\[
h_{00}(s)\|x_0-x_\ast\|+h_{01}(s)\|x^c-x_\ast\|
<
r\big(h_{00}(s)+h_{01}(s)\big)
=
r.
\]
Next, because \(W'\) is bounded and \(y_0^i\) is bounded away from zero on \(W'\),
there exists a constant \(M_v<\infty\) such that the initial velocity $v_0$ satisfies
\[
\|v_0\|\le M_v
\qquad
\text{for all }U_0=(x_0,z_0,y_0)\in W'.
\]
Hence
\[
\|x(t)-x_\ast\|
<
r + T_1\|h_{10}\|_{L^\infty(0,1)}\,M_v.
\]
We now choose \(T_1>0\) so small that
\[
T_1\|h_{10}\|_{L^\infty(0,1)}\,M_v<r.
\]
Then, for every \(U_0\in W'\) and every \(x^c\in B(x_\ast,2\rho_x)\),
\[
\|x(t)-x_\ast\|<2r
\qquad \text{for all } t\in[0,T_1].
\]
Thus
\[
x([0,T_1])\subset B(x_\ast,2r),
\]
uniformly with respect to \(U_0\in W'\) and \(x^c\in B(x_\ast,2\rho_x)\). \qed

\subsection{Proof of Lemma \ref{keylem}}\label{Pr1} Fix \(U_0=(x_0,z_0,y_0)\in W'\). For each \(x\in B(x_\ast,2\rho_x)\), let $y^c(U_0,x)$
denote the value at time \(T_1\) of the \(y\)-component obtained by using on
\([0,T_1]\) the controlled trajectory from Lemma~\ref{lem:initial-hermite-path}
with terminal \(x(T_1)=x\). By construction, $T_1$is such that this trajectory satisfies
\[
x([0,T_1])\subset B(x_\ast,2r).
\]
In particular, \(y^c(U_0,x)\) is well-defined for every
\(x\in B(x_\ast,2\rho_x)\). Now fix \(y^\sharp\in U_y\). If on \([T_1,T_1+S]\) we keep \(x(t)\equiv x\),
\(z(t)\equiv 0\), and choose $h(t)=-\sigma^{-1}\nabla f(x)$,
then the \(y\)-equation becomes
\[
\dot y(t)=b\bigl(-y(t)+g(x)+\sigma^2\mathbf 1\bigr),
\qquad y(T_1)=y^c(U_0,x).
\]
Hence
\begin{equation}\label{eq:y-relax-final-rev}
y(T_1+S)
=
e^{-bS}\,y^c(U_0,x)
+
(1-e^{-bS})\bigl(g(x)+\sigma^2\mathbf 1\bigr).
\end{equation}
We seek \(x^c\in B(x_\ast,2\rho_x)\) such that \(y(T_1+S)=y^\sharp\). By
\eqref{eq:y-relax-final-rev}, this is equivalent to
\begin{equation}\label{eq:implicit-xc-rev}
g(x^c)
=
\frac{y^\sharp-e^{-bS}y^c(U_0,x^c)}{1-e^{-bS}}
-\sigma^2\mathbf 1.
\end{equation}


We are now going to use that \(Dg(x)=2A(x)\) is invertible on \(U_x'=B(x_\ast,4\rho_x)\). Indeed, using this fact we have that the map
\[
g:U_x'\to g(U_x')
\]
is a \(C^1\)-diffeomorphism. In particular,
\[
g^{-1}:g(U_x')\to U_x'
\]
is \(C^1\). Define
\[
\mathcal F_{U_0}(x)
:=
g^{-1}\!\left(
\frac{y^\sharp-e^{-bS}y^c(U_0,x)}{1-e^{-bS}}
-\sigma^2\mathbf 1
\right),
\qquad x\in B(x_\ast,2\rho_x).
\]
We claim that
$$\mathcal F_{U_0}(B(x_\ast,2\rho_x))\subset B(x_\ast,2\rho_x),$$
i.e., \(\mathcal F_{U_0}\) maps \(B(x_\ast,2\rho_x)\) into
itself, for \(S\) large enough, uniformly in \(U_0\in W'\). To see this we first note that since
\[
U_y\subset g(B(x_\ast,2\rho_x))+\sigma^2\mathbf 1,
\]
we have that there exist, for every \(y^\sharp\in U_y\),
\[
x^\sharp_g:=g^{-1}(y^\sharp-\sigma^2\mathbf 1)\in B(x_\ast,2\rho_x).
\]
Because \(B(x_\ast,2\rho_x)\Subset U_x'\), there exists \(\delta_0>0\) such that
the closed \(\delta_0\)-neighborhood of \(B(x_\ast,2\rho_x)\) is still contained
in \(U_x'\). Since \(g^{-1}\) is continuous on the compact set
\[
g(\overline{B(x_\ast,2\rho_x)})+\overline{B(0,\delta_1)}
\]
for sufficiently small \(\delta_1>0\), it follows that if the argument of
\(g^{-1}\) is within \(\delta_1\) of \(y^\sharp-\sigma^2\mathbf 1\), then
\(\mathcal F_{U_0}(x)\in B(x_\ast,2\rho_x)\).

Using that \(W'\) is bounded and that the first-stage paths $x[0,T]$ stay in \(B(x_\ast,2r)\), there exists \(M_y<\infty\) such that
\[
\|y^c(U_0,x)\|\le M_y
\qquad
\text{for all }U_0\in W',\ x\in B(x_\ast,2\rho_x).
\]
Therefore,
\[
\left\|
\frac{y^\sharp-e^{-bS}y^c(U_0,x)}{1-e^{-bS}}
-y^\sharp
\right\|
\le
\frac{e^{-bS}}{1-e^{-bS}}\bigl(\|y^\sharp\|+M_y\bigr).
\]
Since \(y^\sharp\in U_y\), the quantity \(\|y^\sharp\|\) is uniformly bounded on
\(U_y\). Hence, for \(S>0\) sufficiently large, uniformly in \(U_0\in W'\),
\(x\in B(x_\ast,2\rho_x)\), and \(y^\sharp\in U_y\), the argument of \(g^{-1}\)
belongs to \(g(U_x')\) and \(\mathcal F_{U_0}(x)\in B(x_\ast,2\rho_x)\).

Next we show that \(\mathcal F_{U_0}\) is a contraction for \(S\) large enough,
uniformly in \(U_0\in W'\). Since \(g^{-1}\) is \(C^1\) on a neighborhood of the
compact set \(g(\overline{B(x_\ast,2\rho_x)})\), there exists \(C_g>0\) such
that $\|Dg^{-1}\|\le C_g$ there. Hence
\begin{equation}\label{Bg3}
\|\mathcal F_{U_0}(x)-\mathcal F_{U_0}(\tilde x)\|
\le
C_g\,\frac{e^{-bS}}{1-e^{-bS}}\,
\|y^c(U_0,x)-y^c(U_0,\tilde x)\|.
\end{equation}
Moreover, for fixed \(T_1>0\), the map
\[
x\mapsto y^c(U_0,x)
\]
is Lipschitz on \(B(x_\ast,2\rho_x)\), uniformly in \(U_0\in W'\), i.e., there exists
\(L_{T_1}<\infty\) such that
\[
\|y^c(U_0,x)-y^c(U_0,\tilde x)\|
\le
L_{T_1}\|x-\tilde x\|,
\qquad
x,\tilde x\in B(x_\ast,2\rho_x),
\]
for all \(U_0\in W'\). Combining this with \eqref{Bg3} we deduce
\[
\|\mathcal F_{U_0}(x)-\mathcal F_{U_0}(\tilde x)\|
\le
C_g\,\frac{e^{-bS}}{1-e^{-bS}}\,L_{T_1}\,\|x-\tilde x\|.
\]
We now choose \(S>0\) so large that
\[
C_g\,\frac{e^{-bS}}{1-e^{-bS}}\,L_{T_1}<1.
\]
Then \(\mathcal F_{U_0}\) is a contraction on \(B(x_\ast,2\rho_x)\), uniformly in
\(U_0\in W'\). By Banach's fixed point theorem, the map $x\to \mathcal F_{U_0}(x)$ admits a unique fixed point
\[
x^c=x^c_{U_0,y^\sharp}\in B(x_\ast,2\rho_x).
\]
This fixed point satisfies \eqref{eq:implicit-xc-rev}, and substituting into
\eqref{eq:y-relax-final-rev} yields
\[
y(T_1+S)=y^\sharp.
\]
This completes the proof of the lemma.\qed


\subsection{Proof of Lemma \ref{CdG}}\label{Pr2}  The proof is similar to the proof of Lemma \ref{lem:initial-hermite-path} but subject to differences. We here give the complete proof. Set $s:={(t-T_2)}/{\delta}\in[0,1]$. We again use the cubic Hermite basis functions
\[
h_{00}(s)=2s^3-3s^2+1,\qquad
h_{10}(s)=s^3-2s^2+s,
\]
\[
h_{01}(s)=-2s^3+3s^2,\qquad
h_{11}(s)=s^3-s^2.
\]
Define \(x^0\) by the vector-valued cubic Hermite interpolant
\begin{equation}\label{eq:x0-hermite}
x^0(t)
:=
h_{00}(s)\,x^c
+\delta\,h_{10}(s)\,\dot x^0(T_2)
+h_{01}(s)\,x^\sharp
+\delta\,h_{11}(s)\,\dot x^0(T),
\qquad s=\frac{t-T_2}{\delta}.
\end{equation}
With the choices
\[
\dot x^0(T_2)=0,\qquad \dot x^0(T)=v^\sharp,
\]
this simplifies to
\begin{equation}\label{eq:x0-hermite-simplified}
x^0(t)
=
h_{00}(s)\,x^c
+h_{01}(s)\,x^\sharp
+\delta\,h_{11}(s)\,v^\sharp,
\qquad s=\frac{t-T_2}{\delta}.
\end{equation}
Since \(h_{00},h_{01},h_{11}\) are polynomials, we have $x^0\in C^\infty([T_2,T];\mathbb R^m)$. We now verify \eqref{eq:bc-x0}. Using
\[
h_{00}(0)=1,\quad h_{01}(0)=0,\quad h_{11}(0)=0,
\]
and
\[
h_{00}(1)=0,\quad h_{01}(1)=1,\quad h_{11}(1)=0,
\]
it follows directly from \eqref{eq:x0-hermite-simplified} that
\[
x^0(T_2)=x^c,\qquad x^0(T)=x^\sharp.
\]
Differentiating \eqref{eq:x0-hermite} we obtain
\[
\dot x^0(t)
=
\frac{1}{\delta}
\Big(
h_{00}'(s)\,x^c
+\delta h_{10}'(s)\,\dot x^0(T_2)
+h_{01}'(s)\,x^\sharp
+\delta h_{11}'(s)\,\dot x^0(T)
\Big).
\]
Since
\[
h_{00}'(0)=0,\quad h_{10}'(0)=1,\quad h_{01}'(0)=0,\quad h_{11}'(0)=0,
\]
and
\[
h_{00}'(1)=0,\quad h_{10}'(1)=0,\quad h_{01}'(1)=0,\quad h_{11}'(1)=1,
\]
we get
\[
\dot x^0(T_2)=\dot x^0(T_2),\qquad \dot x^0(T)=\dot x^0(T).
\]
Hence, with the above choices,
\[
\dot x^0(T_2)=0,\qquad \dot x^0(T)=v^\sharp,
\]
and the four endpoint conditions in \eqref{eq:bc-x0} are satisfied. It remains to prove the inclusion \(x^0([T_2,T])\subset U_x'\). From
\eqref{eq:x0-hermite-simplified} and the identity
\[
h_{00}(s)+h_{01}(s)=1,
\]
we obtain
directly from
\[
\|x^0(t)-x_\ast\|
\le
|h_{00}(s)|\,\|x^c-x_\ast\|
+
|h_{01}(s)|\,\|x^\sharp-x_\ast\|
+
\delta\,|h_{11}(s)|\,\|v^\sharp\|.
\]
Since
\[
0\le h_{00}(s), h_{01}(s)\le 1,
\qquad
\sup_{s\in[0,1]}|h_{11}(s)|=\frac{4}{27}.
\]
it follows that
\begin{equation}\label{eq:x0-bound-revised}
\|x^0(t)-x_\ast\|
\le
\|x^c-x_\ast\|
+
\|x^\sharp-x_\ast\|
+
\frac{4}{27}\,\delta\,\|v^\sharp\|,
\qquad t\in[T_2,T].
\end{equation}
Now assume
\[
x^c\in B(x_\ast,2\rho_x),\qquad x^\sharp\in B(x_\ast,\rho_x).
\]
Then
\[
\|x^0(t)-x_\ast\|
\le 2\rho_x+\rho_x+
\frac{4}{27}\,\delta\,\|v^\sharp\|,
\qquad t\in[T_2,T].
\]
Therefore, if in addition
\[
\delta\|v^\sharp\|\le \frac{27}{8}\rho_x,
\]
then
\[\|x^0(t)-x_\ast\|
\le 3\rho_x+
\frac{4}{27}\cdot \frac{27}{8}\rho_x\le \frac 72\rho_x<4\rho_x
\qquad t\in[T_2,T].
\]
We can conclude that
\[
\|x^0(t)-x_\ast\|<4\rho_x
\qquad\text{for all }t\in[T_2,T],
\]
and hence
\[
x^0([T_2,T])\subset U_x'.
\]
This proves the lemma.\qed





\subsection{Proof of Lemma \ref{CdG+}}\label{Pr3} Standard smooth dependence of ODE solutions on parameters implies that
\(\Phi_\delta\) is \(C^1\) on \(B(0,r_\alpha(\delta))\).
Differentiating the map \(\alpha\mapsto\Phi_\delta(\alpha)\) yields
\[
D\Phi_\delta(0)
=
b\int_{T_2}^{T}
e^{-b(T-s)}
Dg(x^0(s))\,\psi_\delta(s)\,\d s .
\]
Using the change of variables \(s=T_2+\delta r\) gives
\[
D\Phi_\delta(0)
=
b\int_0^1
e^{-b\delta(1-r)}
Dg(x^0(T_2+\delta r))\,\psi(r)\,\d r .
\]
Since
\[
\sup_{r\in[0,1]}
\|x^0(T_2+\delta r)-x^c\|\to0
\qquad (\delta\to0),
\]
we have
\[
x^0(T_2+\delta r)\to x^c
\quad\text{uniformly in } r\in[0,1].
\]
Because \(Dg\) is continuous,
\[
D\Phi_\delta(0)
\longrightarrow
b\left(\int_0^1\psi(r)\,\d r\right)Dg(x^c)
\qquad (\delta\to0).
\]
Since
\[
\int_0^1\psi(r)\,\d r\neq0
\]
and \(Dg(x^c)=2A(x^c)\) is invertible, the limit matrix is invertible.
Hence there exists \(\delta_0>0\) such that \(D\Phi_\delta(0)\) is
invertible for all \(\delta\in(0,\delta_0)\). Fix such a \(\delta\). Then, by the inverse function theorem there exists
\(\eta_\delta>0\) such that
\[
\Phi_\delta:
B(0,\eta_\delta)\to
\Phi_\delta(B(0,\eta_\delta))
\]
is a diffeomorphism onto an open neighborhood of \(\Phi_\delta(0)\). Since \(x^0(t)\to x^c\) uniformly as \(\delta\to0\), the solution of the
\(y\)-equation driven by \(x^0\) satisfies
\[
\|\Phi_\delta(0)-y^\sharp\|
=
\|y^0(T)-y^\sharp\|
\to0
\qquad (\delta\to0).
\]
Therefore, for sufficiently small \(\delta\), the point \(y^\sharp\)
belongs to the neighborhood
\(\Phi_\delta(B(0,\eta_\delta))\).
Hence there exists $\alpha^\ast\in B(0,\eta_\delta)$ such that
\[
\Phi_\delta(\alpha^\ast)=y^\sharp .
\]
Since \(r_\alpha(\delta)=O(\delta)\), we may choose
\(\alpha^\ast\) satisfying
\[
\|\alpha^\ast\|\le r_\alpha(\delta).
\]
This completes the construction of $\alpha^\ast$.

To complete the proof of the lemma we next verify that the constant $\delta_0$ can be chosen uniformly
with respect to the parameters entering the construction. Recall that $x^c\in B(x_\ast,2\rho_x)\subset U_x'$, where
\[
U_x'=B(x_\ast,4\rho_x)\subset\R^m\setminus\mathcal D_A^\dagger.
\]
Since $Dg(x)=2A(x)$ and $A$ is continuous, it follows that $Dg$ is continuous
on $\overline{B(x_\ast,2\rho_x)}$. Moreover, since
$\det A(x)\neq 0$ on $U_x'$, we have
\[
\det Dg(x)\neq 0
\qquad \text{for all } x\in \overline{B(x_\ast,2\rho_x)}.
\]
By compactness, there exists a constant $\lambda_\ast>0$ such that
\begin{equation}\label{eq:uniform-invertibility}
\|Dg(x)^{-1}\|\le \lambda_\ast^{-1}
\qquad \text{for all } x\in \overline{B(x_\ast,2\rho_x)}.
\end{equation}
In particular, $Dg$ is uniformly invertible on this set. Next, since $Dg$ is continuous on a compact set, it is uniformly continuous.
Hence, for every $\varepsilon_0>0$ there exists $\delta_0>0$ such that
\[
\|x-x^c\|\le \delta_0
\quad\Longrightarrow\quad
\|Dg(x)-Dg(x^c)\|\le \varepsilon_0
\]
uniformly for all $x^c\in B(x_\ast,2\rho_x)$. From the construction of the curve $x^0$ (see Lemma~\ref{CdG}) we have
\[
\sup_{r\in[0,1]}
\|x^0(T_2+\delta r)-x^c\|\;\longrightarrow\;0
\qquad \text{as } \delta\to 0,
\]
uniformly with respect to $x^c\in B(x_\ast,2\rho_x)$ and
$x^\sharp\in B(x_\ast,\rho_x)$. Therefore, for $\delta>0$ sufficiently small (independently of $x^c$),
we obtain
\[
\sup_{r\in[0,1]}
\|Dg(x^0(T_2+\delta r))-Dg(x^c)\|
\le \varepsilon_0.
\]
Choosing $\varepsilon_0>0$ sufficiently small relative to the
uniform bound \eqref{eq:uniform-invertibility}, it follows that
$D\Phi_\delta(0)$ remains invertible for all $\delta\in(0,\delta_0)$,
with constants independent of $x^c$. Finally, the constraint ensuring that
$x^\alpha([T_2,T])\subset U_x$ is controlled by the condition
$\|\alpha\|\le r_\alpha(\delta)=c_0\delta$, where $c_0$ depends only on
$\rho_x$ and $\|\psi\|_\infty$, and is therefore independent of $x^c$. Collecting the above, we conclude that $\delta_0$ can be chosen
uniformly with respect to $x^c\in B(x_\ast,2\rho_x)$,
$y^\sharp\in U_y$, and the initial condition $U_0\in W'$. This completes the proof of the lemma.\qed













\section{Concluding remarks and future research}\label{Conc}

In this paper, we developed an effective continuous-time model for the long-term dynamics of adaptive stochastic optimization, focusing on bias-corrected Adam-type methods. Starting from the finite-sum setting, we identified a canonical scaling and an associated effective closure under which the discrete dynamics give rise to a coupled, time-inhomogeneous SDE governing the parameters, the first-moment tracker, and the per-coordinate second-moment tracker. The bias-correction mechanism persists in the limit through explicit time-dependent coefficients, reflecting its finite-time influence on the dynamics.

Building on this, we analyzed the long-time behavior of the limiting dynamics and studied the existence and uniqueness of invariant measures.  In Theorem~\ref{CorC} we proved, under mild regularity and dissipativity assumptions on the objective function $f$, existence, uniqueness and exponential convergence to equilibrium. An important observation was that our assumptions imply $\mathcal D_A^\dagger \neq \mathbb R^m$, where $A$ refers to the defining matrix $A(x)=\mathrm{Diag}(\nabla f(x))\,H_f(x)$ describing how noise propagates through the dynamics.  $\mathcal D_A^\dagger \neq \mathbb R^m\implies \mathcal D_A \neq \mathbb R^m$ and the latter guaranteed the existence of a region of the state space in which the H\"ormander bracket condition held. This ensured that the system possessed nondegenerate directions of noise propagation and that hypoellipticity could be exploited locally. Still, this had to be complemented by a construction of control paths compensating for the lack of global hypoellipticity, and this was precisely where the stronger condition $\mathcal D_A^\dagger \neq \mathbb R^m$ entered. Underlying our proof of Theorem~\ref{CorC} was a Harris-type argument combined with a minorization condition on Lyapunov sublevel sets, see Theorem~\ref{Thm3-revised-}.

To our knowledge,  Theorem~\ref{CorC}, and Theorem~\ref{Thm3-revised-} are new, although ours proofs, from a top-down perspective, follow classical routes. We nevertheless provided a largely self-contained analysis. In general, the central theme is the study of estimates, regularity, uniqueness, and exponential convergence to equilibrium in settings where (global) hypoellipticity and positivity of densities fail.


The literature devoted to this theme  seem to be quite limited. An interesting recent paper touching on the topics in a context somewhat related to ours, is \cite{DouSalortSmets2025}. In \cite{DouSalortSmets2025} the authors
consider  the linear partial differential equation
\[
\partial_t u - \partial_{xx} u + V(x)\,\partial_y u = 0
\]
for the unknown function
\[
u \equiv u(t,x,y) : \mathbb{R}_+ \times \mathbb{T} \times \mathbb{T} \to \mathbb{R},
\]
where \(V:\mathbb{T}\to\mathbb{R}\) is a given bounded Borel measurable function, and \(\mathbb{T}=\mathbb{R}/\mathbb{Z}\) denotes the unit flat torus. The equation reflects the interplay between linear diffusion in the variable
\(x\) and transport with velocity \(V(x)\) in the transverse variable \(y\). If  \(V\) is smooth, taking iterative commutators reveals that hypoellipticity
occurs when the critical points of \(V\) are at most finitely degenerate. Consequently, if $V$ is constant on some intervals, the equation is infinitely degenerate. Still, in this (simple) context, the authors obtain quantitative mixing-rate estimates under minimal assumptions on the transport field $V$, which may be
highly irregular and/or highly degenerate, and in particular need not satisfy
any hypoellipticity-type condition. We believe that this is an interesting direction for future research.

An other interesting strain of research is represented by \cite{BallyCaramellino2013,BallyCaramellinoPigato2022a,BallyCaramellinoPigato2022b,Pigato2022,Pigato2017} and some of the references therein. The common theme of these works is the study of \emph{tubular estimates for transition densities} of diffusion processes under various regularity and non-degeneracy assumptions on the coefficients. More precisely, the goal is to obtain lower and upper bounds for the density of a diffusion process in $\mathbb{R}^n$ over small (but non-asymptotic) time intervals. The diffusion coefficients may be degenerate at the initial point, but the system satisfies suitable non-degeneracy conditions such as the strong or weak H\"ormander condition. A key feature of these estimates is that they reflect the \emph{anisotropic geometry} induced by the vector fields generating the diffusion. In particular, the bounds are expressed in terms of norms adapted to the Lie algebra generated by the diffusion vector fields. At small times $\delta$, the process propagates at speed $\sqrt{\delta}$ in the directions of the diffusion vector fields $\sigma_j$, while propagation in the directions of the  Lie brackets  occurs at the slower scale. Such density estimates naturally lead to \emph{tube estimates}, i.e., quantitative bounds for the probability that the diffusion remains within a tubular neighborhood of a deterministic skeleton path over a finite time interval. We believe that this is also an interesting direction to study and revisit in future research,  assuming that hypoellipticity fails on some small sets.

From a practical and numerical perspective, it is of interest to characterize the invariant measure $\pi_\infty$, as it provides a rigorous description of the model’s asymptotic statistical behavior. In particular, $\pi_\infty$ describes how the optimizer concentrates near minimizers of $f$, and how dispersion and correlations influence generalization, escape from saddle points, and metastability in high-dimensional regimes. For comparison, in the overdamped Langevin setting with damping parameter $\gamma>0$, the invariant distribution is given by the Gibbs measure
\[
\pi_{\mathrm{od}}(\mathrm{d}x)\propto\exp\!\Bigl(-\tfrac{2\gamma}{\sigma^2}f(x)\Bigr)\,\mathrm{d}x.
\]
For the heavy-ball model, the associated underdamped Langevin dynamics admits the Gibbs law
\[
\pi_{\mathrm{ud}}(\mathrm{d}x\,\mathrm{d}z)\propto \exp\!\Big(-\tfrac{2}{\gamma\sigma^2}\,f(x)\Big)\;
\exp\!\Big(-\tfrac{1}{a\sigma^2}\,\|z\|^2\Big)\,\mathrm{d}x\,\mathrm{d}z.
\]
In Theorem~\ref{CorC}, we established, under suitable assumptions, that the system \eqref{eq:cts-x+}-\eqref{eq:cts-y+} admits a unique invariant measure $\pi_\infty$. Moreover, Corollary~\ref{correa} shows that the long-time behavior of the inhomogeneous system \eqref{eq:cts-x}-\eqref{eq:cts-y} is also governed by $\pi_\infty$. In contrast to the Langevin models, $\pi_\infty$ does not admit a closed-form Gibbs representation. Instead, it reflects anisotropic couplings and nontrivial correlations between $(x,z,y)$ induced by adaptivity and bias correction. To approximate $\pi_\infty$, it may therefore be natural to develop a conditional–Gaussian (Hermite-Galerkin) ansatz in the fast variable $z$, exploiting its approximate Ornstein-Uhlenbeck structure around $\nabla f(x)$.







   \begin{thebibliography}{99}






\bibitem{AgrachevSachkov2004}
A.~A.~Agrachev and Y.~L.~Sachkov.
\newblock \emph{Control Theory from the Geometric Viewpoint}.
\newblock Encyclopaedia of Mathematical Sciences, Vol.~87, Springer, Berlin, 2004.





\bibitem{Bach2024}
F.~Bach.
\newblock \emph{Learning Theory from First Principles}.
\newblock MIT Press, Cambridge, MA, 2024.

\bibitem{BakryGentilLedoux2014}
D.~Bakry, I.~Gentil, and M.~Ledoux.
\newblock \emph{Analysis and Geometry of Markov Diffusion Operators}.
\newblock Springer, Cham, 2014.



\bibitem{BallyCaramellino2013}
V.~Bally and L.~Caramellino.
\newblock Positivity and lower bounds for the density of Wiener functionals.
\newblock \emph{Potential Anal.} \textbf{39} (2013), 141-168.



\bibitem{BallyCaramellinoPigato2022a}
V.~Bally, L.~Caramellino, and P.~Pigato.
\newblock Diffusions under a local strong H\"ormander condition. Part I: density estimates.
\newblock \emph{Stochastic Process. Appl.} \textbf{145} (2022), 117-142.

\bibitem{BallyCaramellinoPigato2022b}
V.~Bally, L.~Caramellino, and P.~Pigato.
\newblock Diffusions under a local strong H\"ormander condition. Part II: tube estimates.
\newblock \emph{Stochastic Process. Appl.} \textbf{145} (2022), 143-185.



\bibitem{BogachevKrylovRoeckner2015}
V.~I.~Bogachev, N.~V.~Krylov, and M.~R\"ockner.
\newblock \emph{Fokker-Planck-Kolmogorov Equations}.
\newblock American Mathematical Society, Providence, RI, 2015.

\bibitem{BouleauHirsch1991}
N.~Bouleau and F.~Hirsch.
\newblock \emph{Dirichlet Forms and Analysis on Wiener Space}.
\newblock De Gruyter Studies in Mathematics, Vol.~14.
\newblock Walter de Gruyter, Berlin, 1991.

\bibitem{DouSalortSmets2025}
X.~Dou, D.~Salort, and D.~Smets.
\newblock On partial diffusion and mixing without hypoellipticity.
\newblock \emph{arXiv:2511.05280}, 2025.

\bibitem{DereichJentzenRiekert2025}
S.~Dereich, A.~Jentzen, and A.~Riekert.
\newblock Sharp higher order convergence rates for the Adam optimizer.
\newblock \emph{arXiv preprint}, 2025.

\bibitem{Dozat2016}
T.~Dozat.
\newblock Incorporating Nesterov momentum into Adam.
\newblock In \emph{Proc. 4th Int. Conf. on Learning Representations (ICLR)}, Workshop Track, 2016.

\bibitem{Doob1953}
J.~L.~Doob.
\newblock \emph{Stochastic Processes}.
\newblock Wiley, New York, 1953.



\bibitem{Duchi2011}
J.~C.~Duchi, E.~Hazan, and Y.~Singer.
\newblock Adaptive subgradient methods for online learning and stochastic optimization.
\newblock \emph{J. Mach. Learn. Res.} \textbf{12} (2011), 2121-2159.

\bibitem{E2020}
W.~E, C.~Ma, S.~Wojtowytsch, and L.~Wu.
\newblock Towards a mathematical understanding of neural network-based machine learning: What we know and what we don't.
\newblock \emph{arXiv:2009.10713}, 2020.

\bibitem{Eberle2016}
A.~Eberle.
\newblock Reflection coupling and contraction rates for diffusions.
\newblock \emph{Probab. Theory Related Fields} \textbf{166} (2016), 851-886.

\bibitem{Eberle2019}
A.~Eberle, A.~Guillin, and R.~Zimmer.
\newblock Coupling and quantitative contraction rates for Langevin dynamics.
\newblock \emph{Ann. Probab.} \textbf{47} (2019), 1982-2010.

\bibitem{ethier-kurtz}
S.~N.~Ethier and T.~G.~Kurtz.
\newblock \emph{Markov Processes: Characterization and Convergence}.
\newblock Wiley, New York, 1986.

\bibitem{GarrigosGower2023}
G.~Garrigos and R.~M.~Gower.
\newblock Handbook of convergence theorems for (stochastic) gradient methods.
\newblock \emph{arXiv:2301.11235}, 2023.

\bibitem{goyal2017accurate}
P.~Goyal et al.
\newblock Accurate, large minibatch SGD: Training ImageNet in 1 hour.
\newblock \emph{arXiv:1706.02677}, 2017.



\bibitem{HairerMattinglyScheutzow2011}
M.~Hairer, J.~C.~Mattingly, and M.~Scheutzow.
\newblock Asymptotic coupling and a general form of Harris' theorem with applications to stochastic delay equations.
\newblock \emph{Probab. Theory Related Fields} \textbf{149} (2011), 223-259.

\bibitem{HairerMattingly2011b}
M.~Hairer and J.~C.~Mattingly.
\newblock Yet another look at Harris' ergodic theorem for Markov chains.
\newblock In \emph{Seminar on Stochastic Analysis, Random Fields and Applications VI}, 109-117.
\newblock Birkh\"auser, Basel, 2011.

\bibitem{Hasminskii2012}
R.~Z.~Has'minskii.
\newblock \emph{Stochastic Stability of Differential Equations}.
\newblock 2nd ed., Springer, Berlin, 2012.

\bibitem{Hardt2016TrainSGD}
M.~Hardt, T.~Ma, and B.~Recht.
\newblock Gradient descent learns linear dynamical systems.
\newblock \emph{J. Mach. Learn. Res.} \textbf{19} (2018), 1-44.

\bibitem{Harris1956}
T.~E.~Harris.
\newblock The existence of stationary measures for certain Markov processes.
\newblock In \emph{Proc. Third Berkeley Symp. on Mathematical Statistics and Probability}, Vol.~2, 113-124.
\newblock University of California Press, 1956.


\bibitem{HerzogMattingly2015}
D.~P. Herzog and J.~C. Mattingly.
\newblock A practical criterion for positivity of transition densities.
\newblock \emph{Nonlinearity} \textbf{28} (2015), 2823-2845.

\bibitem{Hormander1967}
L.~H\"ormander.
\newblock Hypoelliptic second order differential equations.
\newblock \emph{Acta Math.} \textbf{119} (1967), 147-171.

\bibitem{jacod-shiryaev}
J.~Jacod and A.~N.~Shiryaev.
\newblock \emph{Limit Theorems for Stochastic Processes}.
\newblock 2nd ed., Springer, Berlin, 2003.

\bibitem{JentzenIntroDL2023}
A.~Jentzen, B.~Kuckuck, and P.~von Wurstemberger.
\newblock Mathematical introduction to deep learning: Methods, implementations, and theory.
\newblock \emph{arXiv:2310.20360}, 2023.

\bibitem{Jurdjevic1997}
V.~Jurdjevic.
\newblock \emph{Geometric Control Theory}.
\newblock Cambridge Studies in Advanced Mathematics, Vol.~52, Cambridge University Press, Cambridge, 1997.

\bibitem{Karimi2016PL}
H.~Karimi, J.~Nutini, and M.~Schmidt.
\newblock Linear convergence of gradient and proximal-gradient methods under the Polyak-Lojasiewicz condition.
\newblock In \emph{Proc. AISTATS}, 795-803, 2016.

\bibitem{KingmaBa2015}
D.~P.~Kingma and J.~Ba.
\newblock Adam: A method for stochastic optimization.
\newblock In \emph{Proc. ICLR}, 2015.

\bibitem{kurtz-protter}
T.~G.~Kurtz and P.~E.~Protter.
\newblock Weak limit theorems for stochastic integrals and stochastic differential equations.
\newblock \emph{Ann. Probab.} \textbf{19} (1991), 1035-1070.

\bibitem{li2019stochastic}
Q.~Li, C.~Tai, and W.~E.
\newblock Stochastic modified equations and adaptive stochastic gradient algorithms.
\newblock In \emph{Advances in Neural Information Processing Systems}, 2104-2113, 2019.

\bibitem{li2021validity}
Q.~Li, C.~Tai, and W.~E.
\newblock On the validity of modeling SGD with SDEs.
\newblock In \emph{Advances in Neural Information Processing Systems}, 2021.

\bibitem{LindvallRogers1986}
T.~Lindvall and L.~C.~G.~Rogers.
\newblock Coupling of multidimensional diffusions by reflection.
\newblock \emph{Ann. Probab.} \textbf{14} (1986), 860-872.

\bibitem{malladi2022sdes}
S.~Malladi, K.~Lyu, A.~Panigrahi, and S.~Arora.
\newblock On the SDEs and scaling rules for adaptive gradient algorithms.
\newblock In \emph{Advances in Neural Information Processing Systems}, Vol.~35, 2022.

\bibitem{MattinglyStuartHigham2002}
J.~C.~Mattingly, A.~M.~Stuart, and D.~J.~Higham.
\newblock Ergodicity for SDEs and approximations: locally Lipschitz vector fields and degenerate noise.
\newblock \emph{Stochastic Process. Appl.} \textbf{101} (2002), 185-232.





\bibitem{McMahanStreeter2010}
H.~B.~McMahan and M.~J.~Streeter.
\newblock Adaptive bound optimization for online convex optimization.
\newblock In \emph{Proc. COLT}, 244-256, 2010.

\bibitem{meyn1993foster}
S.~P.~Meyn and R.~L.~Tweedie.
\newblock Stability of Markovian processes III: Foster-Lyapunov criteria for continuous-time processes.
\newblock \emph{Adv. Appl. Probab.} \textbf{25} (1993), 518-548.

\bibitem{MeynTweedie2009}
S.~P.~Meyn and R.~L.~Tweedie.
\newblock \emph{Markov Chains and Stochastic Stability}.
\newblock 2nd ed., Cambridge University Press, 2009.



\bibitem{Nualart2006}
D.~Nualart.
\newblock \emph{The Malliavin Calculus and Related Topics}.
\newblock 2nd ed., Springer, Berlin, 2006.


\bibitem{Pigato2017}
P.~Pigato.
\newblock Tube estimates for diffusion processes under a weak H\"ormander condition.
\newblock \emph{Ann. Inst. H. Poincar\'e Probab. Statist.} \textbf{53} (2017), 2008-2033.

\bibitem{Pigato2022}
P.~Pigato.
\newblock Density estimates and short-time asymptotics for a hypoelliptic diffusion process.
\newblock \emph{Stochastic Process. Appl.} \textbf{145} (2022), 117-142.

\bibitem{Reddi2018}
S.~J.~Reddi, S.~Kale, and S.~Kumar.
\newblock On the convergence of Adam and beyond.
\newblock In \emph{Proc. ICLR}, 2018.

\bibitem{Ruder2016}
S.~Ruder.
\newblock An overview of gradient descent optimization algorithms.
\newblock \emph{arXiv:1609.04747}, 2017.

\bibitem{StroockVaradhan1972}
D.~W. Stroock and S.~R.~S. Varadhan.
\newblock On the support of diffusion processes with applications to the strong maximum principle.
\newblock In \emph{Proc. Sixth Berkeley Symp. on Mathematical Statistics and Probability}, Vol.~3 (Probability Theory), 333-359.
\newblock University of California Press, Berkeley, CA, 1972.

\bibitem{StroockVaradhan1979}
D.~W.~Stroock and S.~R.~S.~Varadhan.
\newblock \emph{Multidimensional Diffusion Processes}.
\newblock Springer, Berlin, 1979.

\bibitem{Sun2019}
R.~Sun.
\newblock Optimization for deep learning: theory and algorithms.
\newblock \emph{arXiv:1912.08957}, 2019.

\bibitem{TielemanHinton2012}
T.~Tieleman and G.~Hinton.
\newblock RMSProp: Divide the gradient by a running average of its recent magnitude.
\newblock Lecture notes for \emph{Neural Networks for Machine Learning}, University of Toronto, 2012.

\bibitem{Villani2009}
C.~Villani.
\newblock \emph{Hypocoercivity}.
\newblock American Mathematical Society, Providence, RI, 2009.

\bibitem{Zeiler2012}
M.~D.~Zeiler.
\newblock ADADELTA: An adaptive learning rate method.
\newblock \emph{arXiv:1212.5701}, 2012.

\bibitem{Zhang2023}
Y.~Zhang, C.~Chen, N.~Shi, R.~Sun, and Z.-Q.~Luo.
\newblock Adam can converge without any modification on update rules.
\newblock \emph{arXiv:2302.11505}, 2023.

\end{thebibliography}


\newpage
\appendix




\section{The continuous-time limit: Proof of Theorem \ref{Thm1}}\label{aaapp}

Recall that $t_k:=kh$, and that $\eta, \alpha, \beta, \{\xi_k\}$ are chosen as in \eqref{scale}. We begin with an elementary lemma.

\begin{lemma}\label{lem:bc}
Set
\[
\theta_k^{(a,h)}:=\frac{1-\alpha}{1-\alpha^{k+1}}
=\frac{ah}{1-(1-ah)^{k+1}},\qquad
\theta_k^{(b,h)}:=\frac{1-\beta}{1-\beta^{k+1}}
=\frac{bh}{1-(1-bh)^{k+1}},
\]
with $t_k=kh$. Then, for each fixed $T>0$,
\[
\sup_{k\le T/h}\Big|\,\theta_k^{(a,h)}-h\,c_a(t_{k+1})\,\Big| = O(h^2),
\qquad
\sup_{k\le T/h}\Big|\,\theta_k^{(b,h)}-h\,c_b(t_{k+1})\,\Big| = O(h^2),
\quad\text{as }h\to 0,
\]
where $c_a(t)={a}/{(1-e^{-at})}$ and $c_b(t)={b}/{(1-e^{-bt})}$.
\end{lemma}

\begin{proof}
We prove only the $a$-statement, as the $b$-case is identical. Let $t:=t_{k+1}=(k+1)h$. Using the Taylor expansion of the logarithm with a uniform remainder,
\[
\log(1-u)=-u-\frac{u^2}{2}+R(u),\qquad |R(u)|\le C u^3,
\]
valid for $|u|\le u_0$, and taking $u=ah$, we obtain uniformly for $k\le T/h$,
\[
(k+1)\log(1-ah)
= -a t - \frac{a^2}{2}t h + \rho_{k,h},
\qquad |\rho_{k,h}|\le C' h^2,
\]
since $(k+1)(ah)^3 = O(h^2)$. Exponentiating,
\[
(1-ah)^{k+1}
= e^{-a t}\big(1 - \tfrac{a^2}{2}t h + O(h^2)\big),
\]
uniformly for $k\le T/h$. Hence
\[
1-(1-ah)^{k+1}
= (1-e^{-a t}) + e^{-a t}\,\tfrac{a^2}{2}t h + O(h^2).
\]
Therefore,
\[
\theta_k^{(a,h)}
= \frac{ah}{1-(1-ah)^{k+1}}
= \frac{ah}{1-e^{-a t}}
\big(1 + O(h)\big)
= h\,c_a(t) + O(h^2),
\]
uniformly for $k\le T/h$.
\end{proof}

To begin the proof of Theorem \ref{Thm1}, fix $\delta>0$ and set
\[
k_\delta:=\min\{k\ge 0:\ t_k\ge \delta\}=\lceil \delta/h\rceil.
\]
Since $t_{k_\delta}\to\delta$ as $h\to0$, it suffices to prove convergence on $[t_{k_\delta},T]$. All sums below are taken over $k_\delta\le k\le N-1$, where $t\in[t_N,t_{N+1})$, and the value at time $t_{k_\delta}$ serves as the initial condition.

We organize the proof into several steps.

\subsection{The $z$–equation.}
From \eqref{eqn:sgd_main++}, Lemma~\ref{lem:bc}, and \eqref{scale},
\[
z_{k+1}-z_k
=
h\,c_a(t_{k+1})\bigl(-z_k+\nabla f(x_k)\bigr)
+\sqrt h\,c_a(t_{k+1})\sigma\,\zeta_k
+r_k^{(z)},
\]
where
\[
r_k^{(z)}
=
\bigl(\theta_k^{(a,h)}-h\,c_a(t_{k+1})\bigr)\bigl(-z_k+\nabla f(x_k)\bigr).
\]
Hence, for $t\in[t_N,t_{N+1})\subset[\delta,T]$,
\[
Z^h(t)
=
Z^h(t_{k_\delta})
+\sum_{k=k_\delta}^{N-1} h\,c_a(t_{k+1})\bigl(\nabla f(x_k)-z_k\bigr)
+\sum_{k=k_\delta}^{N-1}\sqrt h\,c_a(t_{k+1})\sigma\,\zeta_k
+R_h^{(z,\delta)}(t),
\]
with
\[
R_h^{(z,\delta)}(t)
=
\sum_{k=k_\delta}^{N-1}
\bigl(\theta_k^{(a,h)}-h\,c_a(t_{k+1})\bigr)\bigl(-z_k+\nabla f(x_k)\bigr).
\]
By Lemma~\ref{lem:bc},
\[
\sup_{k\le T/h}\|r_k^{(z)}\|=O(h^2),
\qquad
\sup_{t\in[\delta,T]}\|R^{(z,\delta)}_h(t)\|=O(h),
\]
in probability. The drift term converges to
\[
\int_\delta^t c_a(s)\bigl(\nabla f(x_s)-z_s\bigr)\,\mathrm{d}s,
\]
provided $(X^h,Z^h)\Rightarrow (x,z)$. For the martingale term
\[
M^h(t) = \sum_{k=k_\delta}^{N-1} \sqrt{h}\,c_a(t_{k+1})\sigma\,\zeta_k,
\]
its quadratic variation satisfies
\[
\langle M^h\rangle_t
= \sigma^2 \sum_{k=k_\delta}^{N-1} h\,(c_a(t_{k+1}))^2
\;\longrightarrow\;
\sigma^2 \int_\delta^t c_a(s)^2\,\mathrm{d}s.
\]
Since Lindeberg’s condition holds, the martingale functional CLT yields
\[
M^h \Rightarrow \int_\delta^t c_a(s)\sigma\,\mathrm{d}B_s
\quad \text{in } \mathbb{D}([\delta,T];\mathbb{R}^m).
\]
Hence any limit point satisfies \eqref{eq:cts-z}.

\subsection{The $y$–equation.}
Using the closure approximation \eqref{scale+}, we replace
\[
\bigl|\partial_{x_i}f(x_k)+\xi_k^i\bigr|^2
\;\leadsto\;
(\partial_{x_i}f(x_k))^2+\sigma^2.
\]
Thus,
\[
y_{k+1}^i-y_k^i
=
\theta_k^{(b,h)}\Bigl(-y_k^i+(\partial_{x_i}f(x_k))^2+\sigma^2\Bigr).
\]
Proceeding as in the $z$–equation,
\[
\frac{y_{k+1}^i-y_k^i}{h}
=
c_b(t_{k+1})\Bigl(-y_k^i+(\partial_{x_i}f(x_k))^2+\sigma^2\Bigr)
+ e_k^{(y,i)},
\]
with $e_k^{(y,i)}=O(h)$ in probability. Summing,
\[
Y^{h,i}(t)
=
Y^{h,i}(t_{k_\delta})
+\sum_{k=k_\delta}^{N-1} h\,c_b(t_{k+1})\Bigl(-y_k^i+(\partial_{x_i}f(x_k))^2+\sigma^2\Bigr)
+R_h^{(y,i)}(t),
\]
where $\sup_{t\in[\delta,T]}|R_h^{(y,i)}(t)|\to0$ in probability. Passing to the limit yields \eqref{eq:cts-y}.

\subsection{The $x$–equation.}
From $\eta=\gamma h$,
\[
x_{k+1}^i-x_k^i = -\gamma h\,\frac{z_{k+1}^i}{\sqrt{y_{k+1}^i}+\varepsilon}.
\]
Replacing $(z_{k+1},y_{k+1})$ by $(z_k,y_k)$ introduces an error of order $O(h)$ in probability. Hence
\[
X^h(t)
=
X^h(t_{k_\delta})
-\gamma\sum_{k=k_\delta}^{N-1} h\,\frac{Z^h(t_k)}{\sqrt{Y^h(t_k)}+\varepsilon}
+R_h^{(x,\delta)}(t),
\]
with $\sup_{t\in[\delta,T]}\|R_h^{(x,\delta)}(t)\|=O(h)$. Passing to the limit yields \eqref{eq:cts-x}.

\subsection{Identification of the limit.}
Define coefficients $b$ and $\sigma$ as in the statement. Under (A1), they are locally Lipschitz with linear growth on $[\delta,T]$, ensuring existence and uniqueness of a strong solution. The processes $(X^h,Z^h,Y^h)$ form an Euler-type scheme with martingale noise. The previous estimates show convergence of drift and diffusion terms. Standard results (e.g.\ Kurtz–Protter \cite{kurtz-protter}) imply tightness in $\mathbb{D}([\delta,T];\mathbb{R}^{3m})$ and identification of the limit as a weak solution. Pathwise uniqueness yields uniqueness in law, and hence convergence of the full sequence to the unique strong solution. \qed

\section{Long-term dynamics of the inhomogeneous system}\label{LLD}



\begin{proposition}\label{Proppa}
Assume that $f$ satisfies condition \textnormal{(A1)} from Subsection~\ref{Subcond} and fix $\varepsilon>0$, $a,b,\gamma,\sigma>0$. Let $(x_t,z_t,y_t)$ be the unique strong
solution of the time-inhomogeneous system \eqref{eq:cts-x}-\eqref{eq:cts-y}, defined for $t\ge \delta>0$, with initial state $(x_\delta,z_\delta,y_\delta)\in\R^{3m}$, driven by an $m$–dimensional Brownian motion $B_t$. For each
$s\ge \delta$, let $(\tilde x^{(s)}_t,\tilde z^{(s)}_t,\tilde y^{(s)}_t)_{t\ge 0}$ be the unique strong
solution of the autonomous system \eqref{eq:cts-x+}–\eqref{eq:cts-y+} started from the same state
and driven by the shifted Brownian motion, i.e.
\[
(\tilde x^{(s)}_0,\tilde z^{(s)}_0,\tilde y^{(s)}_0)=(x_s,z_s,y_s),
\qquad \tilde B^{(s)}_t=B_{s+t}-B_s.
\]
Then for every $T>0$ there exist constants $C_T<\infty$ depending only
on $T$, $L_f$, $a,b,\gamma,\varepsilon,\sigma$ and $\kappa=\min\{a,b\}>0$ such that
\begin{equation}\label{eq:AA-L2}
\mathbb{E}\Big[\sup_{0\le t\le T}\big\|(x_{s+t},z_{s+t},y_{s+t})
-(\tilde x^{(s)}_t,\tilde z^{(s)}_t,\tilde y^{(s)}_t)\big\|^2\Big]
\;\le\; C_T\,e^{-2\kappa s}\,\Big(1+\mathbb{E}\bigl [\|(x_s,z_s,y_s)\|^2\bigr ]\Big),\qquad s\ge \delta.
\end{equation}
In particular, for each fixed $T>0$,
\[
(x_{s+\cdot},z_{s+\cdot},y_{s+\cdot})
\;\longrightarrow\;
(\tilde x^{(s)}_{\cdot},\tilde z^{(s)}_{\cdot},\tilde y^{(s)}_{\cdot})
\quad\text{in }\mathrm L^2\!\big(\Omega;C([0,T];\R^{3m})\big)\ \text{and hence in probability, as } s\to\infty.
\]
\end{proposition}

\begin{proof}
To start the proof we first note that for $c_a(t)=a/(1-e^{-at})$ and $c_b(t)=b/(1-e^{-bt})$ we have, for every fixed $T>0$, that there exists $C_T>0$ such that
\begin{equation}\label{eq:ca_cb_decay}
\sup_{\delta\le u\le \delta+T}\bigl|c_a(s+u)-a\bigr| \le C_T e^{-as},\qquad
\sup_{\delta\le u\le \delta+T}\bigl|c_b(s+u)-b\bigr| \le C_T e^{-bs}.
\end{equation}
Under (A1)  $\nabla f$ satisfies
\[
\|\nabla f(x)-\nabla f(\bar x)\|\le L_f\|x-\bar x\|,\qquad
\|\nabla f(x)\|\le \|\nabla f(0)\|+L_f\|x\|.
\]
For componentwise squares,
\begin{equation}\label{eq:square_diff}
\|\nabla f(x)^{\odot 2}-\nabla f(\bar x)^{\odot 2}\|
\;\le\; \big(\|\nabla f(x)\|+\|\nabla f(\bar x)\|\big)\,\|\nabla f(x)-\nabla f(\bar x)\|.
\end{equation}
For the preconditioner $\Psi(z,y)=z\oslash(\sqrt{y}+\varepsilon)$,
\begin{equation}\label{eq:Psi_diff}
\|\Psi(z,y)-\Psi(\bar z,\bar y)\|
\;\le\; \frac{1}{\varepsilon}\|z-\bar z\|
+\frac{1}{\varepsilon^2}(\|z\|+\|\bar z\|)\,\|\sqrt{y}-\sqrt{\bar y}\|.
\end{equation}
and
\[
\|\sqrt{y}-\sqrt{\bar y}\|\le\|y-\bar y\|^{1/2}\le \delta\|y-\bar y\|+\tfrac{1}{4\delta}.
\]
Standard linear growth and Burkholder–Davis–Gundy and Grönwall arguments yield, for any $T>0$,
\begin{align}\label{eq:moment_bounds}
&\mathbb{E}\Big[\sup_{0\le u\le T}\big(\|x_{s+u}\|^2+\|z_{s+u}\|^2+\|y_{s+u}\|^2\big)\Big]
+\mathbb{E}\Big[\sup_{0\le u\le T}\big(\|\tilde x_u\|^2+\|\tilde z_u\|^2+\|\tilde y_u\|^2\big)\Big]\notag\\
&\le\; C_T\Big(1+\mathbb{E}\bigl [\|(x_s,z_s,y_s)\|^2\bigr ]\Big).
\end{align}
We introduce the differences
\[
\Delta x_t=x_{s+t}-\tilde x_t,\quad
\Delta z_t=z_{s+t}-\tilde z_t,\quad
\Delta y_t=y_{s+t}-\tilde y_t.
\]
Subtracting the SDEs on $[s,s+t]$ and $[0,t]$ with the coupled Brownian motion, one obtains
\[
\Delta x_t = -\gamma\int_0^t\bigl[\Psi(z_{s+u},y_{s+u})-\Psi(\tilde z_u,\tilde y_u)\bigr]\mathrm{d}u
\]
and corresponding expressions for $\Delta z_t$ and $\Delta y_t$. Using \eqref{eq:Psi_diff}–\eqref{eq:moment_bounds}, we obtain
\[
\mathbb{E}\Big[\sup_{0\le u\le t}\|\Delta x_u\|^2\Big]
\;\le\; C \int_0^t \mathbb{E}\big[\|\Delta z_u\|^2+\|\Delta y_u\|^2\big]\mathrm{d}u
+ C_T\big(1+\mathbb{E}\bigl [\|(x_s,z_s,y_s)\|^2\bigr ]\big).
\]
Similarly, for $\Delta z_t$ we use Burkholder–Davis–Gundy and \eqref{eq:ca_cb_decay} to get
\[
\mathbb{E}\Big[\sup_{0\le u\le t}\|\Delta z_u\|^2\Big]
\;\le\; C_T e^{-2as}\big(1+\mathbb{E}\bigl [\|(x_s,z_s,y_s)\|^2\bigr ]\big)
+ C \int_0^t \mathbb{E}\Big[\sup_{0\le r\le u}\|\Delta x_r\|^2+\|\Delta z_r\|^2\Big]\mathrm{d}u.
\]
For $\Delta y_t$,
\[
\mathbb{E}\Big[\sup_{0\le u\le t}\|\Delta y_u\|^2\Big]
\;\le\; C_T e^{-2bs}\big(1+\mathbb{E}\bigl [\|(x_s,z_s,y_s)\|^2\bigr ]\big)
+ C \int_0^t \mathbb{E}\Big[\sup_{0\le r\le u}\|\Delta x_r\|^2+\|\Delta y_r\|^2\Big]\mathrm{d}u.
\]
Finally, let
\[
\Phi(t)=\mathbb{E}\Big[\sup_{0\le u\le t}\big(\|\Delta x_u\|^2+\|\Delta z_u\|^2+\|\Delta y_u\|^2\big)\Big].
\]
Combining the above,
\[
\Phi(t) \;\le\; C_T e^{-2\kappa s}\big(1+\mathbb{E}(\|(x_s,z_s,y_s)\|^2)\big)
+ C\int_0^t \Phi(u)\,\mathrm{d}u,
\qquad \kappa=\min\{a,b\}.
\]
By Grönwall’s inequality,
\[
\Phi(T)\;\le\; C_T e^{-2\kappa s}\big(1+\mathbb{E}(\|(x_s,z_s,y_s)\|^2)\big).
\]
This proves the stated bound and convergence in probability as $s\to\infty$.
\end{proof}



\begin{lemma}\label{thm:same-equilibrium}
Assume the hypotheses of Proposition~\ref{Proppa}.
Let \(X_t=(x_t,z_t,y_t)\) denote the solution of the time-inhomogeneous system
\eqref{eq:cts-x}-\eqref{eq:cts-y}, defined for \(t\ge \delta>0\), with initial distribution
\(\mathrm{Law}(X_\delta)=\mu_\delta\).
Let \(P_t\) be the Markov semigroup associated with the autonomous system
\eqref{eq:cts-x+}-\eqref{eq:cts-y+}. Assume that the semigroup \(P_t\) admits a unique invariant probability measure \(\pi_\infty\). Then every weak limit point of \(\mathrm{Law}(X_t)\) is equal to \(\pi_\infty\).
In particular,
\[
\mathrm{Law}(X_t)\Rightarrow \pi_\infty
\qquad\text{as }t\to\infty.
\]
Hence, the time-inhomogeneous system and the autonomous system have the same
equilibrium distribution.
\end{lemma}

\begin{proof} For the following argument it is important to first note that the family of laws \(\{\mathrm{Law}(X_t):t\ge \delta\}\) is tight. To prove this it is enough to establish a uniform moment bound
\begin{equation}\label{haa}
\sup_{t\ge \delta}\mathbb E[\|X_t\|^2]<\infty.
\end{equation}
Indeed, assuming \eqref{haa} we have by Markov's inequality,
\[
\sup_{t\ge \delta}\mathbb P(\|X_t\|>R)
\le \frac{1}{R^2}\sup_{t\ge \delta}\mathbb E[\|X_t\|^2],
\]
and therefore, for every $\varepsilon>0$, we can choose $R$ sufficiently large so that
\[
\sup_{t\ge \delta}\operatorname{Law}(X_t)(B(0,R)^c)<\varepsilon.
\]
Since closed balls in $\mathbb R^{3m}$ are compact, Prokhorov's theorem now implies that
$\{\operatorname{Law}(X_t):t\ge \delta\}$ is tight. However, \eqref{haa} follows from the results of Appendix \ref{HMI} below, see Remark \ref{Inhomo}.

Having concluded that the family of laws \(\{\mathrm{Law}(X_t):t\ge \delta\}\) is tight, we let \(\nu_t=\mathrm{Law}(X_t)\) for \(t\ge \delta\). The family \(\{\nu_t:t\ge \delta\}\) is tight and hence for every sequence \(\{t_n\}\to\infty\), there exists a subsequence, still denoted
\(\{t_n\}\), and a probability measure \(\nu\) such that
\[
\nu_{t_n}\Rightarrow \nu
\qquad\text{as }n\to\infty.
\]
It is enough to show that \(\nu\) is invariant for \(P_t\) for every \(t\ge 0\).
By uniqueness of the invariant probability measure of \(P_t\), this will imply
\(\nu=\pi_\infty\). Fix \(t>0\), and let \(\varphi:\mathbb R^{3m}\to\mathbb R\) be bounded and continuous.
By Proposition~\ref{Proppa}, we have
\begin{equation}\label{eq:asp-semigroup}
\Big|
\mathbb E[\varphi(X_{s+t})]
-
\mathbb E[P_t\varphi(X_s)]
\Big|
\longrightarrow 0
\qquad\text{as }s\to\infty.
\end{equation}
Applying this with \(s=t_n\), we obtain
\[
\iiint \varphi\,\mathrm{d}\nu_{t_n+t}
-
\iiint P_t\varphi\,\mathrm{d}\nu_{t_n}
\longrightarrow 0.
\]
Passing to yet an other subsequence, if necessary, tightness also gives
\[
\nu_{t_n+t}\Rightarrow \bar\nu,
\]
for some probability measure \(\bar\nu\). Therefore,
\[
\iiint \varphi\,\mathrm{d}\bar\nu
=
\iiint P_t\varphi\,\mathrm{d}\nu.
\]
In particular, choosing a subsequence such that both \(\nu_{t_n}\Rightarrow \nu\) and
\(\nu_{t_n+t}\Rightarrow \nu\), we obtain
\[
\iiint \varphi\,\mathrm{d}\nu
=
\iiint P_t\varphi\,\mathrm{d}\nu
\qquad\text{for all bounded continuous }\varphi.
\]
Hence
\[
\nu P_t=\nu,
\]
so \(\nu\) is an invariant probability measure for the autonomous semigroup \(P_t\).
By uniqueness, \(\nu=\pi_\infty\). Thus every weak subsequential limit of \(\nu_t\) equals \(\pi_\infty\), and therefore
\[
\nu_t=\mathrm{Law}(X_t)\Rightarrow \pi_\infty
\qquad\text{as }t\to\infty.
\]
This proof the lemma.
\end{proof}

\begin{remark}
Although the time-inhomogeneous system does not generate a Markov semigroup,
Proposition~\ref{Proppa} shows that its transition mechanism becomes
asymptotically close to that of the autonomous semigroup \(P_t\).
In particular, for large times, the law of \(X_t\) behaves like \(\mu_0 P_t\),
which explains why both systems share the same invariant distribution.
\end{remark}



\section{Higher moments and integrability}\label{HMI} Let $X_t = ( x_t, z_t, y_t)$ solve the homogeneous system \eqref{eq:cts-x}-\eqref{eq:cts-y}. In the following we consistently assume that
\begin{equation}\label{intimeasureagain}
\mu_0\big(\mathbb R^m\times\mathbb R^m\times\mathbb R^m\big)=1
\quad\text{and}\quad
\mathbb{E}_{\mu_0}\!\big [V(x_0,z_0,y_0)\big]<\infty,
\end{equation}
where $V=V_0$. As ${\mathcal L}={\mathcal L}_0$, considering $\epsilon=0$, \eqref{eq:HL-drift+} states that
\begin{equation}\label{eq:HL-driftagain}
{\mathcal L} V_\upsilon\;\le\;-\lambda' V_\upsilon + K'\,\mathbf 1_{C},
\end{equation}
with $\lambda'$ and $K'$ independent of $\upsilon\in (0,1)$. Applying Dynkin's formula to $V_\upsilon(X_t)$ yields
\[
\mathbb{E}_{(x_0,z_0,y_0)}[V_\upsilon( X_t)]
= \mathbb{E}_{(x_0,z_0,y_0)}[V_\upsilon( X_0)]
+ \mathbb{E}_{(x_0,z_0,y_0)}\!\left[\int_0^t {\mathcal L} V_\upsilon( X_s)\,\mathrm{d}s\right].
\]
Using \eqref{eq:HL-driftagain} in the integrand gives
\[
\frac{\mathrm{d}}{\mathrm{d}t}\,\mathbb{E}_{(x_0,z_0,y_0)}[V_\upsilon( X_t)]
\;\le\; -\lambda'\,\mathbb{E}_{(x_0,z_0,y_0)}[V_\upsilon( X_t)] + K'.
\]
By Grönwall’s inequality we see that there exists a finite constant $C>0$, independent of $(x_0,z_0,y_0)$, $t\geq 0$, and $\upsilon\in (0,1)$, such that
\[
\mathbb{E}_{(x_0,z_0,y_0)}[V_\upsilon( X_t)]\leq C\implies
\sup_{t \ge 0}\,\mathbb{E}_{(x_0,z_0,y_0)}[V_0( X_t)]\leq C\implies
\sup_{t \ge 0}\,\mathbb{E}_{\mu_0}[V_0( X_t)]\leq C.
\]
By \eqref{eq:drift-ineql} and the fact that $\|y\|\le \|y\|_1$, we therefore deduce that
\begin{equation}\label{eq:2nd-moment-unif}
\sup_{t\ge0}\,\mathbb E_{\mu_0}\bigl[\| x_t\|^2+\| z_t\|^2+\| y_t\|\bigr]\ <\ \infty.
\end{equation}
In particular, $\pi_t=\mu_0P_t$ has finite second moments in $ x_t$ and $ z_t$, and finite first moments in $ y_t$. In the following we are going to use a bootstrap argument to control second and higher order moments of $y_t$.


\begin{lemma}\label{lem:power-Lyap} Assume that $f$ satisfies \textnormal{(A1)} and \textnormal{(A2)}. Let $V$ be given by \eqref{V} with $\alpha=1$. Assume that there exist
$\lambda',K'>0$, a compact set $C\subset\R^{3m}$, and constants $c_1,c_2>0$ such that, for all $(x,z,y)\in \mathbb{R}^m \times \mathbb{R}^m \times \mathbb{R}^m$, and for all $\epsilon\in [0,\upsilon]$, $\upsilon\in (0,1)$,
\[
{\mathcal L_\epsilon}V_\upsilon \;\le\; -\lambda' V_\upsilon + K'\,\mathbf 1_{C},
\qquad
V_\upsilon(x,z,y)\;\ge\; \hat c_1\bigl(\|x\|^2+\|z\|^2+\|y\|_{1,\upsilon}\bigr)-\hat c_2.
\]
Then there exists $\eta>0$ small and
constants $\lambda^{''}, K^{''}>0$, all independent of $\epsilon$ and $\upsilon$, such that
\begin{equation}\label{eq:power-drift}
{\mathcal L_\epsilon}\,\Phi(V_\upsilon) \;\le\; -\,\lambda^{''}\,\Phi(V_\upsilon)\;+\; K^{''}\,\mathbf 1_{C},
\qquad
\Phi(s):=(1+s)^{1+\eta}.
\end{equation}
Consequently,
\[
\sup_{t\ge0}\,\mathbb E_{\mu_0}\bigl[\Phi\!\bigl(V_\upsilon(X_t)\bigr)\bigr]\;<\;\infty.
\]
\end{lemma}
\begin{proof}
Let $\Phi(s)=(1+s)^{1+\eta}$ with $\eta\in(0,1)$ to be chosen small.
We compute ${\mathcal L}_\epsilon\Phi(V_\upsilon)$ using the chain rule for generators.
Since the diffusion acts only in the $z$–variables, we obtain
\[
{\mathcal L}_\epsilon\,\Phi(V_\upsilon)
=\Phi'(V_\upsilon)\,{\mathcal L}_\epsilon V_\upsilon
+\frac{a^2\sigma^2}{2}\Phi''(V_\upsilon)
\sum_{i=1}^m(\partial_{z_i}V_\upsilon)^2 .
\]
Here
\[
\Phi'(s)=(1+\eta)(1+s)^{\eta},
\qquad
\Phi''(s)=\eta(1+\eta)(1+s)^{\eta-1}.
\]

\medskip

\noindent
Since $\alpha=1$, we have
\[
V_\upsilon(x,z,y)
= \theta(f(x)-f_\ast)
+ \frac12\|z\|^2
- \beta\, x\cdot z
+ \delta\|y\|_{1,\upsilon},
\]
and we estimate
\[
\sum_{i=1}^m(\partial_{z_i}V_\upsilon)^2
= \|z-\beta x\|^2
\le 2\|z\|^2 + 2\beta^2\|x\|^2.
\]
By the coercive lower bound on $V_\upsilon$ we see that there exist $C_0>0$, $C_1>0$, independent of $\epsilon$ and $\upsilon$,
such that
\[
\|z\|^2+\|x\|^2
\le C_0(1+V_\upsilon)\implies
\sum_{i=1}^m(\partial_{z_i}V_\upsilon)^2
\le C_1(1+V_\upsilon).
\]
Therefore,
\[
\frac{a^2\sigma^2}{2}\Phi''(V_\upsilon)
\sum_{i=1}^m(\partial_{z_i}V_\upsilon)^2
\le C_2(1+V_\upsilon)^{\eta},
\]
for some constant $C_2>0$ independent of $\epsilon,\upsilon$. By assumption,
\[
{\mathcal L}_\epsilon V_\upsilon
\le -\lambda' V_\upsilon + K'\mathbf 1_{C},
\]
uniformly in $\epsilon\in[0,\upsilon]$, $\upsilon\in(0,1)$.
Hence
\[
{\mathcal L}_\epsilon\Phi(V_\upsilon)
\le (1+\eta)(1+V_\upsilon)^{\eta}
\bigl(-\lambda' V_\upsilon + K'\mathbf 1_{C}\bigr)
+ C_2(1+V_\upsilon)^{\eta}.
\]
Enlarge $C$ if necessary so that $V_\upsilon\ge v_0>0$ on $C^c$.
Then on $C^c$,
\[
(1+V_\upsilon)^{\eta}V_\upsilon
\ge \frac{v_0}{1+v_0}(1+V_\upsilon)^{1+\eta}
=: c_\star\,\Phi(V_\upsilon)\implies
(1+V_\upsilon)^{\eta}
\le \frac{1}{1+v_0}\Phi(V_\upsilon).
\]
Therefore, on $C^c$,
\[
{\mathcal L}_\epsilon\Phi(V_\upsilon)
\le -\lambda'(1+\eta)c_\star\,\Phi(V_\upsilon)
+ \frac{C_2}{1+v_0}\Phi(V_\upsilon).
\]
Choosing $\eta>0$ sufficiently small and enlarging $C$ if necessary,
we obtain
\[
{\mathcal L}_\epsilon\Phi(V_\upsilon)
\le -\lambda^{''}\Phi(V_\upsilon)
\qquad \text{on } {C}^c,
\]
for some $\lambda^{''}>0$ independent of $\epsilon,\upsilon$. Since $C$ is compact and ${\mathcal L}_\epsilon\Phi(V_\upsilon)$
is continuous, there exists $K^{''}>0$,
independent of $\epsilon,\upsilon$, such that
\[
{\mathcal L}_\epsilon\Phi(V_\upsilon)
\le K^{''}\mathbf 1_{C}.
\]
Combining the estimates on $C$ and ${C}^c$ yields
\[
{\mathcal L}_\epsilon\Phi(V_\upsilon)
\le -\lambda^{''}\Phi(V_\upsilon)
+ K^{''}\mathbf 1_{C},
\]
which proves \eqref{eq:power-drift}. Applying Dynkin's formula and Grönwall's inequality then gives
\[
\sup_{t\ge0}
\mathbb E_{\mu_0}\bigl[\Phi(V_\upsilon(X_t))\bigr]
<\infty,
\]
uniformly in $\epsilon$ and $\upsilon$. This completes the proof.
\end{proof}



\begin{lemma}\label{lem:power-fixedq}
Assume that $f$ satisfies \textnormal{(A1)} and \textnormal{(A2)}. Let $V$ be given by \eqref{V} with $\alpha=1$. Assume that there exist
$\lambda',K'>0$, a compact set $C\subset\R^{3m}$, and constants $c_1,c_2>0$ such that, for all $(x,z,y)\in \mathbb{R}^m \times \mathbb{R}^m \times \mathbb{R}^m$, and for all $\epsilon\in [0,\upsilon]$, $\upsilon\in (0,1)$,
\[
{\mathcal L_\epsilon}V_\upsilon \;\le\; -\lambda' V_\upsilon + K'\,\mathbf 1_{C},
\qquad
V_\upsilon(x,z,y)\;\ge\; \hat c_1\bigl(\|x\|^2+\|z\|^2+\|y\|_{1,\upsilon}\bigr)-\hat c_2.
\]
Then, for every fixed $q>1$, there exist constants
$\lambda_q,K_q>0$ and a compact set $C_q\supset C$ (all depending on $q$), all independent of $\epsilon$ and $\upsilon$, such that
\begin{equation}\label{eq:powerq-drift+}
{\mathcal L_\epsilon}\,(1+V_\upsilon)^q \;\le\; -\,\lambda_q\,(1+V_\upsilon)^q \;+\; K_q\,\mathbf 1_{C_q}.
\end{equation}
Consequently, for every fixed $q>1$,
\begin{equation}\label{eq:powerq-drift+aa}
\sup_{t\ge0}\ \mathbb E_{\mu_0}\bigl[(1+V_\upsilon( X_t))^q \bigr]<\hat c<\infty
\end{equation}
for a constant $\hat c$ independent of $\upsilon$.
\end{lemma}

\begin{proof}  Fix $q>1$ and define $\Phi(s)=(1+s)^q$. As in the proof of Lemma \ref{lem:power-Lyap} we deduce \[
\frac{a^2\sigma^2}{2}
\Phi''(V_\upsilon)
\sum_{i=1}^m (\partial_{z_i}V_\upsilon)^2
\le C_2 (1+V_\upsilon)^{q-1},
\]
for some constant $C_2>0$ independent of $\epsilon,\upsilon$. By assumption,
\[
{\mathcal L}_\epsilon V_\upsilon
\le -\lambda' V_\upsilon + K' \mathbf 1_C,
\]
uniformly in $\epsilon\in[0,\upsilon]$, $\upsilon\in(0,1)$.
Thus
\[
{\mathcal L}_\epsilon \Phi(V_\upsilon)
\le
- \lambda' q V_\upsilon (1+V_\upsilon)^{q-1}
+ qK'(1+V_\upsilon)^{q-1}\mathbf 1_C
+ C_2 (1+V_\upsilon)^{q-1}.
\]
Enlarge $C$ if necessary so that
\[
V_\upsilon \ge v_0>0
\qquad \text{on } C^c.
\]
Then on $C^c$,
\[
V_\upsilon (1+V_\upsilon)^{q-1}
\ge \frac{v_0}{1+v_0}(1+V_\upsilon)^q.
\]
Hence, on $C^c$,
\[
{\mathcal L}_\epsilon \Phi(V_\upsilon)
\le
- \lambda' q \frac{v_0}{1+v_0}(1+V_\upsilon)^q
+ C_2 (1+V_\upsilon)^{q-1}.
\]
Since
\[
(1+V_\upsilon)^{q-1}
= \frac{1}{1+V_\upsilon}(1+V_\upsilon)^q
\le \frac{1}{1+v_0}(1+V_\upsilon)^q
\quad \text{on } C^c,
\]
we obtain
\[
{\mathcal L}_\epsilon \Phi(V_\upsilon)
\le
-\left(
\lambda' q \frac{v_0}{1+v_0}
- \frac{C_2}{1+v_0}
\right)
(1+V_\upsilon)^q .
\]
Choosing $v_0$ sufficiently large (i.e.\ enlarging $C$ if necessary),
we obtain
\[
{\mathcal L}_\epsilon \Phi(V_\upsilon)
\le -\lambda_q (1+V_\upsilon)^q
\qquad \text{on } C^c,
\]
for some $\lambda_q>0$ independent of $\epsilon,\upsilon$. Since $C$ is compact and all coefficients are continuous,
there exists $K_q>0$ such that
\[
{\mathcal L}_\epsilon (1+V_\upsilon)^q
\le K_q \mathbf 1_{C_q},
\]
where $C_q\supset C$ is a compact set.
Combining with the estimate on $C^c$ yields
\eqref{eq:powerq-drift+}. Applying Dynkin's formula to $(1+V_\upsilon(X_t))^q$
and using \eqref{eq:powerq-drift+},
Grönwall's inequality gives
\[
\sup_{t\ge0}
\mathbb E_{\mu_0}\bigl[(1+V_\upsilon(X_t))^q\bigr]
<\infty,
\]
with a constant independent of $\upsilon$.
This proves \eqref{eq:powerq-drift+aa}. The constants $\lambda_q,K_q$ and the set $C_q$
may depend on $q$, but are independent of $\epsilon$ and $\upsilon$.
\end{proof}



\begin{corollary}\label{cor:y-2+delta}
For any $q\geq 1$,
\[
\sup_{t\ge0}\ \mathbb E_{\mu_0}\,\bigl[\| y_t\|^{q}\bigr]\ <\ \infty.
\]
\end{corollary}

\begin{proof}
It is sufficient to prove the result for $q>2$ fixed. Since $y_t\in\R^m$, we have
\[
\|y_t\|^2=\sum_{i=1}^m |y_t^i|^2
\le \left(\sum_{i=1}^m |y_t^i|\right)^2
= \|y_t\|_1^2\implies
\|y_t\|^{q}
=\bigl(\|y_t\|^2\bigr)^{q/2}
\le \|y_t\|_1^{\,q}.
\]
Recall that
\[
\|y_t\|_{1,\upsilon}
=\sum_{i=1}^m \sqrt{|y_t^i|^2+\upsilon^2}
\ge \sum_{i=1}^m |y_t^i|
=\|y_t\|_1\implies
\|y_t\|^{q}
\le \|y_t\|_1^{\,q}
\le \|y_t\|_{1,\upsilon}^{\,q}.
\]
From the coercive lower bound on $V_\upsilon$,
\[
V_\upsilon(x,z,y)
\ge \hat c_1\bigl(\|x\|^2+\|z\|^2+\|y\|_{1,\upsilon}\bigr)-\hat c_2,
\]
we deduce
\[
\|y\|_{1,\upsilon}
\le \frac{1}{\hat c_1}\bigl(V_\upsilon+\hat c_2\bigr)
\le C(1+V_\upsilon)
\]
for some constant $C>0$ independent of $\epsilon$ and $\upsilon$.
Therefore
\[
\|y_t\|_{1,\upsilon}^{\,q}
\le C (1+V_\upsilon(X_t))^{q}\implies
\|y_t\|^{q}
\le C (1+V_\upsilon(X_t))^{q}.
\]
Applying Lemma~\ref{lem:power-fixedq} with this fixed $q>2$,
we obtain
\[
\sup_{t\ge0}
\mathbb E_{\mu_0}\bigl[(1+V_\upsilon(X_t))^{q}\bigr]
<\infty,
\]
with a bound independent of $\upsilon$.
Consequently,
\[
\sup_{t\ge0}
\mathbb E_{\mu_0}\bigl[\|y_t\|^{q}\bigr]
\le
C
\sup_{t\ge0}
\mathbb E_{\mu_0}\bigl[(1+V_\upsilon(X_t))^{q}\bigr]
<\infty.
\]
This completes the proof.
\end{proof}


\begin{remark}\label{Inhomo} We claim that the argument of this appendix can be revisited to also conclude that if \(X_t=(x_t,z_t,y_t)\) denotes the solution of the time-inhomogeneous system
\eqref{eq:cts-x}-\eqref{eq:cts-y}, defined for \(t\ge \delta>0\), with initial distribution
\(\mathrm{Law}(X_\delta)=\mu_\delta\), then, for every $q>2$,
\[
\sup_{t\ge\delta}
\mathbb E_{\mu_\delta}\bigl[(1+V_\upsilon(X_t))^{q}\bigr]
<\infty,
\]
with a bound independent of $\upsilon$. This implies uniform moment bound
\begin{equation*}
\sup_{t\ge \delta}\mathbb E_{\mu_\delta}[\|X_t\|^2]<\infty.
\end{equation*}
\end{remark}




\end{document}

